% Limites et continuité d’une fonction — Mathématiques 1re Tech — Cours signé OnBuch+
% Programme officiel MINESEC. Compiler avec : tectonic N13-limites-continuite-fonction.tex
\newcommand{\DOCMATIERE}{Mathématiques}
\newcommand{\DOCNIVEAU}{1re Tech}
\newcommand{\DOCMODULE}{Mathématiques – classe de Première technique}
\newcommand{\DOCLECON}{N13}
\newcommand{\DOCTITRE}{Limites et continuité d’une fonction}
% OnBuch+ — préambule « cours signés » ENRICHI (compilé avec tectonic / XeLaTeX)
% Système visuel « Pop » d'OnBuch+ : fond crème (jamais blanc), bordures encre
% épaisses, ombres « dures » décalées, titres Archivo Black en capitales.
% Version MATHÉMATIQUES : graphes (pgfplots/tikz), boîtes pédagogiques
% (définition, propriété, méthode, attention, le savais-tu, exercice résolu,
% exercice, TP, bilan), page de garde et signature OnBuch+ ancrée en bas.
%
% Macros attendues dans main.tex AVANT \input{preamble} :
%   \DOCMATIERE  \DOCNIVEAU  \DOCMODULE  \DOCLECON (numéro)  \DOCTITRE
\documentclass[11pt]{article}

\usepackage{fontspec}
\usepackage[french]{babel}
\usepackage[a4paper,margin=2cm,top=3.4cm,bottom=2.8cm,headheight=1.4cm,footskip=1.3cm]{geometry}
\usepackage{amsmath,amssymb,amsfonts}
\usepackage{mathtools} % \xmapsto, \coloneqq... souvent utilisés par les modèles
\usepackage{cancel} % simplifications barrées
% \abs{z} (module, valeur absolue) et \norm{u} (norme), utilisés par les modèles
\providecommand{\abs}{}\let\abs\relax\DeclarePairedDelimiter{\abs}{\lvert}{\rvert}
\providecommand{\norm}{}\let\norm\relax\DeclarePairedDelimiter{\norm}{\lVert}{\rVert}
\usepackage[table]{xcolor} % [table] : fournit aussi \rowcolors (alternance de lignes)
\usepackage{pagecolor}
\usepackage{tikz}
\usetikzlibrary{arrows.meta,positioning,calc,shapes.geometric,shapes.misc,shapes.arrows,decorations.pathmorphing,decorations.pathreplacing,decorations.markings,patterns,shadows,mindmap,trees,fit,backgrounds,matrix,intersections,angles,quotes}
\usepackage{pgfplots}
\usepackage{pgf-pie} % diagrammes circulaires \pie (statistiques)
\pgfplotsset{compat=1.18}
\usepgfplotslibrary{fillbetween,groupplots} % aires entre courbes (calcul intégral)
\usepackage{enumitem}
\usepackage{fancyhdr}
% [most] charge toutes les bibliothèques tcolorbox (listings, minted, raster,
% documentation...) et gonfle inutilement la mémoire TeX sur un document long
% (~300 boîtes) : on ne charge que ce qui est réellement utilisé.
\usepackage[skins,breakable]{tcolorbox}
\usepackage{graphicx}
\usepackage{colortbl}
\usepackage{booktabs}
\usepackage{tabularx}
\usepackage{diagbox} % case d'en-tête diagonale des tableaux à double entrée
% Colonnes extensibles centrée / gauche / droite (C, L, R), souvent utilisées par les modèles
\newcolumntype{C}{>{\centering\arraybackslash}X}
\newcolumntype{L}{>{\raggedright\arraybackslash}X}
\newcolumntype{R}{>{\raggedleft\arraybackslash}X}
\usepackage{array}
\usepackage{multicol}
\usepackage{siunitx}
% parse-numbers=false : accepte aussi \SI{2,5 \times 10^{-3}}{...}
\sisetup{locale=FR,output-decimal-marker={,},group-minimum-digits=4,parse-numbers=false}
\DeclareSIUnit\ohm{\ensuremath{\symup{\Omega}}} % U+2126 absent de la police
\usepackage{qrcode}
% \degree : commande fréquemment utilisée par les modèles pour un angle ou une
% température en dehors de \SI{}{} (non fournie par les paquets chargés).
\providecommand{\degree}{^{\circ}}
% \qed (fin de démonstration), utilisé par les modèles sans amsthm
\providecommand{\qed}{\hfill$\square$}
% \barre : double barre des valeurs interdites dans un tableau de variations (tkz-tab)
\providecommand{\barre}{\ensuremath{\|}}
% environnement proof (amsthm non chargé) utilisé par les modèles
\ifdefined\proof\else\newenvironment{proof}[1][Démonstration]{\par\noindent\textit{#1.}\ }{\qed\par}\fi
\usepackage{lastpage}
\usepackage{hyperref}
\hypersetup{hidelinks}

% ============================================================
% Palette « Pop » OnBuch+ — mêmes hex que la classe P (plus_ui.dart)
% ============================================================
\definecolor{poporange}{HTML}{F59321}
\definecolor{popdark}{HTML}{E07A0C}
\definecolor{popcream}{HTML}{FFF6E8}
\definecolor{popink}{HTML}{1C1714}
\definecolor{poppurple}{HTML}{7A5AE0}
\definecolor{popgreen}{HTML}{2E9E6A}
\definecolor{poppink}{HTML}{E0457B}
\definecolor{popblue}{HTML}{2F7DE1}
\definecolor{popgold}{HTML}{C98A12}
\definecolor{popmuted}{HTML}{8A7F76}
\definecolor{popline}{HTML}{EADFD2}
\definecolor{popgray}{HTML}{8A7F76}
\colorlet{popgrey}{popgray}
% Alias tolérants (le modèle nomme parfois les couleurs différemment)
\colorlet{popwhite}{white}
\colorlet{popblack}{popink}
\colorlet{popred}{poppink}
\colorlet{popyellow}{popgold}
% Teintes claires pour fonds de tableaux / zones
\colorlet{popblueL}{popblue!10!white}
\colorlet{poporangeL}{poporange!14!white}
\colorlet{popgreenL}{popgreen!12!white}
\colorlet{poppurpleL}{poppurple!12!white}
\colorlet{poppinkL}{poppink!12!white}
\colorlet{popgoldL}{popgold!14!white}

\pagecolor{popcream}

% ============================================================
% Typographie
% ============================================================
\defaultfontfeatures{Ligatures=TeX}
\setmainfont{PlusJakartaSans-Medium.ttf}[Path=../../../fonts/,BoldFont=PlusJakartaSans-Bold.ttf]
% Maths en sans-serif (Fira Math) avec chiffres et lettres droites de Plus
% Jakarta, pour que les formules se fondent dans le texte.
\usepackage[math-style=ISO,bold-style=ISO]{unicode-math}
\setmathfont{FiraMath-Regular.otf}[Path=../../../fonts/]
\setmathfont{PlusJakartaSans-Medium.ttf}[Path=../../../fonts/,range={up/num,up/{Latin,latin},\mathup}]
\usepackage{icomma} % virgule décimale sans espace en mode math
% Caractères mathématiques Unicode absents des polices de texte (Plus Jakarta,
% Archivo Black) : rendus par leur équivalent mathématique, en texte comme en
% formule — sinon ils s'affichent en carré vide (ex. « Arithmétique dans ℤ »).
\usepackage{newunicodechar}
\newunicodechar{ℝ}{\ensuremath{\mathbb{R}}}
\newunicodechar{ℤ}{\ensuremath{\mathbb{Z}}}
\newunicodechar{ℂ}{\ensuremath{\mathbb{C}}}
\newunicodechar{ℕ}{\ensuremath{\mathbb{N}}}
\newunicodechar{ℚ}{\ensuremath{\mathbb{Q}}}
\newunicodechar{𝔻}{\ensuremath{\mathbb{D}}}
\newunicodechar{α}{\ensuremath{\alpha}}
\newunicodechar{β}{\ensuremath{\beta}}
\newunicodechar{γ}{\ensuremath{\gamma}}
\newunicodechar{δ}{\ensuremath{\delta}}
\newunicodechar{ε}{\ensuremath{\varepsilon}}
\newunicodechar{θ}{\ensuremath{\theta}}
\newunicodechar{λ}{\ensuremath{\lambda}}
\newunicodechar{μ}{\ensuremath{\mu}}
\newunicodechar{σ}{\ensuremath{\sigma}}
\newunicodechar{τ}{\ensuremath{\tau}}
\newunicodechar{φ}{\ensuremath{\varphi}}
\newunicodechar{ψ}{\ensuremath{\psi}}
\newunicodechar{ω}{\ensuremath{\omega}}
\newunicodechar{Σ}{\ensuremath{\Sigma}}
% majuscules produites par \MakeUppercase dans les titres (α-aminés -> Α)
\newunicodechar{Α}{\ensuremath{\alpha}}
\newunicodechar{Β}{\ensuremath{\beta}}
\newunicodechar{Γ}{\ensuremath{\Gamma}}
\newunicodechar{Δ}{\ensuremath{\Delta}}
\newunicodechar{Ε}{\ensuremath{\varepsilon}}
\newunicodechar{Θ}{\ensuremath{\Theta}}
\newunicodechar{Λ}{\ensuremath{\Lambda}}
\newunicodechar{Μ}{\ensuremath{\mu}}
\newunicodechar{Τ}{\ensuremath{\tau}}
\newunicodechar{Φ}{\ensuremath{\Phi}}
\newunicodechar{Ψ}{\ensuremath{\Psi}}
\newunicodechar{Ω}{\ensuremath{\Omega}}
\newunicodechar{↦}{\ensuremath{\mapsto}}
\newunicodechar{∈}{\ensuremath{\in}}
\newunicodechar{∉}{\ensuremath{\notin}}
\newunicodechar{∧}{\ensuremath{\wedge}}
\newunicodechar{⇔}{\ensuremath{\Leftrightarrow}}
\newunicodechar{⟺}{\ensuremath{\Longleftrightarrow}}
\newunicodechar{⇒}{\ensuremath{\Rightarrow}}
\newunicodechar{⟹}{\ensuremath{\Longrightarrow}}
\newunicodechar{∘}{\ensuremath{\circ}}
\newunicodechar{∪}{\ensuremath{\cup}}
\newunicodechar{∩}{\ensuremath{\cap}}
\newunicodechar{≡}{\ensuremath{\equiv}}
\newunicodechar{⊂}{\ensuremath{\subset}}
\newunicodechar{⊥}{\ensuremath{\perp}}
\newunicodechar{∃}{\ensuremath{\exists}}
\newunicodechar{∀}{\ensuremath{\forall}}
\newunicodechar{∛}{\ensuremath{\sqrt[3]{\,}}}
\newunicodechar{∜}{\ensuremath{\sqrt[4]{\,}}}
\newunicodechar{⃗}{\ensuremath{{}^{\to}}}
\newunicodechar{ⁿ}{\ifmmode^{n}\else\textsuperscript{\ensuremath{n}}\fi}
\newunicodechar{ˣ}{\ifmmode^{x}\else\textsuperscript{\ensuremath{x}}\fi}
\newunicodechar{ᵇ}{\ifmmode^{b}\else\textsuperscript{\ensuremath{b}}\fi}
\newunicodechar{ʳ}{\ifmmode^{r}\else\textsuperscript{\ensuremath{r}}\fi}
\newunicodechar{ᵘ}{\ifmmode^{u}\else\textsuperscript{\ensuremath{u}}\fi}
\newunicodechar{ʸ}{\ifmmode^{y}\else\textsuperscript{\ensuremath{y}}\fi}
\newunicodechar{ᵃ}{\ifmmode^{a}\else\textsuperscript{\ensuremath{a}}\fi}
\newunicodechar{ᵉ}{\ifmmode^{e}\else\textsuperscript{\ensuremath{e}}\fi}
\newunicodechar{ᵏ}{\ifmmode^{k}\else\textsuperscript{\ensuremath{k}}\fi}
\newunicodechar{ᵖ}{\ifmmode^{p}\else\textsuperscript{\ensuremath{p}}\fi}
\newunicodechar{ᴺ}{\ifmmode^{N}\else\textsuperscript{\ensuremath{N}}\fi}
\newunicodechar{ᶜ}{\ifmmode^{c}\else\textsuperscript{\ensuremath{c}}\fi}
\newunicodechar{ʰ}{\ifmmode^{h}\else\textsuperscript{\ensuremath{h}}\fi}
\newunicodechar{ᵠ}{\ifmmode^{q}\else\textsuperscript{\ensuremath{q}}\fi}
\newunicodechar{ₙ}{\ifmmode_{n}\else\textsubscript{\ensuremath{n}}\fi}
\newunicodechar{ₐ}{\ifmmode_{a}\else\textsubscript{\ensuremath{a}}\fi}
\newunicodechar{ᵢ}{\ifmmode_{i}\else\textsubscript{\ensuremath{i}}\fi}
\newunicodechar{ᵣ}{\ifmmode_{r}\else\textsubscript{\ensuremath{r}}\fi}
\newunicodechar{ₖ}{\ifmmode_{k}\else\textsubscript{\ensuremath{k}}\fi}
\newunicodechar{ᵦ}{\ifmmode_{\beta}\else\textsubscript{\ensuremath{\beta}}\fi}
\newunicodechar{ₓ}{\ifmmode_{x}\else\textsubscript{\ensuremath{x}}\fi}
\newfontfamily\popdisplay{ArchivoBlack-Regular.ttf}[Path=../../../fonts/]
\newfontfamily\poplogo{SpaceGrotesk-Bold.ttf}[Path=../../../fonts/]
\newfontfamily\popsemi{PlusJakartaSans-SemiBold.ttf}[Path=../../../fonts/]
\newfontfamily\popxbold{PlusJakartaSans-ExtraBold.ttf}[Path=../../../fonts/]
\newfontfamily\popxboldls{PlusJakartaSans-ExtraBold.ttf}[Path=../../../fonts/,LetterSpace=3.0]

\setlist[enumerate]{leftmargin=1.7em,itemsep=5pt,topsep=4pt}
\setlist[itemize]{leftmargin=1.7em,itemsep=4pt,topsep=4pt,label={\color{poporange}\textbullet}}
\setlength{\parindent}{0pt}
\setlength{\parskip}{5pt}
\renewcommand{\arraystretch}{1.25}

% pgfplots : style Pop par défaut
\pgfplotsset{
  every axis/.append style={axis line style={popink,line width=1pt},tick style={popink},
    grid style={popline,line width=0.5pt},label style={font=\small},tick label style={font=\footnotesize}},
  popcurve/.style={poporange,line width=1.8pt,no markers,smooth},
}
% Même style disponible dans un \draw TikZ simple (hors pgfplots)
\tikzset{popcurve/.style={poporange,line width=1.8pt}}

% ============================================================
% Pill / squiggle
% ============================================================
\newcommand{\poppill}[3]{%
  \tikz[baseline=-0.55ex]\node[draw=popink,line width=1.2pt,rounded corners=2.6pt,%
    fill=#2,inner xsep=6.5pt,inner ysep=3pt]{%
    \popxboldls\fontsize{7}{7}\selectfont\color{#3}\MakeUppercase{#1}};%
}
\newcommand{\popsquiggle}[1]{%
  \tikz[baseline=-0.4ex]{
    \draw[#1,line width=2.6pt,line cap=round]
      plot [smooth] coordinates {(0,0.09) (0.34,0.01) (0.68,0.21) (1.02,0.01) (1.36,0.10)};
  }%
}
% Mot-clé surligné (marqueur orange sous le texte)
\newcommand{\cle}[1]{{\popxbold\color{popdark}#1}}

% ============================================================
% En-tête / pied
% ============================================================
\pagestyle{fancy}
\fancyhf{}
\renewcommand{\headrulewidth}{0pt}
\renewcommand{\footrulewidth}{0pt}
\lhead{\raisebox{-0.15ex}{\poplogo\fontsize{13}{13}\selectfont\color{popink}OnBuch\color{poporange}+}}
\rhead{\poppill{\DOCMATIERE\ \textbullet\ \DOCNIVEAU\ \textbullet\ Leçon \DOCLECON}{white}{popink}}
\lfoot{\popxbold\fontsize{7.6}{7.6}\selectfont\color{popmuted}\MakeUppercase{Cours signé OnBuch+}\ \textbullet\ {\popsemi\DOCTITRE}}
\rfoot{\tikz[baseline=-0.6ex]\node[rounded corners=8.5pt,fill=popink,minimum height=17pt,minimum width=17pt,inner xsep=5pt,inner sep=0]{\popxbold\fontsize{8}{8}\selectfont\color{popcream}\thepage\,/\,\pageref*{LastPage}};}

% ============================================================
% Titres
% ============================================================
\newcommand{\chaptitre}[2]{%
  \begin{tcolorbox}[enhanced,colback=poporange,colframe=popink,boxrule=2.6pt,arc=11pt,
      drop shadow=popink,left=15pt,right=15pt,top=12pt,bottom=11pt,before skip=3pt,after skip=18pt]
    \poppill{#2}{popink}{popcream}\par\vspace{9pt}
    {\raggedright\hyphenpenalty=10000\exhyphenpenalty=10000\popdisplay\fontsize{22}{24}\selectfont\color{popcream}\MakeUppercase{#1}\par}
  \end{tcolorbox}
}
% Section numérotée : pastille orange avec le numéro + titre Archivo
\newcounter{popsec}
\newcommand{\coursec}[1]{%
  \stepcounter{popsec}\setcounter{popsub}{0}%
  \par\vspace{16pt}\noindent
  \begin{minipage}{\linewidth}
  \tikz[baseline=-0.7ex]\node[draw=popink,line width=1.6pt,fill=poporange,rounded corners=4pt,minimum size=22pt,inner sep=0]{\popdisplay\fontsize{12}{12}\selectfont\color{popcream}\thepopsec};%
  \hspace{8pt}{\popdisplay\fontsize{15}{17}\selectfont\color{popink}\MakeUppercase{#1}}\par
  \vspace{4pt}\noindent{\color{poporange}\rule{36pt}{3.6pt}}
  \end{minipage}\par\nopagebreak\vspace{8pt}
}
\newcounter{popsub}
\newcommand{\courssub}[1]{%
  \stepcounter{popsub}%
  \par\vspace{9pt}\noindent{\popxbold\fontsize{11.5}{13}\selectfont\color{popink}{\color{poporange}\thepopsec.\thepopsub}\hspace{6pt}#1}\par\nopagebreak\vspace{4pt}
}

% ============================================================
% Boîtes pédagogiques — même recette : fond blanc, bordure épaisse colorée,
% ombre dure encre, badge plein. Titre optionnel [#1].
% ============================================================
\tcbset{popbase/.style={breakable,enhanced,colback=white,boxrule=2.3pt,arc=9pt,
  shadow={3pt}{-3pt}{0pt}{popink},left=14pt,right=14pt,top=11pt,bottom=11pt,before skip=11pt,after skip=15pt,
  fontupper=\popsemi\fontsize{10.6}{14.5}\selectfont\color{popink},
  fontlower=\popsemi\fontsize{10.6}{14.5}\selectfont\color{popink}}}
% Coche dessinée (le glyphe ✓ n'existe pas dans les polices du document)
\newcommand{\popcheck}[1]{\tikz[baseline=-0.5ex]\draw[#1,line width=1.6pt,line cap=round,line join=round](-0.13,0.0)--(-0.04,-0.09)--(0.13,0.1);}
\newcommand{\popboxhead}[3]{\poppill{#1}{#2}{white}\ifx\relax#3\relax\else\hspace{6pt}{\popxbold\fontsize{10.6}{12}\selectfont\color{popink}#3}\fi\par\vspace{7pt}}

\newtcolorbox{aretenir}[1][]{popbase,colframe=popgreen,before upper={\popboxhead{À retenir}{popgreen}{#1}}}
\newtcolorbox{exemplebox}[1][]{popbase,colframe=poporange,before upper={\popboxhead{Exemple}{poporange}{#1}}}
\newtcolorbox{definition}[1][]{popbase,colframe=popblue,before upper={\popboxhead{Définition}{popblue}{#1}}}
\newtcolorbox{propriete}[1][]{popbase,colframe=poppurple,before upper={\popboxhead{Propriété}{poppurple}{#1}}}
\newtcolorbox{methode}[1][]{popbase,colframe=popgold,before upper={\popboxhead{Méthode}{popgold}{#1}}}
\newtcolorbox{attention}[1][]{popbase,colframe=poppink,before upper={\popboxhead{Attention}{poppink}{#1}}}
\newtcolorbox{experience}[1][]{popbase,colframe=popblue,colback=popblueL,before upper={\popboxhead{Expérience}{popblue}{#1}}}
% « Le savais-tu ? » : carte violette pleine (contexte, culture, Cameroun)
\newtcolorbox{savaistu}[1][]{popbase,colback=poppurple,colframe=popink,
  fontupper=\popsemi\fontsize{10.4}{14.2}\selectfont\color{white},
  before upper={\poppill{Le savais-tu\,?}{white}{poppurple}\ifx\relax#1\relax\else\hspace{6pt}{\popxbold\color{white}#1}\fi\par\vspace{7pt}}}
% Exercice résolu : énoncé en haut, \tcblower puis la solution sur fond crème
\newcounter{popexo}\newcounter{popexr}
\newtcolorbox{exoresolu}[1][]{popbase,colframe=poporange,colbacklower=popcream,
  segmentation style={popink,line width=1.2pt,dashed},
  before upper={\stepcounter{popexr}\popboxhead{Exercice résolu \thepopexr}{poporange}{#1}},
  before lower={\poppill{Solution}{popink}{popcream}\par\vspace{6pt}}}
% Exercice (non résolu en place) — la correction est dans la partie Corrigés
\newtcolorbox{exercice}[1][]{popbase,colframe=popink,
  before upper={\stepcounter{popexo}\popboxhead{Exercice \thepopexo}{popink}{#1}}}
% Corrigé
\newtcolorbox{corrige}[1][]{popbase,colframe=popgreen,colback=popgreenL,
  before upper={\popboxhead{Corrigé}{popgreen}{#1}}}
% Objectifs / prérequis / bilan
\newtcolorbox{objectifs}{popbase,colframe=popink,colback=poporangeL,
  before upper={\popboxhead{Objectifs de la leçon}{popink}{}}}
\newtcolorbox{prerequis}{popbase,colframe=popink,colback=popblueL,
  before upper={\popboxhead{Prérequis}{popblue}{}}}
\newtcolorbox{bilan}{popbase,colframe=popink,colback=white,boxrule=2.8pt,
  before upper={\popboxhead{Fiche bilan}{poporange}{L'essentiel à connaître pour l'examen}}}
% Figure encadrée avec légende
\newtcolorbox{popfigure}[1][]{enhanced,breakable=false,colback=white,colframe=popline,boxrule=1.4pt,arc=8pt,
  left=10pt,right=10pt,top=10pt,bottom=8pt,before skip=10pt,after skip=12pt,halign=center,
  lower separated=false,halign lower=center,
  fontlower=\popsemi\fontsize{9}{11.5}\selectfont\color{popmuted},
  #1}
\newcommand{\legende}[1]{\par\vspace{6pt}{\popsemi\fontsize{9}{11.5}\selectfont\color{popmuted}#1}}

% ============================================================
% Signature OnBuch+
% ============================================================
\newcommand{\poplogoinline}[1]{{\poplogo\fontsize{#1}{#1}\selectfont\color{popink}OnBuch\color{poporange}+}}
\newcommand{\signatureonbuch}{%
  \begin{tcolorbox}[enhanced,colback=popink,colframe=popink,boxrule=2.2pt,arc=10pt,
      drop shadow=poporange,left=14pt,right=14pt,top=10pt,bottom=10pt,before skip=0pt,after skip=0pt]
    \tikz[baseline=-0.7ex]\node[circle,fill=popgreen,draw=popcream,line width=1.3pt,minimum size=22pt,inner sep=0]{\popcheck{white}};%
    \hspace{9pt}\begin{minipage}[c]{0.62\linewidth}
      {\popxbold\fontsize{12}{13}\selectfont\color{popcream}Signé\ \textbullet\ {\poplogo OnBuch\color{poporange}+}}\par\vspace{2pt}
      {\raggedright\popsemi\fontsize{8}{10.5}\selectfont\color{popline}\MakeUppercase{Équipe pédagogique OnBuch+}\\\MakeUppercase{Conforme au programme officiel MINESEC}\par}
    \end{minipage}\hfill
    {\poplogo\fontsize{22}{22}\selectfont\color{popcream}OnBuch\color{poporange}+}
  \end{tcolorbox}%
}

% ============================================================
% Page de garde
% #1 titre · #2 matière · #3 niveau · #4 module · #5 n° leçon · #6 durée ·
% #7 description / objectifs (texte ou itemize)
% ============================================================
\newcommand{\pagegarde}[7]{%
  \thispagestyle{empty}
  \newgeometry{a4paper,margin=1.7cm,top=1.6cm,bottom=1.4cm}%
  \begin{tikzpicture}[remember picture,overlay]
    \fill[poporange!18] ([xshift=-3.2cm,yshift=-3.6cm]current page.north east) circle (4.6cm);
    \fill[poppurple!14] ([xshift=2.4cm,yshift=7.6cm]current page.south west) circle (3.8cm);
  \end{tikzpicture}%
  {\poplogo\fontsize{30}{30}\selectfont\color{popink}OnBuch\color{poporange}+}\hfill
  \raisebox{0.6ex}{\poppill{Cours signé \textbullet\ Premium}{popink}{popcream}}\par
  \vspace{1pt}\popsquiggle{poppurple}\par
  \vspace{8pt}
  \begin{tcolorbox}[enhanced,colback=poporange,colframe=popink,boxrule=2.8pt,arc=13pt,
      drop shadow=popink,left=18pt,right=18pt,top=12pt,bottom=13pt,after skip=10pt]
    \poppill{#2 \textbullet\ #3}{popink}{popcream}\hspace{5pt}\poppill{#4}{white}{popink}\par\vspace{8pt}
    {\popdisplay\fontsize{14}{14}\selectfont\color{popink}LEÇON #5}\par\vspace{4pt}
    {\raggedright\hyphenpenalty=10000\exhyphenpenalty=10000\popdisplay\fontsize{25}{27}\selectfont\color{popcream}\MakeUppercase{#1}\par}
    \vspace{8pt}
    {\popxbold\fontsize{9.5}{9.5}\selectfont\color{popink}\MakeUppercase{Durée conseillée : #6}}
  \end{tcolorbox}
  \begin{tcolorbox}[enhanced,colback=white,colframe=popink,boxrule=2.2pt,arc=10pt,
      drop shadow=popink,left=16pt,right=16pt,top=10pt,bottom=10pt,after skip=8pt]
    \poppill{Dans cette leçon}{poppurple}{white}\par\vspace{5pt}
    \popsemi\fontsize{9.3}{12.6}\selectfont\color{popink}#7
  \end{tcolorbox}
  \noindent\begin{minipage}[t]{0.487\linewidth}
    \begin{tcolorbox}[enhanced,colback=popgreen,colframe=popink,boxrule=2.2pt,arc=10pt,
        drop shadow=popink,left=11pt,right=11pt,top=10pt,bottom=10pt,height=3.1cm,valign=center]
      \tcbox[colback=white,colframe=popink,boxrule=1.3pt,arc=5pt,boxsep=0pt,nobeforeafter,
        left=3pt,right=3pt,top=3pt,bottom=3pt,valign=center]{\qrcode[height=1.9cm]{https://play.google.com/store/apps/details?id=com.luvvix.onbuch&hl=fr}}%
      \hspace{8pt}\begin{minipage}[c]{0.5\linewidth}
        {\popdisplay\fontsize{11}{12.5}\selectfont\color{white}\MakeUppercase{Télécharge OnBuch}}\par\vspace{4pt}
        {\raggedright\popsemi\fontsize{8.4}{11}\selectfont\color{white}Gratuit \textbullet\ Google Play\\Tuteur IA, annales, concours, résultats\par}
      \end{minipage}
    \end{tcolorbox}
  \end{minipage}\hfill
  \begin{minipage}[t]{0.487\linewidth}
    \begin{tcolorbox}[enhanced,colback=poppurple,colframe=popink,boxrule=2.2pt,arc=10pt,
        drop shadow=popink,left=11pt,right=11pt,top=10pt,bottom=10pt,height=3.1cm,valign=center]
      \tikz[baseline=-0.7ex]\node[circle,fill=white,draw=popink,line width=1.3pt,minimum size=1.6cm,inner sep=0]{\popdisplay\fontsize{24}{24}\selectfont\color{poppurple}+};%
      \hspace{8pt}\begin{minipage}[c]{0.6\linewidth}
        {\popdisplay\fontsize{11}{12.5}\selectfont\color{white}\MakeUppercase{Passe à OnBuch+}}\par\vspace{4pt}
        {\raggedright\popsemi\fontsize{8.4}{11}\selectfont\color{white}Tous les cours signés, Tuteur IA illimité, fiches PDF — onglet OnBuch+ de l'app\par}
      \end{minipage}
    \end{tcolorbox}
  \end{minipage}
  \vspace{8pt}
  \popcontenu
  \vfill
  \signatureonbuch
  \restoregeometry
  \newpage
}

% Rangée « Ce cours contient » (page de garde)
\newcommand{\poptile}[3]{% #1 couleur #2 gros texte #3 légende
  \begin{tcolorbox}[enhanced,colback=#1,colframe=popink,boxrule=2pt,arc=9pt,shadow={2.5pt}{-2.5pt}{0pt}{popink},
      left=6pt,right=6pt,top=9pt,bottom=9pt,halign=center,valign=center,height=1.75cm,width=\linewidth,nobeforeafter]
    {\popdisplay\fontsize{12.5}{13}\selectfont\color{white}\MakeUppercase{#2}}\par\vspace{3pt}
    {\centering\popsemi\fontsize{7.8}{9.5}\selectfont\color{white}#3\par}
  \end{tcolorbox}}
\newcommand{\popcontenu}{%
  \par\noindent{\popxboldls\fontsize{7.5}{7.5}\selectfont\color{popmuted}CE COURS SIGNÉ CONTIENT}\par\vspace{7pt}
  \noindent\begin{minipage}[t]{0.235\linewidth}\poptile{popblue}{Cours}{complet, illustré, pas à pas}\end{minipage}\hfill
  \begin{minipage}[t]{0.235\linewidth}\poptile{poporange}{Méthodes}{et exercices résolus}\end{minipage}\hfill
  \begin{minipage}[t]{0.235\linewidth}\poptile{poppink}{Exercices}{type examen + corrigés}\end{minipage}\hfill
  \begin{minipage}[t]{0.235\linewidth}\poptile{popgreen}{Fiche bilan}{l'essentiel à retenir}\end{minipage}\par}

% Sommaire
\newtcolorbox{sommaire}{enhanced,colback=white,colframe=popink,boxrule=2.2pt,arc=10pt,
  drop shadow=popink,left=18pt,right=18pt,top=13pt,bottom=13pt,before skip=6pt,after skip=20pt,
  before upper={\poppill{Sommaire}{poppurple}{white}\par\vspace{9pt}\popsemi\fontsize{10.6}{14.5}\selectfont\color{popink}}}

% Signature de fin de cours — poussée tout en bas de la dernière page
\newcommand{\findecours}{%
  \par\vfill\nopagebreak
  \begin{center}\popsquiggle{poporange}\end{center}\vspace{4pt}
  \signatureonbuch
}
\AtBeginDocument{\def\llbracket{\mathopen{[\mkern-3.5mu[}}\def\rrbracket{\mathclose{]\mkern-3.5mu]}}} % intervalles d'entiers (glyphe ⟦ absent de la police)
% Glyphes absents de Fira Math : redéfinis à partir de glyphes disponibles
\AtBeginDocument{%
  \def\setminus{\mathbin{\backslash}}\let\smallsetminus\setminus
  \def\vdots{\mathord{\vbox{\baselineskip3.2pt\lineskiplimit0pt\kern4pt\hbox{.}\hbox{.}\hbox{.}}}}%
  \def\ddots{\mathinner{\mkern1mu\raise6.5pt\hbox{.}\mkern2mu\raise3.6pt\hbox{.}\mkern2mu\raise0.7pt\hbox{.}\mkern1mu}}%
  \def\triangle{\mathord{\scalebox{0.9}{$\symup{\Delta}$}}}%
  \def\nabla{\mathord{\raisebox{\depth}{\scalebox{1}[-1]{$\symup{\Delta}$}}}}%
  \def\checkmark{\popcheck{popdark}}%
  \def\star{\mathbin{\ast}}\def\bigstar{\mathord{\ast}}%
  \def\diamond{\mathbin{\scalebox{0.6}{$\Diamond$}}}%
  \def\bigcup{\mathop{\mathchoice{\raisebox{-0.3em}{\scalebox{1.7}{$\cup$}}}{\scalebox{1.25}{$\cup$}}{\cup}{\cup}}\displaylimits}%
  \def\bigcap{\mathop{\mathchoice{\raisebox{-0.3em}{\scalebox{1.7}{$\cap$}}}{\scalebox{1.25}{$\cap$}}{\cap}{\cap}}\displaylimits}%
  \def\lesseqqgtr{\mathrel{\lessgtr}}\def\bigoplus{\mathop{\oplus}\displaylimits}%
}
\providecommand{\ssi}{si et seulement si} % abréviation fréquente
\AtBeginDocument{\providecommand{\wideparen}{\overparen}} % arc de cercle

%% FIGFIX-BEGIN v4 — correctifs globaux des figures (docs/onbuch-generation/tools/figfix)
\usepackage{adjustbox}\usepackage{etoolbox}
% toute figure TikZ/pgfplots plus large que la ligne est réduite à la largeur disponible
\BeforeBeginEnvironment{tikzpicture}{\begin{adjustbox}{max width=\linewidth}}
\AfterEndEnvironment{tikzpicture}{\end{adjustbox}}
% légendes pgfplots lisibles par-dessus les courbes
\pgfplotsset{every axis/.append style={legend style={fill=white,fill opacity=0.92,text opacity=1,draw=black!15,inner sep=3pt}}}
% étiquettes d'arêtes (syntaxe edge["..."]) avec fond blanc
\tikzset{every edge quotes/.append style={fill=white,inner sep=1.2pt}}
%% FIGFIX-END
\begin{document}

\pagegarde{Limites et continuité d’une fonction}{Mathématiques}{1re Tech}{Mathématiques – classe de Première technique}{N13}{2 séances (fiche de progression harmonisée)}{Cette leçon introduit les notions fondamentales de limite d'une fonction en un point et de continuité, outils essentiels de l'analyse. Elle permet d'étudier le comportement local des fonctions, de justifier l'existence de solutions à des équations et de prolonger certaines fonctions, en lien direct avec des situations concrètes de la vie courante.
\vspace{4pt}
\begin{multicols}{2}\small
\begin{itemize}
\item À la fin de cette leçon, je sais déterminer graphiquement la limite d'une fonction en un point ou à l'infini.
\item À la fin de cette leçon, je sais calculer la limite en un point des fonctions usuelles (polynômes, fonctions rationnelles simples, valeur absolue).
\item À la fin de cette leçon, je sais définir et reconnaître une fonction continue en un point et sur un intervalle.
\item À la fin de cette leçon, je sais interpréter graphiquement la continuité (courbe tracée « sans lever le crayon ») et repérer une discontinuité.
\item À la fin de cette leçon, je sais utiliser les propriétés des fonctions continues (opérations, image d'un intervalle, théorème des valeurs intermédiaires).
\item À la fin de cette leçon, je sais justifier l'existence d'une solution à une équation du type f(x) = k sur un intervalle.
\end{itemize}\end{multicols}}

\begin{sommaire}
\begin{multicols}{2}
\begin{enumerate}
\item Limite d'une fonction en un point
\item Limites infinies et limites à l'infini (complément utile)
\item Continuité en un point et sur un intervalle
\item Propriétés des fonctions continues
\item Prolongement par continuité
\item Synthèse et applications à la vie courante
\item Activité d'intégration
\item Méthodes et exercices résolus
\item Exercices
\item Corrigés des exercices
\item Fiche bilan
\end{enumerate}
\end{multicols}
\end{sommaire}

\begin{objectifs}
\begin{itemize}
\item À la fin de cette leçon, je sais déterminer graphiquement la limite d'une fonction en un point ou à l'infini.
\item À la fin de cette leçon, je sais calculer la limite en un point des fonctions usuelles (polynômes, fonctions rationnelles simples, valeur absolue).
\item À la fin de cette leçon, je sais définir et reconnaître une fonction continue en un point et sur un intervalle.
\item À la fin de cette leçon, je sais interpréter graphiquement la continuité (courbe tracée « sans lever le crayon ») et repérer une discontinuité.
\item À la fin de cette leçon, je sais utiliser les propriétés des fonctions continues (opérations, image d'un intervalle, théorème des valeurs intermédiaires).
\item À la fin de cette leçon, je sais justifier l'existence d'une solution à une équation du type f(x) = k sur un intervalle.
\item À la fin de cette leçon, je sais déterminer si une fonction admet un prolongement par continuité en un point et le construire.
\item À la fin de cette leçon, je sais rédiger et justifier une démarche complète et appliquer ces notions à des situations concrètes de la vie courante.
\end{itemize}
\end{objectifs}

\begin{prerequis}
\begin{itemize}
\item Notion de fonction, ensemble de définition, image et antécédent (Seconde)
\item Représentation graphique et lecture graphique d'une fonction (Seconde)
\item Fonctions usuelles : fonctions affines, polynômes, fonction inverse, valeur absolue
\item Calcul algébrique : factorisation, développement, fractions
\item Intervalles de ℝ et inégalités (début d'année de Première)
\end{itemize}
\end{prerequis}

\newpage

\coursec{Limite d'une fonction en un point}

\courssub{Situation d'accroche et problématique}

Mama Ngo Bell tient son stand de beignets au marché Mokolo, à Douala. Son tarif est simple : 100 FCFA le beignet. Mais pour fidéliser sa clientèle, elle a une règle spéciale : \emph{à partir de 10 beignets achetés, elle offre une réduction forfaitaire de 200 FCFA}. Le prix total payé en fonction du nombre $x$ de beignets achetés est donc :
\[
p(x) = \begin{cases} 100x & \text{si } x < 10 \\ 100x - 200 & \text{si } x \geq 10 \end{cases}
\]
Un client qui achète 9 beignets paie 900 FCFA ; celui qui en achète 10 paie seulement 800 FCFA ! Le prix fait un \cle{saut} brutal en $x = 10$. Le graphe de cette fonction est-il \og en un seul morceau \fg{} ou présente-t-il une coupure ?

Autre image : un moto-taxi s'approche du carrefour Total à Yaoundé. Sa distance au carrefour se rapproche de zéro, de plus en plus, sans qu'on ait besoin de préciser l'instant exact de l'arrivée. Comment décrire mathématiquement ce comportement \og quand $x$ s'approche d'une valeur $a$, que devient $f(x)$ ? \fg{}

\textbf{Problématique :} comment définir et calculer rigoureusement la valeur vers laquelle $f(x)$ se rapproche lorsque $x$ se rapproche d'un nombre $a$, même lorsque $f(a)$ n'existe pas ?

\courssub{Approche intuitive de la limite}

Considérons la fonction $f$ définie pour tout réel $x \neq 1$ par
\[
f(x) = \frac{x^2 - 1}{x - 1}.
\]
Cette fonction n'est \textbf{pas définie en 1} : remplacer $x$ par 1 donnerait une division par zéro, ce qui est interdit. Pourtant, on peut s'interroger : \emph{que vaut $f(x)$ lorsque $x$ se rapproche de 1 ?} Calculons quelques valeurs :

\begin{center}
\begin{tabularx}{\linewidth}{|l|X|X|X|X|X|X|}
\hline
\rowcolor{popblueL} $x$ & 0,9 & 0,99 & 0,999 & 1,001 & 1,01 & 1,1 \\
\hline $f(x)$ & 1,9 & 1,99 & 1,999 & 2,001 & 2,01 & 2,1 \\
\hline
\end{tabularx}
\end{center}

On observe que plus $x$ est proche de 1 (par valeurs inférieures \emph{ou} supérieures), plus $f(x)$ est proche de 2. On dit que \og $f(x)$ tend vers 2 quand $x$ tend vers 1 \fg{}. Remarque essentielle : pour $x \neq 1$, on peut simplifier :
\[
f(x) = \frac{(x-1)(x+1)}{x-1} = x + 1,
\]
et quand $x$ s'approche de 1, $x+1$ s'approche bien de 2. La limite existe donc, \textbf{même si $f(1)$ n'existe pas}.

\begin{definition}[Limite finie en un point (définition intuitive)]
Soit $f$ une fonction définie au voisinage d'un réel $a$ (sauf éventuellement en $a$ lui-même) et $\ell$ un réel. On dit que $f(x)$ \cle{tend vers} $\ell$ quand $x$ tend vers $a$ si $f(x)$ devient aussi proche de $\ell$ que l'on veut, dès que $x$ est suffisamment proche de $a$. On note :
\[
\lim_{x \to a} f(x) = \ell.
\]
\end{definition}

\begin{attention}[Limite et valeur en $a$ : deux choses différentes]
Erreur fréquente : confondre $f(a)$ et $\displaystyle\lim_{x \to a} f(x)$.
\begin{itemize}
\item La \textbf{limite} décrit le comportement de $f(x)$ quand $x$ \emph{s'approche} de $a$, sans jamais l'atteindre.
\item $f(a)$ est la \textbf{valeur} de la fonction exactement en $a$ : elle peut exister ou non, et même être différente de la limite.
\end{itemize}
Dans l'exemple ci-dessus, $\displaystyle\lim_{x \to 1} f(x) = 2$ alors que $f(1)$ n'existe pas. La limite peut donc exister même si la fonction n'est pas définie en $a$.
\end{attention}

\courssub{Lecture graphique d'une limite}

Sur une courbe représentative, lire $\displaystyle\lim_{x \to a} f(x)$ revient à observer vers quelle ordonnée les points de la courbe se rapprochent lorsque leurs abscisses se rapprochent de $a$. Si la courbe a un \og trou \fg{} en $a$ (point creux), la limite est l'ordonnée de ce trou.

\begin{popfigure}
\begin{center}
\begin{tikzpicture}
\begin{axis}[
    width=0.85\linewidth, height=6.5cm,
    axis lines=middle,
    xlabel={$x$}, ylabel={$y$},
    xmin=-1.5, xmax=3.5, ymin=-1, ymax=4.5,
    xtick={1}, ytick={2},
    grid=both, grid style={popline!40},
    clip=false
]
\addplot[popcurve,domain=-1.2:3.2,samples=100]{x+1};
\addplot[only marks,mark=*,mark size=2.5pt,color=popblue,fill=white] coordinates {(1,2)};
\draw[dashed,poporange] (axis cs:1,0) -- (axis cs:1,2) -- (axis cs:0,2);
\node[popdark,anchor=south west] at (axis cs:1.05,2.05) {\scriptsize point creux $(1\,;\,2)$};
\end{axis}
\end{tikzpicture}
\end{center}
\legende{Courbe de $f(x)=\dfrac{x^2-1}{x-1}$ : c'est la droite d'équation $y=x+1$, privée du point $(1\,;\,2)$. La limite en 1 vaut 2, bien que $f(1)$ n'existe pas.}
\end{popfigure}

\courssub{Limites des fonctions usuelles en un point}

\begin{propriete}[Limites des fonctions usuelles — admises]
\begin{itemize}
\item Toute \cle{fonction polynôme} $P$ vérifie : $\displaystyle\lim_{x \to a} P(x) = P(a)$ pour tout réel $a$. En particulier $\displaystyle\lim_{x \to a} x = a$ et $\displaystyle\lim_{x \to a} k = k$ (fonction constante).
\item La \cle{fonction inverse} : pour tout $a \neq 0$, $\displaystyle\lim_{x \to a} \frac{1}{x} = \frac{1}{a}$.
\item La \cle{fonction valeur absolue} : pour tout réel $a$, $\displaystyle\lim_{x \to a} |x| = |a|$.
\end{itemize}
Autrement dit, pour ces fonctions usuelles, la limite en $a$ se calcule par simple \textbf{substitution} de $x$ par $a$ (lorsque $a$ est dans l'ensemble de définition).
\end{propriete}

\begin{exemplebox}[Calculs directs par substitution]
\begin{itemize}
\item $\displaystyle\lim_{x \to 3} (2x^2 - 5x + 1) = 2(3)^2 - 5(3) + 1 = 18 - 15 + 1 = 4$.
\item $\displaystyle\lim_{x \to -2} \frac{1}{x} = -\frac{1}{2}$.
\item $\displaystyle\lim_{x \to -5} |x| = |-5| = 5$.
\end{itemize}
\end{exemplebox}

\courssub{Limite à gauche, limite à droite}

Il arrive que le comportement de $f(x)$ dépende du \emph{côté} par lequel $x$ s'approche de $a$.

\begin{definition}[Limites latérales]
\begin{itemize}
\item On dit que $f(x)$ tend vers $\ell_1$ quand $x$ tend vers $a$ \cle{par valeurs inférieures} (à gauche) si $f(x)$ se rapproche de $\ell_1$ quand $x$ se rapproche de $a$ avec $x < a$. On note $\displaystyle\lim_{\substack{x \to a \\ x < a}} f(x) = \ell_1$ ou $\displaystyle\lim_{x \to a^-} f(x) = \ell_1$.
\item On définit de même la limite \cle{à droite} ($x > a$), notée $\displaystyle\lim_{x \to a^+} f(x)$.
\end{itemize}
La limite en $a$ existe si et seulement si les limites à gauche et à droite existent et sont \textbf{égales}.
\end{definition}

Reprenons le tarif de Mama Ngo Bell : $p(x) = 100x$ si $x < 10$ et $p(x) = 100x - 200$ si $x \geq 10$.
\[
\lim_{x \to 10^-} p(x) = 100 \times 10 = 1000 \quad ; \quad \lim_{x \to 10^+} p(x) = 100 \times 10 - 200 = 800.
\]
Les deux limites latérales sont \textbf{différentes} : la fonction $p$ n'a \textbf{pas de limite} en 10. Graphiquement, la courbe présente un \og saut \fg{} en $x = 10$.

\begin{popfigure}
\begin{center}
\begin{tikzpicture}
\begin{axis}[
    width=0.85\linewidth, height=7cm,
    axis lines=middle,
    xlabel={$x$ (beignets)}, ylabel={$p(x)$ (FCFA)},
    xmin=-0.5, xmax=15, ymin=-100, ymax=1600,
    xtick={5,10}, ytick={500,800,1000},
    grid=both, grid style={popline!40},
    clip=false
]
\addplot[popcurve,domain=0:9.99,samples=50]{100*x};
\addplot[popcurve,domain=10:15,samples=50,color=poporange,very thick]{100*x-200};
\addplot[only marks,mark=*,mark size=2.5pt,color=popblue,fill=white] coordinates {(10,1000)};
\addplot[only marks,mark=*,mark size=2.5pt,color=poporange] coordinates {(10,800)};
\draw[dashed,popmuted] (axis cs:10,0) -- (axis cs:10,1000);
\node[popdark,anchor=west] at (axis cs:10.3,1000) {\scriptsize point creux $(10\,;\,1000)$};
\node[popdark,anchor=west] at (axis cs:10.3,800) {\scriptsize point plein $(10\,;\,800)$};
\end{axis}
\end{tikzpicture}
\end{center}
\legende{Tarif par paliers des beignets : limite à gauche 1000, limite à droite 800 en $x=10$. Les limites latérales diffèrent : pas de limite en 10.}
\end{popfigure}

\begin{exemplebox}[Fonction partie entière]
La \cle{fonction partie entière}, notée $E(x)$, associe à tout réel $x$ le plus grand entier inférieur ou égal à $x$. Par exemple $E(2,7) = 2$ et $E(3) = 3$. Alors :
\[
\lim_{x \to 3^-} E(x) = 2 \quad \text{et} \quad \lim_{x \to 3^+} E(x) = 3.
\]
Les limites latérales diffèrent : $E$ n'a pas de limite en 3. C'est le cas de toutes les fonctions \og en escaliers \fg{} (tarifs postaux, taxes par tranches).
\end{exemplebox}

\courssub{Opérations sur les limites finies}

\begin{propriete}[Opérations sur les limites — admise]
Soient $f$ et $g$ deux fonctions telles que $\displaystyle\lim_{x \to a} f(x) = \ell$ et $\displaystyle\lim_{x \to a} g(x) = \ell'$ (limites finies). Alors :
\begin{itemize}
\item \textbf{Somme :} $\displaystyle\lim_{x \to a} \big(f(x) + g(x)\big) = \ell + \ell'$ ;
\item \textbf{Produit :} $\displaystyle\lim_{x \to a} \big(f(x) \times g(x)\big) = \ell \times \ell'$ ;
\item \textbf{Quotient :} si $\ell' \neq 0$, $\displaystyle\lim_{x \to a} \frac{f(x)}{g(x)} = \frac{\ell}{\ell'}$.
\end{itemize}
\end{propriete}

\begin{exemplebox}[Application des opérations]
Calculons $\displaystyle\lim_{x \to 2} \frac{3x + 1}{x + 4}$.
Le numérateur tend vers $3 \times 2 + 1 = 7$ et le dénominateur vers $2 + 4 = 6 \neq 0$. Donc :
\[
\lim_{x \to 2} \frac{3x + 1}{x + 4} = \frac{7}{6}.
\]
\end{exemplebox}

\begin{attention}[Condition sur le dénominateur]
La règle du quotient exige que la limite du dénominateur soit \textbf{non nulle}. Si le dénominateur tend vers 0, on ne peut pas appliquer directement la propriété : on obtient une forme dite \emph{indéterminée} qu'il faut traiter autrement (voir ci-dessous).
\end{attention}

\courssub{Lever une indétermination \og 0/0 \fg{} par factorisation}

Que se passe-t-il si, en remplaçant $x$ par $a$, on obtient $\dfrac{0}{0}$ ? Ce n'est \textbf{pas} un résultat : c'est une \cle{forme indéterminée}. Cela signifie que le calcul direct échoue et qu'il faut transformer l'expression.

\begin{methode}[Lever une indétermination 0/0 en $a$]
\begin{enumerate}
\item Vérifier qu'en remplaçant $x$ par $a$, numérateur et dénominateur valent tous deux 0.
\item \textbf{Factoriser} $(x - a)$ au numérateur et au dénominateur (identités remarquables, factorisation de polynômes).
\item \textbf{Simplifier} par $(x - a)$, ce qui est licite car $x \neq a$ dans le calcul de limite.
\item Calculer la limite de l'expression simplifiée par substitution.
\end{enumerate}
\end{methode}

\begin{exoresolu}[Limite de $\dfrac{x^2 - 4}{x - 2}$ en 2]
Calculer $\displaystyle\lim_{x \to 2} \frac{x^2 - 4}{x - 2}$.
\tcblower
\textbf{Étape 1 :} en remplaçant $x$ par 2 : numérateur $= 4 - 4 = 0$, dénominateur $= 0$. Forme indéterminée $\dfrac{0}{0}$.

\textbf{Étape 2 :} on factorise le numérateur grâce à l'identité remarquable $a^2 - b^2 = (a-b)(a+b)$ :
\[
x^2 - 4 = (x - 2)(x + 2).
\]
\textbf{Étape 3 :} pour $x \neq 2$, on simplifie :
\[
\frac{x^2 - 4}{x - 2} = \frac{(x-2)(x+2)}{x-2} = x + 2.
\]
\textbf{Étape 4 :} on conclut :
\[
\lim_{x \to 2} \frac{x^2 - 4}{x - 2} = \lim_{x \to 2} (x + 2) = 4.
\]
\end{exoresolu}

\begin{exercice}[À toi de jouer]
\begin{enumerate}
\item Calculer $\displaystyle\lim_{x \to 3} \frac{x^2 - 9}{x - 3}$.
\item Calculer $\displaystyle\lim_{x \to -1} \frac{x^2 + 2x + 1}{x + 1}$.
\item La fonction $g$ définie par $g(x) = 2x$ si $x < 1$ et $g(x) = 5 - x$ si $x \geq 1$ admet-elle une limite en 1 ? Justifier.
\end{enumerate}
\end{exercice}

\begin{corrige}[Exercice 1]
\begin{enumerate}
\item Forme $\dfrac{0}{0}$. On factorise : $x^2 - 9 = (x-3)(x+3)$. Pour $x \neq 3$ : $\dfrac{x^2-9}{x-3} = x + 3$, donc la limite vaut $3 + 3 = 6$.
\item Forme $\dfrac{0}{0}$. On reconnaît $x^2 + 2x + 1 = (x+1)^2$. Pour $x \neq -1$ : $\dfrac{(x+1)^2}{x+1} = x + 1$, donc la limite vaut $-1 + 1 = 0$.
\item Limite à gauche : $\displaystyle\lim_{x \to 1^-} g(x) = 2 \times 1 = 2$. Limite à droite : $\displaystyle\lim_{x \to 1^+} g(x) = 5 - 1 = 4$. Comme $2 \neq 4$, $g$ \textbf{n'admet pas de limite} en 1.
\end{enumerate}
\end{corrige}

\begin{aretenir}[L'essentiel de la section]
\begin{itemize}
\item $\displaystyle\lim_{x \to a} f(x) = \ell$ signifie : $f(x)$ devient aussi proche de $\ell$ qu'on veut quand $x$ est suffisamment proche de $a$.
\item La limite en $a$ est indépendante de $f(a)$ : elle peut exister même si $f(a)$ n'existe pas.
\item La limite existe si et seulement si les limites à gauche et à droite sont égales.
\item Pour les fonctions usuelles, on calcule par substitution ; les limites se combinent par somme, produit et quotient (dénominateur non nul).
\item Face à une forme $\dfrac{0}{0}$ : factoriser $(x-a)$, simplifier, puis substituer.
\end{itemize}
\end{aretenir}

\begin{savaistu}[Weierstrass, le père de la rigueur]
La notion de limite, utilisée de façon intuitive par Newton et Leibniz au XVII\textsuperscript{e} siècle pour inventer le calcul différentiel, n'a été formalisée rigoureusement qu'au XIX\textsuperscript{e} siècle, notamment par le mathématicien allemand Karl Weierstrass (1815--1897). Sa définition précise est le fondement de toute l'analyse moderne : dérivées, intégrales, études de fonctions. Sans limites, pas de vitesses instantanées, pas d'optimisation, pas de modélisation des phénomènes continus qui régissent la mécanique, l'électricité ou l'économie !
\end{savaistu}

\coursec{Limites infinies et limites à l'infini (complément utile)}

Dans la section précédente, nous avons appris à reconnaître la \cle{limite finie} d'une fonction en un point $a$ : lorsque $x$ se rapproche de $a$, les valeurs $f(x)$ se rapprochent d'un nombre réel $\ell$. Mais que se passe-t-il lorsque, au contraire, les valeurs $f(x)$ \emph{explosent} et deviennent de plus en plus grandes ? Et que devient une fonction lorsque $x$ lui-même devient très grand ? C'est l'objet de cette section, traitée de manière \emph{intuitive et graphique} : elle prépare les notions de \cle{Terminale} (asymptotes, tableaux de variations complets) sans exiger de formalisation rigoureuse.

\courssub{Limite infinie en un point}

\textbf{Intuition.} Imagine un technicien qui mesure l'intensité $I = \dfrac{U}{R}$ dans un circuit. Si la résistance $R$ se rapproche de zéro (court-circuit), l'intensité devient énorme : c'est exactement ce phénomène que décrit une limite infinie. La fonction $\dfrac{1}{R}$ ne se rapproche d'aucun nombre fini : elle \emph{dépasse toutes les bornes}.

\begin{definition}[Limite infinie en un point (définition intuitive)]
On dit que $f(x)$ tend vers $+\infty$ quand $x$ tend vers $a$, et on note
\[ \lim_{x \to a} f(x) = +\infty, \]
lorsque $f(x)$ devient \cle{aussi grand que l'on veut} dès que $x$ est suffisamment proche de $a$ (avec $x \neq a$).

De même, $\lim\limits_{x \to a} f(x) = -\infty$ signifie que $f(x)$ devient aussi grand que l'on veut \emph{en négatif} (c'est-à-dire inférieur à tout nombre négatif donné) lorsque $x$ est proche de $a$.
\end{definition}

\begin{exemplebox}[La fonction $x \mapsto \dfrac{1}{x^2}$ en $0$]
Calculons quelques valeurs :
\begin{center}
\begin{tabularx}{\linewidth}{|l|X|X|X|X|X|}
\hline
\rowcolor{popblueL} $x$ & $1$ & $0{,}1$ & $0{,}01$ & $-0{,}001$ & $10^{-6}$ \\
\hline $\dfrac{1}{x^2}$ & $1$ & $100$ & $10\,000$ & $10^{6}$ & $10^{12}$ \\
\hline
\end{tabularx}
\end{center}
Que $x$ s'approche de $0$ par la droite ou par la gauche, le carré $x^2$ est positif et de plus en plus petit, donc $\dfrac{1}{x^2}$ devient de plus en plus grand. On écrit :
\[ \lim_{x \to 0} \frac{1}{x^2} = +\infty. \]
Ici, les deux « côtés » donnent le même résultat : $+\infty$.
\end{exemplebox}

\courssub{Limites à gauche et à droite : le cas de la fonction inverse}

Pour la fonction \cle{inverse} $x \mapsto \dfrac{1}{x}$, la situation est plus subtile : le signe de $\dfrac{1}{x}$ dépend du signe de $x$. On distingue donc deux approches de $0$ :

\begin{itemize}
\item $x \to 0^{+}$ : $x$ s'approche de $0$ par \emph{valeurs supérieures} ($x > 0$) ;
\item $x \to 0^{-}$ : $x$ s'approche de $0$ par \emph{valeurs inférieures} ($x < 0$).
\end{itemize}

\begin{exemplebox}[Les deux limites de $\dfrac{1}{x}$ en $0$ diffèrent]
Par valeurs supérieures : si $x = 0{,}1$ alors $\dfrac{1}{x} = 10$ ; si $x = 0{,}001$ alors $\dfrac{1}{x} = 1000$ ; si $x = 10^{-9}$ alors $\dfrac{1}{x} = 10^{9}$. Les valeurs deviennent aussi grandes que l'on veut :
\[ \lim_{x \to 0^{+}} \frac{1}{x} = +\infty. \]
Par valeurs inférieures : si $x = -0{,}1$ alors $\dfrac{1}{x} = -10$ ; si $x = -0{,}001$ alors $\dfrac{1}{x} = -1000$ ; si $x = -10^{-9}$ alors $\dfrac{1}{x} = -10^{9}$. Les valeurs deviennent négatives et aussi grandes que l'on veut en valeur absolue :
\[ \lim_{x \to 0^{-}} \frac{1}{x} = -\infty. \]
\textbf{Conclusion :} les deux limites sont différentes ($+\infty$ à droite, $-\infty$ à gauche). On dit alors que la fonction inverse \emph{n'a pas de limite en $0$} au sens global : il faut préciser de quel côté on s'approche.
\end{exemplebox}

\begin{popfigure}
\begin{center}
\begin{tikzpicture}
\begin{axis}[
    width=0.85\linewidth, height=7cm,
    axis lines=middle,
    xlabel={$x$}, ylabel={$y$},
    xmin=-4, xmax=4, ymin=-5, ymax=5,
    xtick={-3,-2,-1,1,2,3}, ytick={-4,-2,2,4},
    grid=both, grid style={popline!40},
    legend style={at={(0.98,0.05)}, anchor=south east, font=\small}
]
\addplot[popcurve, domain=0.21:4, samples=200] {1/x};
\addplot[popcurve, domain=-4:-0.21, samples=200] {1/x};
\addplot[poporange, dashed, thick] coordinates {(0,-5) (0,5)};
\legend{$y=\dfrac{1}{x}$, , asymptote $x=0$}
\end{axis}
\end{tikzpicture}
\end{center}
\legende{Courbe de la fonction inverse : quand $x \to 0^{+}$, la courbe monte vers $+\infty$ ; quand $x \to 0^{-}$, elle descend vers $-\infty$. Aux extrémités ($x \to \pm\infty$), elle se rapproche de l'axe des abscisses.}
\end{popfigure}

\begin{attention}[Ne jamais écrire « $f(0) = +\infty$ »]
Une limite infinie n'est \textbf{pas} une valeur de la fonction. Écrire $f(0) = +\infty$ est une double erreur :
\begin{enumerate}
\item $+\infty$ n'est pas un nombre réel, donc ce n'est pas une image possible ;
\item la fonction $x \mapsto \dfrac{1}{x}$ n'est tout simplement \emph{pas définie} en $0$ : $f(0)$ n'existe pas.
\end{enumerate}
La bonne rédaction est : $\lim\limits_{x \to 0^{+}} \dfrac{1}{x} = +\infty$. Retiens aussi qu'une limite infinie signifie que la fonction \cle{n'a pas de limite finie} en ce point : dire « la limite est $+\infty$ » est une façon commode de décrire le comportement, pas un nombre.
\end{attention}

\textbf{Interprétation graphique : asymptote verticale.} Sur la figure ci-dessus, la courbe de la fonction inverse se « colle » à la droite verticale d'équation $x = 0$ (l'axe des ordonnées) sans jamais la toucher. On dit que cette droite est une \cle{asymptote verticale} à la courbe. Ce vocabulaire est signalé ici à titre de \emph{complément} : il sera étudié en détail en Terminale.

\courssub{Limites quand $x$ tend vers $+\infty$ ou $-\infty$}

On s'intéresse maintenant non plus au comportement près d'un point, mais au comportement \emph{aux extrémités} : que devient $f(x)$ quand $x$ devient très grand (on note $x \to +\infty$) ou très grand en négatif ($x \to -\infty$) ?

\begin{definition}[Limite à l'infini (définition intuitive)]
\begin{itemize}
\item On dit que $f(x)$ tend vers $+\infty$ quand $x$ tend vers $+\infty$ si $f(x)$ devient aussi grand que l'on veut dès que $x$ est suffisamment grand. On note $\lim\limits_{x \to +\infty} f(x) = +\infty$.
\item On dit que $f(x)$ tend vers un réel $\ell$ quand $x$ tend vers $+\infty$ si $f(x)$ se rapproche aussi près que l'on veut de $\ell$ dès que $x$ est suffisamment grand. On note $\lim\limits_{x \to +\infty} f(x) = \ell$.
\end{itemize}
On définit de même les limites quand $x \to -\infty$.
\end{definition}

\begin{exemplebox}[Deux exemples fondamentaux]
\textbf{1. La fonction carré $x \mapsto x^2$.} Si $x = 10$, $x^2 = 100$ ; si $x = 1000$, $x^2 = 10^{6}$ ; si $x = -10^{4}$, $x^2 = 10^{8}$. Le carré devient aussi grand que l'on veut, des deux côtés :
\[ \lim_{x \to +\infty} x^2 = +\infty \qquad \text{et} \qquad \lim_{x \to -\infty} x^2 = +\infty. \]

\textbf{2. La fonction inverse $x \mapsto \dfrac{1}{x}$.} Si $x = 1000$, $\dfrac{1}{x} = 0{,}001$ ; si $x = 10^{6}$, $\dfrac{1}{x} = 10^{-6}$ ; si $x = -10^{6}$, $\dfrac{1}{x} = -10^{-6}$. Les valeurs se rapprochent de $0$ :
\[ \lim_{x \to +\infty} \frac{1}{x} = 0 \qquad \text{et} \qquad \lim_{x \to -\infty} \frac{1}{x} = 0. \]
\end{exemplebox}

\begin{popfigure}
\begin{center}
\begin{tikzpicture}
\begin{axis}[
    width=0.8\linewidth, height=7cm,
    axis lines=middle,
    xlabel={$x$}, ylabel={$y$},
    xmin=-3.5, xmax=3.5, ymin=-1, ymax=9,
    xtick={-3,-2,-1,1,2,3}, ytick={2,4,6,8},
    grid=both, grid style={popline!40}
]
\addplot[popcurve, domain=-3:3, samples=200] {x^2};
\end{axis}
\end{tikzpicture}
\end{center}
\legende{Courbe de $x \mapsto x^2$ : les deux branches montent indéfiniment quand $x \to +\infty$ et quand $x \to -\infty$ : les deux limites valent $+\infty$.}
\end{popfigure}

\textbf{Lecture graphique : comportement des branches.} Sur une courbe, les limites à l'infini se lisent en regardant \emph{les extrémités des branches} :
\begin{itemize}
\item si la branche \emph{monte} indéfiniment vers la droite, alors $\lim\limits_{x \to +\infty} f(x) = +\infty$ (cas de $x^2$) ;
\item si la branche \emph{descend} indéfiniment, la limite est $-\infty$ ;
\item si la branche se rapproche d'une \emph{droite horizontale} $y = \ell$, la limite est le nombre fini $\ell$ (cas de $\dfrac{1}{x}$, dont la courbe se colle à l'axe des abscisses : on parlera en Terminale d'\emph{asymptote horizontale}).
\end{itemize}

\begin{exoresolu}[Lecture graphique et calcul de limites]
On considère la fonction $f$ définie sur $\mathbb{R} \setminus \{0\}$ par $f(x) = \dfrac{1}{x^2} + 1$.
\begin{enumerate}
\item Déterminer $\lim\limits_{x \to 0} f(x)$.
\item Déterminer $\lim\limits_{x \to +\infty} f(x)$ et $\lim\limits_{x \to -\infty} f(x)$.
\end{enumerate}
\tcblower
\textbf{Solution détaillée.}
\begin{enumerate}
\item Quand $x \to 0$, on a vu que $\dfrac{1}{x^2} \to +\infty$ (des deux côtés, car $x^2 > 0$). Ajouter $1$ ne change rien à ce comportement : les valeurs deviennent aussi grandes que l'on veut. Donc
\[ \lim_{x \to 0} f(x) = +\infty. \]
La droite d'équation $x = 0$ est asymptote verticale à la courbe (vocabulaire de Terminale).
\item Quand $x \to +\infty$, $x^2$ devient très grand, donc $\dfrac{1}{x^2}$ se rapproche de $0$, et $f(x) = \dfrac{1}{x^2} + 1$ se rapproche de $1$ :
\[ \lim_{x \to +\infty} f(x) = 1. \]
Le même raisonnement vaut quand $x \to -\infty$ (car $x^2$ est le même pour $x$ et $-x$) :
\[ \lim_{x \to -\infty} f(x) = 1. \]
La courbe se rapproche de la droite horizontale $y = 1$ aux deux extrémités.
\end{enumerate}
\end{exoresolu}

\begin{attention}[Pièges fréquents]
\begin{itemize}
\item Ne pas confondre les deux situations : « $x \to 0$ » (on s'approche d'un point) et « $x \to +\infty$ » (on part vers les extrémités) donnent des résultats très différents pour la même fonction : $\lim\limits_{x \to 0^{+}} \dfrac{1}{x} = +\infty$ mais $\lim\limits_{x \to +\infty} \dfrac{1}{x} = 0$.
\item Ne pas oublier le signe : $\dfrac{1}{x}$ tend vers $+\infty$ à droite de $0$ mais vers $-\infty$ à gauche. Toujours vérifier le signe de la quantité qui devient petite.
\item Une limite infinie n'est pas une limite finie : on ne peut pas faire de calculs ordinaires avec $+\infty$ comme avec un nombre.
\end{itemize}
\end{attention}

\begin{exercice}[Comportements aux bornes]
Pour chacune des fonctions suivantes, déterminer intuitivement (à l'aide de quelques valeurs bien choisies) les limites demandées :
\begin{enumerate}
\item $f(x) = \dfrac{1}{x-2}$ : limites quand $x \to 2^{+}$ et quand $x \to 2^{-}$ ;
\item $g(x) = x^2 + 3$ : limites quand $x \to +\infty$ et quand $x \to -\infty$ ;
\item $h(x) = \dfrac{-1}{x^2}$ : limite quand $x \to 0$.
\end{enumerate}
\end{exercice}

\begin{corrige}[Exercice 1]
\begin{enumerate}
\item Quand $x \to 2^{+}$, $x - 2$ est positif et tend vers $0$, donc $\dfrac{1}{x-2} \to +\infty$. Quand $x \to 2^{-}$, $x - 2$ est négatif et tend vers $0$, donc $\dfrac{1}{x-2} \to -\infty$.
\item $x^2 \to +\infty$ dans les deux cas, donc $g(x) = x^2 + 3 \to +\infty$ quand $x \to +\infty$ et quand $x \to -\infty$.
\item $x^2 > 0$ tend vers $0$, donc $\dfrac{1}{x^2} \to +\infty$, et avec le signe moins : $h(x) \to -\infty$ quand $x \to 0$.
\end{enumerate}
\end{corrige}

\begin{savaistu}[Et en Terminale ?]
Toutes les notions vues ici de manière intuitive seront \emph{formalisées} en classe de Terminale : définitions rigoureuses des limites, théorèmes sur les opérations (somme, produit, quotient de limites), étude complète des \cle{asymptotes} (verticales et horizontales) et construction de \cle{tableaux de variations} où les limites apparaissent aux bornes du domaine. Ces outils serviront ensuite à étudier des fonctions concrètes : coût moyen de production en gestion, intensité dans un circuit, vitesse d'un mobile. Les physiciens et ingénieurs utilisent quotidiennement ces idées : par exemple, la loi d'Ohm $I = \dfrac{U}{R}$ montre qu'un court-circuit ($R \to 0$) fait « exploser » l'intensité — d'où l'importance des fusibles dans les installations électriques de nos maisons !
\end{savaistu}

\coursec{Continuité en un point et sur un intervalle}

\courssub{Transition : des limites à la continuité}
Dans les sections précédentes, nous avons étudié le comportement d'une fonction \emph{au voisinage} d'un point ou à l'infini, sans nécessairement nous soucier de la valeur prise \emph{en} ce point. La notion de \cle{continuité} vient combler ce vide : elle exige une parfaite adéquation entre la valeur de la fonction en un point et sa limite en ce point. En termes concrets, une fonction est continue si sa courbe se trace « sans lever le crayon ». Cette propriété est fondamentale en technique : elle garantit qu'une petite variation de la cause (entrée) produit une petite variation de l'effet (sortie), ce qui est essentiel pour la stabilité des systèmes (électronique, mécanique, régulation thermique).

\courssub{Définition de la continuité en un point}
Commençons par l'intuition : on dit qu'une fonction est continue en $a$ si, lorsque $x$ se rapproche de $a$, $f(x)$ se rapproche de $f(a)$. Autrement dit, il n'y a ni « trou », ni « saut », ni « asymptote verticale » en $a$.

\begin{definition}[Continuité en un point]
Soit $f$ une fonction définie sur un intervalle $I$ et $a \in I$.
On dit que $f$ est \textbf{continue en $a$} si et seulement si les trois conditions suivantes sont réunies :
\begin{enumerate}
    \item $f$ est \textbf{définie en $a$} ($a$ appartient à l'ensemble de définition $\mathcal{D}_f$).
    \item $f$ admet une \textbf{limite finie} en $a$ : $\displaystyle\lim_{x \to a} f(x) = \ell \in \mathbb{R}$.
    \item Cette limite est \textbf{égale à la valeur de la fonction} en $a$ : $\ell = f(a)$.
\end{enumerate}
En notation condensée :
\[
f \text{ continue en } a \quad\Longleftrightarrow\quad \lim_{x \to a} f(x) = f(a).
\]
Si l'une de ces conditions n'est pas vérifiée, on dit que $f$ est \textbf{discontinue} en $a$ (ou qu'elle présente une \cle{discontinuité} en $a$).
\end{definition}

\begin{attention}[Erreur fréquente : « Définie » ne veut pas dire « Continue »]
Beaucoup d'élèves confondent « $f$ est définie en $a$ » et « $f$ est continue en $a$ ».
\emph{Contre-exemple :} $f(x) = \begin{cases} x^2 & \text{si } x \neq 2 \\ 5 & \text{si } x = 2 \end{cases}$.
$f$ est bien définie en $2$ ($f(2)=5$), mais $\lim_{x\to 2} f(x) = 4 \neq 5$. La fonction n'est \textbf{pas} continue en $2$ (discontinuité « par trou » : le point $(2\,;\,5)$ est isolé).
\textbf{Règle d'or :} vérifier \textbf{systématiquement} l'égalité $\lim_{x\to a} f(x) = f(a)$, et non seulement l'existence de $f(a)$.
\end{attention}

\courssub{Interprétation graphique}
Graphiquement, la continuité en $a$ signifie que le point $A(a\,;\,f(a))$ appartient à la courbe représentative $\mathcal{C}_f$ et que la courbe ne présente aucune rupture au voisinage de $A$. On peut tracer la courbe autour de $A$ sans lever le crayon.

\begin{popfigure}
\begin{tikzpicture}
\begin{axis}[
    name=plot1,
    width=0.48\linewidth, height=6cm,
    axis lines=middle,
    xlabel=$x$, ylabel=$y$,
    title={Fonction continue : $f(x)=x^2$},
    title style={font=\small},
    xmin=-1, xmax=4, ymin=-1, ymax=10,
    xtick={0,1,2,3}, ytick={0,2,4,6,8},
    tick label style={font=\small},
    label style={font=\small}
]
\addplot[popcurve, domain=-1:3, samples=100] {x^2};
\addplot[only marks, mark=*, mark size=2, popblue] coordinates {(2,4)};
\node[above right, popblue, font=\small] at (axis cs:2,4) {$A(2\,;\,4)$};
\end{axis}
\begin{axis}[
    at={($(plot1.east)+(1.6cm,0)$)}, anchor=west,
    width=0.48\linewidth, height=6cm,
    axis lines=middle,
    xlabel=$x$, ylabel=$y$,
    title={Fonction discontinue en $2$ : saut},
    title style={font=\small},
    xmin=-1, xmax=4, ymin=-1, ymax=6,
    xtick={0,1,2,3}, ytick={0,1,2,3,4,5},
    tick label style={font=\small},
    label style={font=\small}
]
\addplot[popcurve, domain=-1:2, samples=50] {x+1};
\addplot[only marks, mark=o, mark size=2.5, popgreen, thick] coordinates {(2,3)};
\addplot[popcurve, domain=2:3.5, samples=50] {0.5*x+3};
\addplot[only marks, mark=*, mark size=2, poporange] coordinates {(2,4)};
\node[above left, popgreen, font=\small] at (axis cs:2,3) {limite à gauche $=3$};
\node[right, poporange, font=\small] at (axis cs:2,4) {$f(2)=4$};
\end{axis}
\end{tikzpicture}
\legende{À gauche : $f(x)=x^2$ est continue en $2$ (la limite est égale à la valeur). À droite : fonction par morceaux avec un saut en $2$ ; $\lim_{x\to 2^-}f(x)=3$ (point creux) alors que $f(2)=4$ (point plein). Les limites latérales étant différentes, la limite n'existe pas en $2$ : $f$ n'est pas continue en $2$.}
\end{popfigure}

\courssub{Exemples fondamentaux de fonctions continues}
Certaines fonctions « de base » sont continues sur tout leur ensemble de définition. Ce résultat est admis au Probatoire : on ne redémontre pas la continuité des polynômes à chaque exercice.

\begin{propriete}[Fonctions continues usuelles (admis)]
\begin{itemize}
    \item Toute \textbf{fonction polynôme} est continue sur $\mathbb{R}$.
    \item Toute \textbf{fonction rationnelle} $f(x) = \dfrac{P(x)}{Q(x)}$ (quotient de polynômes) est continue sur son ensemble de définition $\mathcal{D}_f = \{ x \in \mathbb{R} \mid Q(x) \neq 0 \}$.
    \item La fonction \textbf{valeur absolue} $x \mapsto |x|$ est continue sur $\mathbb{R}$.
    \item La fonction \textbf{racine carrée} $x \mapsto \sqrt{x}$ est continue sur $[0;+\infty[$.
    \item La fonction \textbf{partie entière} $x \mapsto E(x)$ est continue sur chaque intervalle $]n;n+1[$ ($n\in\mathbb{Z}$) et discontinue en tout nombre entier.
\end{itemize}
\end{propriete}

\begin{exemplebox}[La fonction valeur absolue : continue en 0]
La fonction $f(x)=|x|$ est définie par $f(x)=x$ si $x\ge 0$ et $f(x)=-x$ si $x<0$.
\begin{itemize}
    \item Pour $x>0$ : $f(x)=x$ (polynôme), donc $f$ est continue sur $]0;+\infty[$.
    \item Pour $x<0$ : $f(x)=-x$ (polynôme), donc $f$ est continue sur $]-\infty;0[$.
    \item En $x=0$ : $f(0)=0$. De plus $\lim_{x\to 0^-} f(x) = \lim_{x\to 0^-} (-x) = 0$ et $\lim_{x\to 0^+} f(x) = \lim_{x\to 0^+} x = 0$.
    Donc $\lim_{x\to 0} f(x) = 0 = f(0)$ : $f$ est \textbf{continue en $0$}.
\end{itemize}
Graphiquement, la courbe forme un « angle » (un pic) en $0$, mais il n'y a ni trou ni saut : la continuité est bien assurée.
\end{exemplebox}

\begin{popfigure}
\begin{tikzpicture}
\begin{axis}[
    width=0.7\linewidth, height=6cm,
    axis lines=middle,
    xlabel=$x$, ylabel=$y$,
    xmin=-3, xmax=3, ymin=-0.5, ymax=3,
    xtick={-2,-1,0,1,2}, ytick={0,1,2},
    tick label style={font=\small},
    label style={font=\small}
]
\addplot[popcurve, domain=-3:3, samples=100] {abs(x)};
\addplot[only marks, mark=*, mark size=2.5, popblue] coordinates {(0,0)};
\node[above right, popblue, font=\small] at (axis cs:0,0) {$O(0\,;\,0)$};
\node[above left, popmuted, font=\small] at (axis cs:-1.5,1.5) {$y=-x$};
\node[above right, popmuted, font=\small] at (axis cs:1.5,1.5) {$y=x$};
\end{axis}
\end{tikzpicture}
\legende{Courbe de $f(x)=|x|$. La fonction est continue sur $\mathbb{R}$ (aucune rupture), avec un angle en $0$.}
\end{popfigure}

\begin{exemplebox}[Fonction rationnelle : continuité sur l'ensemble de définition]
Soit $f(x) = \dfrac{x^2-1}{x-1}$.
Le dénominateur s'annule en $x=1$, donc $\mathcal{D}_f = \mathbb{R} \setminus \{1\}$.
$f$ est continue sur $]-\infty; 1[$ et sur $]1; +\infty[$ (fonction rationnelle sur son ensemble de définition).
En $x=1$, $f$ \textbf{n'est pas définie} : la question de la continuité en $1$ ne se pose pas directement (on parlera de \emph{prolongement par continuité} dans la section 5, car $\lim_{x\to 1}f(x)=2$).
\end{exemplebox}

\courssub{Méthode : étudier la continuité d'une fonction définie par morceaux}
C'est le cas typique de l'examen (Probatoire). La fonction change d'expression de part et d'autre d'une valeur $a$. Pour étudier la continuité en $a$, on compare les limites latérales (à gauche et à droite) et la valeur $f(a)$.

\begin{methode}[Continuité en $a$ pour une fonction par morceaux]
Soit $f$ définie par :
\[
f(x) = \begin{cases}
f_1(x) & \text{si } x < a \\
f_2(x) & \text{si } x \ge a
\end{cases}
\]
Pour étudier la continuité en $a$ :
\begin{enumerate}
    \item \textbf{Calculer $f(a)$} (valeur imposée par la définition, ici l'expression du second cas puisque $a \ge a$).
    \item \textbf{Calculer la limite à gauche :} $\ell_g = \lim_{x \to a^-} f(x) = \lim_{x \to a^-} f_1(x)$.
    \item \textbf{Calculer la limite à droite :} $\ell_d = \lim_{x \to a^+} f(x) = \lim_{x \to a^+} f_2(x)$.
    \item \textbf{Conclure :}
    \begin{itemize}
        \item Si $\ell_g = \ell_d = f(a)$ : $f$ est \textbf{continue} en $a$.
        \item Si $\ell_g \neq \ell_d$ : la limite n'existe pas en $a$, $f$ présente un \textbf{saut} ; elle n'est pas continue en $a$.
        \item Si $\ell_g = \ell_d \neq f(a)$ : la limite existe mais diffère de la valeur ; $f$ présente un \textbf{trou} ; elle n'est pas continue en $a$.
        \item Si l'une des limites est infinie : \textbf{asymptote verticale} d'équation $x=a$ ; $f$ n'est pas continue en $a$.
    \end{itemize}
\end{enumerate}
\end{methode}

\begin{exoresolu}[Continuité d'une fonction par morceaux]
\textbf{Énoncé :} Étudier la continuité en $1$ de la fonction $f$ définie sur $\mathbb{R}$ par :
\[
f(x) = \begin{cases}
2x + 1 & \text{si } x < 1 \\
4x - 1 & \text{si } x \ge 1
\end{cases}
\]
\tcblower
\textbf{Solution détaillée :}
\begin{enumerate}
    \item \textbf{Valeur en $1$ :} comme $1 \ge 1$, on utilise la deuxième expression : $f(1) = 4\times 1 - 1 = 3$.
    \item \textbf{Limite à gauche ($x \to 1^-$) :} pour $x<1$, $f(x)=2x+1$, donc
    $\ell_g = \lim_{x \to 1^-} (2x+1) = 2\times 1 + 1 = 3$.
    \item \textbf{Limite à droite ($x \to 1^+$) :} pour $x>1$, $f(x)=4x-1$, donc
    $\ell_d = \lim_{x \to 1^+} (4x-1) = 4\times 1 - 1 = 3$.
    \item \textbf{Conclusion :} $\ell_g = \ell_d = 3 = f(1)$.
    \textbf{La fonction $f$ est continue en $1$.}
\end{enumerate}
\emph{Remarque :} bien que les expressions changent, les deux branches « se rejoignent » parfaitement au point $(1\,;\,3)$. La courbe ne présente pas de saut.
\end{exoresolu}

\begin{exoresolu}[Discontinuité : saut de fonction]
\textbf{Énoncé :} Étudier la continuité en $0$ de la fonction $g$ définie sur $\mathbb{R}$ par :
\[
g(x) = \begin{cases} x+2 & \text{si } x < 0 \\ x^2+1 & \text{si } x \ge 0 \end{cases}
\]
\tcblower
\textbf{Solution :}
\begin{itemize}
    \item Valeur en $0$ : $g(0) = 0^2+1 = 1$.
    \item Limite à gauche : $\lim_{x\to 0^-} g(x) = \lim_{x\to 0^-} (x+2) = 2$.
    \item Limite à droite : $\lim_{x\to 0^+} g(x) = \lim_{x\to 0^+} (x^2+1) = 1$.
\end{itemize}
On obtient $\ell_g = 2 \neq 1 = \ell_d$. Les limites latérales sont différentes, donc la limite de $g$ en $0$ n'existe pas.
\textbf{$g$ n'est pas continue en $0$ : elle présente un saut d'amplitude $|2-1|=1$.}
\end{exoresolu}

\courssub{Continuité sur un intervalle}
La continuité en un point se généralise naturellement à un ensemble de points.

\begin{definition}[Continuité sur un intervalle]
Soit $I$ un intervalle de $\mathbb{R}$ (ouvert, fermé, demi-ouvert, borné ou non).
Une fonction $f$ est \textbf{continue sur $I$} si elle est continue \textbf{en tout point de $I$}.
\begin{itemize}
    \item Si $I = [a;b]$ (fermé) : $f$ est continue sur $]a;b[$, continue \textbf{à droite} en $a$ (c'est-à-dire $\lim_{x\to a^+}f(x)=f(a)$) et continue \textbf{à gauche} en $b$ (c'est-à-dire $\lim_{x\to b^-}f(x)=f(b)$).
    \item Si $I = ]a;b[$ (ouvert) : $f$ est continue en tout point de $]a;b[$.
\end{itemize}
\end{definition}

\begin{propriete}[Fonctions continues sur un intervalle]
\begin{itemize}
    \item Les fonctions polynômes sont continues sur $\mathbb{R}$, donc sur tout intervalle de $\mathbb{R}$.
    \item Les fonctions rationnelles sont continues sur tout intervalle \textbf{inclus dans leur ensemble de définition} (ne contenant aucun zéro du dénominateur).
    \item $x \mapsto \sqrt{x}$ est continue sur $[0;+\infty[$.
    \item $x \mapsto |x|$ est continue sur $\mathbb{R}$.
\end{itemize}
\end{propriete}

\begin{attention}[Piège : intervalle et ensemble de définition]
La fonction $f(x)=\dfrac{1}{x}$ est continue sur $]0;+\infty[$ et sur $]-\infty;0[$, mais \textbf{pas} sur $[-1;1]$ car $0 \notin \mathcal{D}_f$ : la fonction n'est même pas définie en $0$.
\textbf{Toujours vérifier que l'intervalle d'étude est inclus dans $\mathcal{D}_f$ avant d'affirmer la continuité sur cet intervalle.}
\end{attention}

\courssub{Applications à la vie courante et situations techniques}
La continuité n'est pas une abstraction pure ; elle modélise des phénomènes réels du quotidien et des métiers techniques.

\begin{savaistu}[Continuité et discontinuité dans la vie réelle]
\begin{center}
\begin{tabularx}{\linewidth}{|l|X|X|}\hline
\rowcolor{popblueL}
\textbf{Phénomène} & \textbf{Modèle mathématique} & \textbf{Type} \\ \hline
Température de l'air au cours d'une journée & $T(t)$ : température à l'instant $t$ & \textbf{Continue} : la température ne saute pas brutalement de \SI{20}{\degree C} à \SI{30}{\degree C} sans passer par les valeurs intermédiaires. \\ \hline
Hauteur d'eau dans un réservoir qui se remplit & $h(t)$ : hauteur à l'instant $t$ & \textbf{Continue} (si l'ouverture de la vanne est progressive). \\ \hline
Tarif d'envoi d'un colis selon le poids & $P(m)$ : prix pour une masse $m$ & \textbf{Discontinue (par paliers)} : par exemple 0--5 kg $\to$ 1000 F ; 5--10 kg $\to$ 1800 F. Saut en $m=5$. \\ \hline
Tension aux bornes d'un condensateur & $u(t)$ & \textbf{Continue} : la tension ne peut pas varier instantanément (une variation instantanée exigerait un courant infini). \\ \hline
Commande « tout-ou-rien » (thermostat, fin de course) & $c(t) \in \{0;1\}$ & \textbf{Discontinue} : basculement brutal entre 0 et 1. \\ \hline
\end{tabularx}
\end{center}
\textbf{Enjeu technique :} en régulation (automatique), on cherche souvent à rendre les systèmes continus pour éviter les à-coups mécaniques (usure, vibrations) ou les surtensions électriques. L'étude de la continuité permet de valider la faisabilité d'une loi de commande.
\end{savaistu}

\begin{exemplebox}[Modélisation d'un tarif d'électricité]
Un fournisseur propose : 50 F/kWh pour les 100 premiers kWh, puis 75 F/kWh au-delà.
Soit $C(x)$ le coût pour $x$ kWh consommés ($x \ge 0$) :
\[
C(x) = \begin{cases}
50x & \text{si } 0 \le x \le 100 \\
5000 + 75(x-100) = 75x - 2500 & \text{si } x > 100
\end{cases}
\]
Étudions la continuité en $100$ (changement de tranche).
\begin{itemize}
    \item Valeur : $C(100) = 50 \times 100 = 5000$ F.
    \item Limite à gauche : $\lim_{x\to 100^-} C(x) = 50\times 100 = 5000$.
    \item Limite à droite : $\lim_{x\to 100^+} C(x) = 75\times 100 - 2500 = 5000$.
\end{itemize}
Les limites latérales coïncident avec la valeur : \textbf{la fonction coût est continue en 100}. Il n'y a pas de « saut » de facture au passage de la tranche : le système de tarification est bien conçu.
\end{exemplebox}

\courssub{Synthèse de la section}
\begin{aretenir}[Ce qu'il faut retenir sur la continuité]
\begin{itemize}
    \item \textbf{Définition :} $f$ continue en $a \iff \lim_{x\to a} f(x) = f(a)$ (trois conditions : $f$ définie en $a$, limite finie en $a$, égalité des deux).
    \item \textbf{Interprétation graphique :} courbe sans rupture (on ne lève pas le crayon).
    \item \textbf{Fonctions continues sur leur ensemble de définition :} polynômes (sur $\mathbb{R}$), fonctions rationnelles (sur $\mathcal{D}_f$), valeur absolue (sur $\mathbb{R}$), racine carrée (sur $[0;+\infty[$).
    \item \textbf{Méthode pour une fonction par morceaux en $a$ :} comparer $\lim_{x\to a^-}f(x)$, $\lim_{x\to a^+}f(x)$ et $f(a)$.
    \item \textbf{Types de discontinuité :} saut (limites latérales finies différentes), trou (limite finie différente de $f(a)$), asymptote verticale (limite infinie).
    \item \textbf{Continuité sur un intervalle $I$ :} continuité en tout point de $I$ (attention aux bornes si $I$ est fermé : limites latérales adaptées).
\end{itemize}
\end{aretenir}

\begin{exercice}[Entraînement : continuité en un point]
\textbf{1.} Étudier la continuité en $0$ des fonctions suivantes :
\begin{itemize}
    \item[a)] $f(x) = \begin{cases} x^2+1 & \text{si } x<0 \\ 2 & \text{si } x=0 \\ \sqrt{x+1} & \text{si } x>0 \end{cases}$
    \item[b)] $g(x) = \dfrac{x^2-4}{x-2}$ en $2$ (on précisera d'abord $\mathcal{D}_g$).
\end{itemize}
\textbf{2.} Soit $h(x) = \begin{cases} 3x-2 & \text{si } x \le 1 \\ ax^2+1 & \text{si } x > 1 \end{cases}$, où $a$ est un réel.
Déterminer la valeur de $a$ pour que $h$ soit continue en $1$.
\textbf{3.} (Contexte technique) La vitesse $v(t)$ (en m/s) d'un véhicule est modélisée par :
\[
v(t) = \begin{cases} 2t & \text{si } 0 \le t < 5 \\ 10 & \text{si } 5 \le t \le 15 \\ -t+25 & \text{si } 15 < t \le 25 \end{cases}
\]
Étudier la continuité de $v$ en $t=5$ et en $t=15$. Interpréter physiquement les résultats.
\end{exercice}

\begin{corrige}[Exercice : continuité en un point]
\textbf{1.a)} $f(0)=2$. $\lim_{x\to 0^-}(x^2+1)=1$ et $\lim_{x\to 0^+}\sqrt{x+1}=\sqrt{1}=1$. Les limites latérales sont égales ($=1$) mais différentes de $f(0)=2$. \textbf{$f$ n'est pas continue en $0$ : discontinuité par trou.}
\textbf{1.b)} Le dénominateur s'annule en $2$, donc $\mathcal{D}_g = \mathbb{R}\setminus\{2\}$. $g$ n'est pas définie en $2$, donc \textbf{$g$ n'est pas continue en $2$}. Comme $\lim_{x\to 2}g(x)=\lim_{x\to 2}(x+2)=4$ (après simplification par $x-2$), cette discontinuité est un trou, qui pourra être « réparé » par prolongement par continuité (section 5).
\textbf{2.} $h(1)=3\times1-2=1$. Limite à gauche : $\lim_{x\to 1^-}(3x-2)=1$. Limite à droite : $\lim_{x\to 1^+}(ax^2+1)=a+1$. Continuité en $1 \iff a+1=1 \iff a=0$.
\textbf{3.} En $t=5$ : $v(5)=10$ ; $\lim_{t\to 5^-}2t=10$ ; $\lim_{t\to 5^+}10=10$. Les trois valeurs coïncident : \textbf{$v$ est continue en $5$}. Interprétation : le véhicule accélère progressivement jusqu'à 10 m/s, puis garde une vitesse constante, sans choc.
En $t=15$ : $v(15)=10$ ; $\lim_{t\to 15^-}10=10$ ; $\lim_{t\to 15^+}(-t+25)=-15+25=10$. \textbf{$v$ est continue en $15$}. Interprétation : fin de la phase à vitesse constante et début d'un freinage progressif ; le véhicule ne subit aucune variation brutale de vitesse, ce qui garantit le confort et la sécurité des passagers.
\end{corrige}

\coursec{Propriétés des fonctions continues}

Après avoir précisé ce que signifie la continuité en un point et sur un intervalle, nous allons voir que cette propriété est \emph{stable} par les opérations usuelles et qu'elle entraîne une conséquence majeure : le théorème des valeurs intermédiaires. Ce théorème est l'outil principal pour prouver l'existence de solutions à des équations sans nécessairement les calculer explicitement — une situation fréquente en technique (dimensionnement, seuils de déclenchement, équilibres thermiques).

\courssub{Stabilité de la continuité par les opérations algébriques}

L'intuition est simple : si deux grandeurs physiques varient sans « saut » (continuité), alors leur somme, leur produit ou leur quotient (tant qu'on ne divise pas par zéro) varient aussi sans saut.

\begin{propriete}[Opérations sur les fonctions continues — Admis]
Soient $f$ et $g$ deux fonctions définies sur un intervalle $I$ et continues sur $I$.
Alors les fonctions suivantes sont continues sur $I$ :
\begin{itemize}
    \item la somme $f+g$ ;
    \item le produit $f \times g$ ;
    \item le quotient $\dfrac{f}{g}$ \textbf{sur l'ensemble des $x \in I$ tels que $g(x) \neq 0$}.
\end{itemize}
\end{propriete}

\begin{attention}[Condition sur le dénominateur]
Pour le quotient $\frac{f}{g}$, la continuité n'est garantie qu'\emph{là où le dénominateur ne s'annule pas}. Si $g(x_0)=0$, le quotient n'est même pas défini en $x_0$ ; on ne peut pas parler de continuité en ce point.
\end{attention}

\begin{exemplebox}[Continuité de fonctions usuelles]
Toutes les fonctions de référence (polynômes, fonctions rationnelles sur leur ensemble de définition, racine carrée sur $[0;+\infty[$, fonctions trigonométriques) sont continues sur leurs intervalles de définition.
\par
Par composition et opérations, toute fonction obtenue par combinaison algébrique de ces fonctions de base est continue sur son domaine de définition.
\par
Exemple : $h(x) = \dfrac{x^2 + \sqrt{x}}{x-1}$ est continue sur $[0;1[ \cup ]1;+\infty[$.
\end{exemplebox}

\courssub{Théorème des valeurs intermédiaires (TVI)}

C'est le résultat le plus puissant de ce chapitre. Il formalise l'idée qu'une courbe continue ne peut pas « sauter » par-dessus une valeur intermédiaire.

\begin{propriete}[Théorème des valeurs intermédiaires — Admis au Probatoire]
Soit $f$ une fonction continue sur un intervalle \textbf{fermé} $[a ; b]$ (avec $a < b$).
Soit $k$ un nombre réel \textbf{compris entre $f(a)$ et $f(b)$} (c'est-à-dire $f(a) \le k \le f(b)$ ou $f(b) \le k \le f(a)$).
Alors il existe \textbf{au moins un} réel $c \in [a ; b]$ tel que $f(c) = k$.
\end{propriete}

\begin{aretenir}[Interprétation graphique]
Si l'on trace la courbe $\mathcal{C}_f$ d'une fonction continue sur $[a ; b]$, et que l'on trace une droite horizontale $y = k$ passant entre les points $A(a;f(a))$ et $B(b;f(b))$, alors cette droite \textbf{coupe obligatoirement la courbe} en au moins un point d'abscisse $c \in [a ; b]$.
\end{aretenir}

\begin{popfigure}
\begin{tikzpicture}
\begin{axis}[
    width=\linewidth, height=7cm,
    axis lines=middle,
    xlabel=$x$, ylabel=$y$,
    xmin=-0.5, xmax=4.5,
    ymin=-1, ymax=4,
    xtick={1,3}, xticklabels={$a$,$b$},
    ytick={1,3}, yticklabels={$f(a)$,$f(b)$},
    clip=false,
    domain=1:3, samples=200,
]
% Function: continuous curve from (1,1) to (3,3) with a bump
\addplot[popcurve, thick, popblue] { (x-1)*(x-3)*(x-2)/2 + (x-1) + 1 };
% Horizontal line y = k
\addplot[poporange, dashed, thick, domain=0.5:4] {2.2};
\addplot[poporange, dashed, thick, domain=0.5:4] {2.2} node[right, pos=0.95] {$y=k$};
% Points A and B
\addplot[only marks, mark=*, popblue, mark size=2.5] coordinates {(1,1) (3,3)};
% Intersection points (approximate visual)
\addplot[only marks, mark=*, poporange, mark size=2.5] coordinates {(1.4,2.2) (2.6,2.2)};
% Labels
\node[popblue, anchor=north east] at (axis cs:1,1) {$A$};
\node[popblue, anchor=south west] at (axis cs:3,3) {$B$};
\node[poporange, anchor=south west] at (axis cs:1.4,2.2) {$C_1$};
\node[poporange, anchor=south east] at (axis cs:2.6,2.2) {$C_2$};
% Vertical lines for c1, c2
\draw[popmuted, thin, dashed] (axis cs:1.4,0) -- (axis cs:1.4,2.2);
\draw[popmuted, thin, dashed] (axis cs:2.6,0) -- (axis cs:2.6,2.2);
\node[popmuted, anchor=north] at (axis cs:1.4,0) {$c_1$};
\node[popmuted, anchor=north] at (axis cs:2.6,0) {$c_2$};
\end{axis}
\end{tikzpicture}
\legende{Illustration du TVI : la droite $y=k$ coupe la courbe en au moins un point ($c_1$ et $c_2$ ici).}
\end{popfigure}

\begin{attention}[Hypothèses indispensables et non-unicité]
\begin{itemize}
    \item \textbf{Continuité sur $[a;b]$ :} Si $f$ n'est pas continue (ex: asymptote verticale, saut), le théorème est \textbf{faux}. La courbe peut « sauter » par-dessus $k$.
    \item \textbf{Intervalle fermé $[a;b]$ :} L'intervalle doit être borné et fermé.
    \item \textbf{Non-unicité :} Le TVI garantit l'\emph{existence} d'au moins un $c$, mais \textbf{pas l'unicité}. Il peut y avoir plusieurs solutions (voir figure ci-dessus avec $c_1$ et $c_2$).
\end{itemize}
\end{attention}

\courssub{Corollaire : Recherche de zéros (Changement de signe)}

C'est l'application directe la plus utilisée au Probatoire et en technique.

\begin{propriete}[Corollaire du TVI — Zéros d'une fonction continue]
Soit $f$ continue sur $[a ; b]$.
Si $f(a)$ et $f(b)$ sont de \textbf{signes contraires} (c'est-à-dire $f(a) \times f(b) < 0$), alors l'équation $f(x) = 0$ admet \textbf{au moins une solution} dans l'intervalle \textbf{ouvert} $]a ; b[$.
\end{propriete}

\begin{methode}[{Prouver l'existence d'une solution dans $]a;b[$}]
Pour montrer que l'équation $f(x)=0$ a une solution dans $]a;b[$ :
\begin{enumerate}
    \item Vérifier que $f$ est \textbf{continue} sur $[a;b]$ (citer la nature de $f$ : polynôme, rationnelle sans annulation du dénominateur, etc.).
    \item Calculer $f(a)$ et $f(b)$.
    \item Vérifier que $f(a) \times f(b) < 0$ (signes contraires).
    \item Conclure par le corollaire du TVI : « Il existe au moins un $c \in ]a;b[$ tel que $f(c)=0$. »
\end{enumerate}
\end{methode}

\begin{exoresolu}[Localisation d'une solution : $x^3 + x - 1 = 0$]
\textbf{Énoncé :} Montrer que l'équation $x^3 + x - 1 = 0$ admet une unique solution dans $]0 ; 1[$.
\par
\textbf{Solution :}
\begin{enumerate}
    \item On pose $f(x) = x^3 + x - 1$. C'est une fonction polynomiale, donc \textbf{continue sur $\mathbb{R}$}, et en particulier sur $[0 ; 1]$.
    \item On calcule les valeurs aux bornes :
    \[ f(0) = 0^3 + 0 - 1 = -1 < 0 \]
    \[ f(1) = 1^3 + 1 - 1 = 1 > 0 \]
    \item On a $f(0) \times f(1) = -1 < 0$. Les valeurs sont de signes contraires.
    \item \textbf{Par le corollaire du TVI}, l'équation $f(x)=0$ admet \textbf{au moins une solution} dans $]0 ; 1[$.
    \item \textbf{Unicité (complément) :} $f'(x) = 3x^2 + 1 > 0$ pour tout $x \in \mathbb{R}$. $f$ est strictement croissante sur $\mathbb{R}$, donc \textbf{injective}. L'équation $f(x)=0$ ne peut avoir \textbf{qu'une seule} solution sur $\mathbb{R}$, donc une unique dans $]0;1[$.
\end{enumerate}
\end{exoresolu}

\courssub{Application : Encadrement d'une solution par balayage (Table de valeurs)}

Le TVI nous dit \emph{qu'il existe} une solution, mais pas \emph{où} elle se trouve précisément. La méthode de \textbf{balayage} (ou recherche dichotomique manuelle) permet d'encadrer la solution dans un intervalle de plus en plus petit.

\begin{methode}[Encadrer une solution par balayage (pas de 0,1 ou 0,01)]
On cherche un intervalle $[a;b]$ tel que $f(a)f(b)<0$. On calcule ensuite les valeurs de $f$ pour des abscisses intermédiaires (par pas de 0,1, puis 0,01) jusqu'à trouver un changement de signe sur un sous-intervalle plus petit.
\end{methode}

\begin{exemplebox}[{Balayage pour $f(x)=x^3+x-1$ sur $]0;1[$}]
On construit une table de valeurs par pas de 0,1 :
\begin{center}
\begin{tabular}{|c|*{11}{c|}}\hline
$x$ & 0,0 & 0,1 & 0,2 & 0,3 & 0,4 & 0,5 & 0,6 & 0,7 & 0,8 & 0,9 & 1,0 \\ \hline
$f(x)$ & -1,000 & -0,899 & -0,792 & -0,673 & -0,536 & -0,375 & -0,184 & \textbf{0,043} & 0,312 & 0,629 & 1,000 \\ \hline
\end{tabular}
\end{center}
\par
On observe un changement de signe entre $x=0,6$ ($f<0$) et $x=0,7$ ($f>0$).
\par
La solution $\alpha$ est encadrée : $\alpha \in ]0,6 ; 0,7[$.
\par
On affine par pas de 0,01 entre 0,6 et 0,7 :
\begin{center}
\begin{tabular}{|c|*{11}{c|}}\hline
$x$ & 0,60 & 0,61 & 0,62 & 0,63 & 0,64 & 0,65 & 0,66 & 0,67 & 0,68 & 0,69 & 0,70 \\ \hline
$f(x)$ & -0,184 & -0,163 & -0,142 & -0,120 & -0,098 & -0,075 & -0,053 & -0,029 & \textbf{-0,006} & \textbf{0,019} & 0,043 \\ \hline
\end{tabular}
\end{center}
\par
Changement de signe entre $0,68$ et $0,69$.
\par
\textbf{Encadrement final :} $\alpha \in ]0,68 ; 0,69[$ (amplitude $0,01$).
\end{exemplebox}

\begin{popfigure}
\begin{tikzpicture}
\begin{axis}[
    width=\linewidth, height=6cm,
    axis lines=middle,
    xlabel=$x$, ylabel=$f(x)$,
    xmin=0.63, xmax=0.72,
    ymin=-0.1, ymax=0.05,
    xtick={0.64,0.65,0.66,0.67,0.68,0.69,0.70},
    xticklabels={0,64,0,65,0,66,0,67,0,68,0,69,0,70},
    ytick={-0.08,-0.04,0,0.04},
    yticklabels={-0,08,-0,04,0,0,04},
    domain=0.63:0.72, samples=200,
    tick label style={font=\footnotesize},
    x tick label style={rotate=45, anchor=east},
]
\addplot[popcurve, thick, popblue] { x^3 + x - 1 };
% Horizontal axis y=0
\addplot[popline, thick, domain=0.63:0.72] {0};
% Vertical lines for scan
\draw[popmuted, very thin, dashed] (axis cs:0.64,0) -- (axis cs:0.64,-0.09786);
\draw[popmuted, very thin, dashed] (axis cs:0.65,0) -- (axis cs:0.65,-0.07537);
\draw[popmuted, very thin, dashed] (axis cs:0.66,0) -- (axis cs:0.66,-0.05250);
\draw[popmuted, very thin, dashed] (axis cs:0.67,0) -- (axis cs:0.67,-0.02924);
\draw[popmuted, very thin, dashed] (axis cs:0.68,0) -- (axis cs:0.68,-0.00557);
\draw[popmuted, very thin, dashed] (axis cs:0.69,0) -- (axis cs:0.69,0.01851);
\draw[popmuted, very thin, dashed] (axis cs:0.70,0) -- (axis cs:0.70,0.04300);
% Highlight interval [0.68, 0.69]
\draw[popgreen, thick] (axis cs:0.68,-0.1) -- (axis cs:0.68,0.05);
\draw[popgreen, thick] (axis cs:0.69,-0.1) -- (axis cs:0.69,0.05);
\node[popgreen, anchor=south, font=\small] at (axis cs:0.685,0.04) {Solution $\in ]0,68;0,69[$};
% Points
\addplot[only marks, mark=*, popblue] coordinates {(0.68,-0.005568) (0.69,0.018509)};
\end{axis}
\end{tikzpicture}
\legende{Balayage graphique : la courbe coupe l'axe des abscisses entre 0,68 et 0,69.}
\end{popfigure}

\courssub{Prolongement par continuité}

Lorsque une fonction n'est pas définie en un point $x_0$ mais admet une limite finie en ce point, on peut « combler le trou » pour la rendre continue.

\begin{definition}[Prolongement par continuité]
Soit $f$ une fonction définie sur un intervalle $I$ sauf en $x_0 \in I$.
Si $\displaystyle\lim_{x \to x_0} f(x) = \ell \in \mathbb{R}$, on définit la fonction $\tilde{f}$ sur $I$ par :
\[
\tilde{f}(x) = \begin{cases}
f(x) & \text{si } x \neq x_0 \\
\ell   & \text{si } x = x_0
\end{cases}
\]
La fonction $\tilde{f}$ est continue sur $I$ (et en particulier en $x_0$). On dit que $\tilde{f}$ est le \textbf{prolongement par continuité} de $f$ en $x_0$.
\end{definition}

\begin{methode}[Prolonger une fonction par continuité en $x_0$]
\begin{enumerate}
    \item Vérifier que $f$ n'est pas définie en $x_0$ (ou que $f(x_0) \neq \lim_{x\to x_0}f(x)$).
    \item Calculer la limite $\ell = \lim_{x \to x_0} f(x)$.
    \item Si $\ell \in \mathbb{R}$, définir $\tilde{f}(x_0) = \ell$. La fonction prolongée est continue en $x_0$.
    \item Si la limite est infinie ou n'existe pas, \textbf{aucun prolongement par continuité n'est possible}.
\end{enumerate}
\end{methode}

\begin{exoresolu}[Prolongement par continuité : $f(x) = \dfrac{x^2-4}{x-2}$]
\textbf{Énoncé :} Étudier la possibilité de prolonger par continuité la fonction $f(x) = \dfrac{x^2-4}{x-2}$ en $x_0 = 2$.
\par
\textbf{Solution :}
\begin{enumerate}
    \item $f$ n'est pas définie en $2$ (dénominateur nul).
    \item Pour $x \neq 2$, on factorise : $f(x) = \dfrac{(x-2)(x+2)}{x-2} = x+2$.
    \item $\displaystyle\lim_{x \to 2} f(x) = \lim_{x \to 2} (x+2) = 4 \in \mathbb{R}$.
    \item On définit $\tilde{f}(2) = 4$. La fonction prolongée $\tilde{f}(x) = x+2$ (pour tout $x \in \mathbb{R}$) est continue sur $\mathbb{R}$.
\end{enumerate}
\end{exoresolu}

\begin{exoresolu}[Cas impossible : $g(x) = \dfrac{1}{x}$ en $0$]
\textbf{Énoncé :} Peut-on prolonger $g(x)=\frac{1}{x}$ par continuité en $0$ ?
\par
\textbf{Solution :} $\lim_{x \to 0^-} g(x) = -\infty$ et $\lim_{x \to 0^+} g(x) = +\infty$. La limite n'est pas finie. \textbf{Aucun prolongement par continuité n'est possible} en $0$.
\end{exoresolu}

\courssub{Applications concrètes : Seuils et franchissements}

En technique, le TVI justifie rigoureusement qu'un système continu franchit obligatoirement un seuil s'il passe d'un état « en-dessous » à un état « au-dessus ».

\begin{exoresolu}[Niveau d'un réservoir]
\textbf{Énoncé :} Un réservoir d'eau est rempli par une pompe. Le volume d'eau $V(t)$ (en $m^3$) en fonction du temps $t$ (en heures) est modélisé par la fonction continue $V(t) = 2t^2 + 3t$ sur l'intervalle $[0 ; 2]$. Le réservoir a une capacité de $10~m^3$. Un capteur déclenche une alarme quand le volume atteint exactement $8~m^3$. Prouver que l'alarme se déclenchera nécessairement entre $t=1$ h et $t=2$ h.
\par
\textbf{Solution :}
\begin{enumerate}
    \item $V$ est polynomiale, donc \textbf{continue sur $[0;2]$}.
    \item On calcule : $V(1) = 2(1)^2 + 3(1) = 5~m^3$.
    \item $V(2) = 2(2)^2 + 3(2) = 8 + 6 = 14~m^3$.
    \item La valeur seuil $k = 8$ est comprise entre $V(1)=5$ et $V(2)=14$ ($5 < 8 < 14$).
    \item \textbf{Par le TVI}, il existe au moins un instant $c \in [1 ; 2]$ tel que $V(c) = 8$.
    \item \textbf{Conclusion :} L'alarme se déclenchera bien au cours de la 2\textsuperscript{e} heure.
\end{enumerate}
\end{exoresolu}

\begin{exoresolu}[Température d'un four industriel]
\textbf{Énoncé :} La température $\theta(t)$ (en $^\circ$C) d'un four en phase de chauffe est donnée par $\theta(t) = 20 + 200(1 - e^{-0,2t})$ pour $t \in [0 ; 30]$ min. On veut savoir si la température a atteint $150^\circ$C durant les 10 premières minutes.
\par
\textbf{Solution :}
\begin{enumerate}
    \item $\theta$ est continue sur $[0;30]$ (somme, produit, composée de fonctions continues).
    \item $\theta(0) = 20 + 200(1-1) = 20^\circ$C.
    \item $\theta(10) = 20 + 200(1 - e^{-2}) \approx 20 + 200(1 - 0,135) \approx 20 + 173 = 193^\circ$C.
    \item $150$ est entre $20$ et $193$.
    \item \textbf{Par le TVI}, $\exists c \in [0;10]$ tel que $\theta(c) = 150$. Le seuil a été atteint.
\end{enumerate}
\end{exoresolu}

\begin{savaistu}[Théorème de Bolzano — Histoire et réalité]
Le théorème des valeurs intermédiaires est aussi appelé \textbf{théorème de Bolzano} (Bernard Bolzano, 1781–1848, mathématicien tchèque). Il l'a démontré rigoureusement en 1817, bien avant Cauchy et Weierstrass.
\par
C'est ce théorème qui garantit qu'à un moment donné de la journée, la température extérieure a été \emph{exactement} $30,000\dots^\circ$C (si le matin il faisait $25^\circ$C et l'après-midi $35^\circ$C), ou qu'il existe deux points antipodaux sur Terre qui ont \emph{exactement} la même température et la même pression atmosphérique en ce moment même (théorème de Borsuk-Ulam, généralisation topologique du TVI). En technique, c'est la base mathématique de tout système de régulation par seuils (thermostats, pressostats, fins de course).
\end{savaistu}

\begin{attention}[Pièges classiques au Probatoire]
\begin{itemize}
    \item \textbf{Oublier de vérifier la continuité :} Écrire « Par le TVI... » sans justifier que $f$ est continue sur $[a;b]$ fait perdre la majorité des points.
    \item \textbf{Confondre intervalle ouvert et fermé :} L'hypothèse porte sur $[a;b]$ (fermé), la conclusion donne $c \in [a;b]$. Le corollaire des zéros donne $c \in ]a;b[$ (ouvert) car $f(a)$ et $f(b)$ sont non nuls (signes contraires).
    \item \textbf{Dire « l'unique solution » :} Le TVI ne donne \textbf{jamais} l'unicité. Pour l'unicité, il faut ajouter la monotonie (dérivée de signe constant).
    \item \textbf{Appliquer le TVI sur un intervalle non fermé ou non borné :} Ex: $f(x)=1/x$ sur $]0;1]$. $f(1)=1$, $\lim_{0^+}f=+\infty$. La valeur $k=2$ est « entre » les deux, mais il n'y a pas de $c$ dans $]0;1]$ tel que $f(c)=2$ ? Si, $c=0,5$. Mais si on prend $k=0$, $0$ n'est pas atteint. L'intervalle doit être \textbf{fermé} $[a;b]$.
    \item \textbf{Prolongement par continuité :} Ne pas confondre « limite finie » (prolongement possible) et « limite infinie » (prolongement impossible). Vérifier toujours que la limite est un nombre réel.
\end{itemize}
\end{attention}

\coursec{Prolongement par continuité}

Dans la section précédente, nous avons étudié les propriétés des fonctions continues : une fonction continue sur un intervalle fermé y possède des bornes, atteint ses extremums et ne peut pas passer d'une valeur positive à une valeur négative sans s'annuler (théorème des valeurs intermédiaires). Mais toutes ces belles propriétés supposent que la fonction est \emph{déjà} continue. Que faire lorsqu'une fonction présente un \emph{petit défaut} : elle n'est pas définie en un seul point $a$, alors que tout se passe bien autour ? Peut-on \og réparer \fg{} la fonction en lui donnant une valeur bien choisie en $a$ ? C'est exactement l'objet du \cle{prolongement par continuité}, une notion très concrète : il s'agit de \og reboucher le trou \fg{} d'une courbe.

\courssub{Situation : une fonction avec un \og trou \fg{} en un point}

\courssub{Intuition : le trou dans la courbe}

Considérons la fonction définie par
\[ f(x) = \dfrac{x^2 - 1}{x - 1}. \]
Cette fonction n'est pas définie en $x = 1$, car le dénominateur s'y annule : on ne peut pas diviser par zéro. Pourtant, pour tout $x \neq 1$, on peut simplifier :
\[ f(x) = \dfrac{(x-1)(x+1)}{x-1} = x + 1 \qquad (x \neq 1). \]
La courbe de $f$ est donc la droite d'équation $y = x + 1$, \emph{privée du point d'abscisse $1$}. Il y a un \og trou \fg{} en ce point. Mais lorsque $x$ se rapproche de $1$, $f(x)$ se rapproche de $2$ : la fonction admet une limite finie en $1$ :
\[ \lim_{x \to 1} f(x) = \lim_{x \to 1} (x+1) = 2. \]
L'idée du prolongement par continuité est alors toute simple : puisque la courbe \og voudrait \fg{} passer par le point $(1\,;\,2)$, donnons à la fonction la valeur $2$ en $1$. Le trou est rebouché, et la fonction obtenue est continue partout, y compris en $1$.

\begin{popfigure}
\begin{center}
\begin{tikzpicture}
\begin{axis}[
  width=0.85\linewidth, height=6.5cm,
  axis lines=middle, xlabel=$x$, ylabel=$y$,
  xmin=-1.5, xmax=3.5, ymin=-1, ymax=5,
  xtick={-1,1,2,3}, ytick={1,2,3,4},
  grid=both, grid style={popline!40},
  legend style={at={(0.02,0.98)}, anchor=north west, font=\small}
]
\addplot[popcurve,domain=-1.2:3.2,samples=100]{x+1};
\addlegendentry{$y = f(x)$, trou en $(1\,;\,2)$}
\addplot[only marks, mark=o, mark size=2.5pt, color=poporange, thick] coordinates {(1,2)};
\addplot[only marks, mark=*, mark size=2pt, color=popgreen] coordinates {(1,2)};
\node[color=popgreen, font=\small, anchor=west] at (axis cs:1.1,1.7) {$g(1)=2$ : le trou est rebouché};
\end{axis}
\end{tikzpicture}
\end{center}
\legende{La courbe de $f(x)=\dfrac{x^2-1}{x-1}$ est la droite $y=x+1$ avec un trou (point creux) en $(1\,;\,2)$. Le point plein vert représente la valeur $g(1)=2$ qui prolonge $f$ par continuité.}
\end{popfigure}

\courssub{Définition du prolongement par continuité}

\begin{definition}[Prolongement par continuité]
Soit $f$ une fonction définie sur un intervalle $I$, \textbf{sauf en un point} $a$ de $I$. On suppose que $f$ admet en $a$ une \cle{limite finie} $\ell$ :
\[ \lim_{x \to a} f(x) = \ell \quad (\ell \in \mathbb{R}). \]
La fonction $g$ définie sur $I$ tout entier par
\[ g(x) = \begin{cases} f(x) & \text{si } x \neq a \\ \ell & \text{si } x = a \end{cases} \]
est appelée le \cle{prolongement par continuité} de $f$ en $a$. La fonction $g$ est alors \textbf{continue en} $a$ (et donc continue sur $I$ si $f$ l'était en dehors de $a$).
\end{definition}

\begin{aretenir}[Condition d'existence]
Le prolongement par continuité de $f$ en $a$ existe \textbf{si et seulement si} $f$ admet en $a$ une limite \emph{finie}. La valeur donnée à $g$ en $a$ est alors nécessairement cette limite : c'est le seul choix qui rende $g$ continue en $a$.
\end{aretenir}

\textbf{Pourquoi $g$ est-elle continue en $a$ ?} C'est immédiat à partir de la définition de la continuité vue dans la section 3 : par construction, $g(a) = \ell$ et $\displaystyle\lim_{x \to a} g(x) = \lim_{x \to a} f(x) = \ell$ (puisque $g$ et $f$ coïncident partout sauf en $a$, et que la limite en $a$ ne dépend pas de la valeur en $a$). On a donc bien $\displaystyle\lim_{x \to a} g(x) = g(a)$, ce qui est exactement la condition de continuité de $g$ en $a$.

\begin{methode}[Prolonger une fonction par continuité en un point]
Pour déterminer si une fonction $f$, non définie en $a$, est prolongeable par continuité en $a$ :
\begin{enumerate}
\item Vérifier que $f$ n'est pas définie en $a$ (en général, le dénominateur s'annule en $a$).
\item Calculer $\displaystyle\lim_{x \to a} f(x)$. Si l'on obtient une forme indéterminée du type $\dfrac{0}{0}$, \textbf{lever l'indétermination} en factorisant le numérateur et le dénominateur par $(x-a)$, puis en simplifiant.
\item Si la limite est un nombre fini $\ell$ : $f$ est prolongeable par continuité en $a$, et l'on pose $g(a) = \ell$.
\item Si la limite est infinie, ou si les limites à gauche et à droite sont différentes : le prolongement est \textbf{impossible}.
\end{enumerate}
\end{methode}

\courssub{Exemples fondateurs}

\begin{exoresolu}[Prolongement de $f(x) = \dfrac{x^2 - 9}{x - 3}$ en $3$]
Soit $f$ la fonction définie par $f(x) = \dfrac{x^2 - 9}{x - 3}$.
\begin{enumerate}
\item Quel est l'ensemble de définition de $f$ ?
\item Montrer que $f$ est prolongeable par continuité en $3$ et définir ce prolongement $g$.
\end{enumerate}
\tcblower
\textbf{Solution détaillée.}
\begin{enumerate}
\item $f$ est définie pour tout réel $x$ tel que $x - 3 \neq 0$, c'est-à-dire sur $\mathbb{R} \setminus \{3\}$. La fonction n'est pas définie en $3$.
\item Calculons la limite de $f$ en $3$. En $x = 3$, le numérateur vaut $9 - 9 = 0$ et le dénominateur vaut $0$ : c'est une forme indéterminée $\dfrac{0}{0}$. On lève l'indétermination en factorisant le numérateur (identité remarquable $a^2 - b^2 = (a-b)(a+b)$) :
\[ f(x) = \dfrac{x^2 - 9}{x - 3} = \dfrac{(x-3)(x+3)}{x-3} = x + 3 \qquad (x \neq 3). \]
Donc :
\[ \lim_{x \to 3} f(x) = \lim_{x \to 3} (x+3) = 6. \]
La limite est finie : $f$ est prolongeable par continuité en $3$. Le prolongement est la fonction $g$ définie sur $\mathbb{R}$ par :
\[ g(x) = \begin{cases} \dfrac{x^2 - 9}{x - 3} & \text{si } x \neq 3 \\ 6 & \text{si } x = 3 \end{cases} \]
c'est-à-dire simplement $g(x) = x + 3$ pour tout réel $x$. La fonction $g$ est continue sur $\mathbb{R}$ tout entier.
\end{enumerate}
\end{exoresolu}

\begin{exemplebox}[Un exemple issu de l'atelier]
Un technicien mesure le débit moyen d'une pompe : $q(t) = \dfrac{V(t)}{t}$ litres par minute, où $V(t) = 2t + t^2$ est le volume écoulé après $t$ minutes. En $t = 0$, le débit moyen n'est pas défini (division par zéro). Mais pour $t \neq 0$ :
\[ q(t) = \dfrac{2t + t^2}{t} = 2 + t \quad \xrightarrow[t \to 0]{} \quad 2. \]
On prolonge $q$ par continuité en $0$ par $q(0) = 2$ : c'est le \emph{débit instantané} de la pompe au démarrage, soit $2$~L/min. Le prolongement par continuité donne ainsi un sens physique à une quantité qui n'était pas définie.
\end{exemplebox}

\courssub{Contre-exemple : quand le prolongement est impossible}

\begin{attention}[Le prolongement n'existe pas toujours]
Le prolongement par continuité exige une limite \textbf{finie} en $a$. Deux situations rendent le prolongement impossible :
\begin{itemize}
\item la limite en $a$ est \textbf{infinie} (par exemple $x \mapsto \dfrac{1}{x}$ en $0$) : la courbe s'envole vers $+\infty$ ou $-\infty$, il n'y a aucun point à placer pour reboucher le trou ;
\item les \textbf{limites à gauche et à droite} en $a$ sont finies mais \textbf{différentes} (par exemple $x \mapsto \dfrac{|x|}{x}$ en $0$, qui vaut $-1$ à gauche et $+1$ à droite) : la courbe fait un \og saut \fg{} qu'aucune valeur unique ne peut réparer.
\end{itemize}
\end{attention}

Étudions le cas de la fonction inverse. La fonction $h(x) = \dfrac{1}{x}$ n'est pas définie en $0$, et :
\[ \lim_{\substack{x \to 0 \\ x > 0}} \dfrac{1}{x} = +\infty \qquad \text{et} \qquad \lim_{\substack{x \to 0 \\ x < 0}} \dfrac{1}{x} = -\infty. \]
Aucune valeur finie $h(0)$ ne peut rendre $h$ continue en $0$ : la courbe \og part \fg{} vers $+\infty$ d'un côté et vers $-\infty$ de l'autre. Le trou est infranchissable.

\begin{popfigure}
\begin{center}
\begin{tikzpicture}
\begin{axis}[
  width=0.85\linewidth, height=7cm,
  axis lines=middle, xlabel=$x$, ylabel=$y$,
  xmin=-4, xmax=4, ymin=-6, ymax=6,
  xtick={-3,-2,-1,1,2,3}, ytick={-4,-2,2,4},
  grid=both, grid style={popline!40},
  legend style={at={(0.02,0.02)}, anchor=south west, font=\small}
]
\addplot[popcurve,domain=0.17:3.8,samples=100]{1/x};
\addplot[popcurve,domain=-3.8:-0.17,samples=100]{1/x};
\addlegendentry{$y = \dfrac{1}{x}$}
\draw[dashed, poporange, thick] (axis cs:0,-6) -- (axis cs:0,6);
\node[color=poporange, font=\small, anchor=west] at (axis cs:0.15,-5) {asymptote $x=0$ : pas de prolongement};
\end{axis}
\end{tikzpicture}
\end{center}
\legende{La courbe de $x \mapsto \dfrac{1}{x}$ s'envole vers $+\infty$ à droite de $0$ et vers $-\infty$ à gauche : impossible de \og reboucher le trou \fg{} en $0$.}
\end{popfigure}

\begin{exoresolu}[Le prolongement est-il possible ?]
Pour chacune des fonctions suivantes, dire si elle admet un prolongement par continuité en $0$ et, le cas échéant, le définir.
\begin{enumerate}
\item $f(x) = \dfrac{x^2 + 2x}{x}$ ;
\item $g(x) = \dfrac{1}{x^2}$ ;
\item $h(x) = \dfrac{|x|}{x}$.
\end{enumerate}
\tcblower
\textbf{Solution détaillée.}
\begin{enumerate}
\item Pour $x \neq 0$ : $f(x) = \dfrac{x(x+2)}{x} = x + 2$, donc $\displaystyle\lim_{x \to 0} f(x) = 2$, limite finie. Le prolongement existe : $\tilde{f}(x) = x + 2$ pour tout $x$, avec $\tilde{f}(0) = 2$.
\item $\displaystyle\lim_{x \to 0} \dfrac{1}{x^2} = +\infty$ : la limite est infinie, donc $g$ \textbf{n'admet pas} de prolongement par continuité en $0$.
\item Pour $x > 0$, $h(x) = \dfrac{x}{x} = 1$ ; pour $x < 0$, $h(x) = \dfrac{-x}{x} = -1$. Les limites à gauche ($-1$) et à droite ($+1$) sont différentes : $h$ n'a pas de limite en $0$, donc \textbf{pas de prolongement} par continuité.
\end{enumerate}
\end{exoresolu}

\courssub{Interprétation graphique : reboucher le trou}

\begin{aretenir}[Lecture graphique]
Graphiquement, prolonger $f$ par continuité en $a$ revient à ajouter à la courbe le \textbf{point manquant} $(a\,;\,\ell)$, là où la courbe \og arrive naturellement \fg{}. On dit que l'on \og rebouche le trou \fg{}. Après prolongement, la courbe peut être tracée sans lever le crayon, y compris au voisinage de $a$ : c'est exactement l'image de la continuité.
\end{aretenir}

\begin{attention}[Pièges fréquents]
\begin{itemize}
\item Ne pas confondre \og $f$ n'est pas définie en $a$ \fg{} avec \og $f$ n'a pas de limite en $a$ \fg{} : une fonction non définie en $a$ peut très bien avoir une limite finie en $a$ (c'est justement le cas favorable).
\item La valeur du prolongement est la \textbf{limite} de $f$ en $a$, jamais une valeur choisie au hasard : poser $g(3) = 5$ dans l'exemple résolu ci-dessus ne rendrait pas $g$ continue.
\item Penser à \textbf{simplifier avant} de conclure : écrire directement $\dfrac{0}{0}$ et déclarer le prolongement impossible est une erreur ; il faut lever l'indétermination par factorisation.
\item Ne pas oublier de préciser l'ensemble de définition du prolongement : $g$ est définie en $a$, contrairement à $f$.
\end{itemize}
\end{attention}

\begin{savaistu}[Vers la notion de dérivée]
Le prolongement par continuité est loin d'être une simple curiosité : il est au cœur de la notion de \textbf{dérivée}, que tu rencontreras dans la suite de ton parcours technique. Pour une fonction $f$ définie en $a$, le \emph{taux d'accroissement}
\[ \tau(x) = \dfrac{f(x) - f(a)}{x - a} \]
n'est jamais défini en $x = a$ (le dénominateur s'y annule). Dire que ce taux admet une limite finie en $a$, c'est exactement dire qu'il admet un prolongement par continuité en $a$ ; la valeur du prolongement est alors le nombre dérivé $f'(a)$, c'est-à-dire la pente de la tangente à la courbe au point d'abscisse $a$. C'est aussi ce prolongement que le technicien calcule lorsqu'il détermine une vitesse instantanée à partir de vitesses moyennes, ou un débit instantané à partir de débits moyens, comme dans l'exemple de la pompe.
\end{savaistu}

\begin{exercice}[À toi de jouer]
\begin{enumerate}
\item Soit $f(x) = \dfrac{x^2 - 5x + 6}{x - 2}$. Montrer que $f$ admet un prolongement par continuité en $2$ et le définir. (Indication : factoriser le numérateur en cherchant ses racines.)
\item Soit $g(x) = \dfrac{x + 1}{x^2}$. La fonction $g$ admet-elle un prolongement par continuité en $0$ ? Justifier.
\item Un four industriel chauffe selon la loi $T(t) = \dfrac{3t^2 + 12t}{t}$ (en degrés, $t$ en heures, $t > 0$). Quelle température initiale $T(0)$ faut-il afficher pour que la fonction $T$ soit continue en $0$ ?
\end{enumerate}
\end{exercice}

\begin{corrige}[Exercice : À toi de jouer]
\begin{enumerate}
\item Le numérateur s'annule en $2$ (car $4 - 10 + 6 = 0$), donc il se factorise par $(x-2)$ : $x^2 - 5x + 6 = (x-2)(x-3)$. Pour $x \neq 2$ : $f(x) = x - 3$, d'où $\displaystyle\lim_{x \to 2} f(x) = -1$. Le prolongement est $\tilde{f}(x) = x - 3$ avec $\tilde{f}(2) = -1$.
\item $\displaystyle\lim_{x \to 0} \dfrac{x+1}{x^2} = +\infty$ (le numérateur tend vers $1$, le dénominateur vers $0$ en restant positif). Limite infinie : pas de prolongement par continuité en $0$.
\item Pour $t > 0$ : $T(t) = \dfrac{t(3t + 12)}{t} = 3t + 12 \xrightarrow[t \to 0]{} 12$. Il faut afficher $T(0) = 12$ degrés : c'est le prolongement par continuité de $T$ en $0$.
\end{enumerate}
\end{corrige}

\coursec{Synthèse et applications à la vie courante}

\courssub{Bilan unifié : limite, continuité, TVI et prolongement}

Après avoir étudié séparément les limites, la continuité, le théorème des valeurs intermédiaires et le prolongement par continuité, il est temps de \textbf{relier ces notions} entre elles. Elles forment une chaîne logique qui permet de passer d'une information locale (comportement près d'un point) à une information globale (existence d'une valeur, régularité d'un phénomène).

\begin{popfigure}
\begin{center}
\begin{tikzpicture}[node distance=1.2cm and 1.6cm, every node/.style={font=\small, align=center}]
\node[draw, rounded corners, fill=popblueL, thick, align=center] (limite){Limite en un point\\$\lim\limits_{x\to x_0}f(x)=\ell$};
\node[draw, rounded corners, fill=popgreenL, thick, right=of limite, align=center] (cv){Continuité en $x_0$\\$\lim\limits_{x\to x_0}f(x)=f(x_0)$};
\node[draw, rounded corners, fill=poporangeL, thick, below=of cv, align=center] (tvi){Théorème des valeurs\\intermédiaires (TVI)\\$f$ continue sur $[a,b]$, $k$ entre $f(a)$ et $f(b)$\\$\Rightarrow\ \exists c\in[a,b],\ f(c)=k$};
\node[draw, rounded corners, fill=poppinkL, thick, left=of tvi, align=center] (prol){Prolongement par continuité\\$f$ admet une limite finie $\ell$ en $x_0\notin D_f$\\$\Rightarrow\ \tilde f$ continue sur $D_f\cup\{x_0\}$};

\draw[->, thick, popblue] (limite) -- node[above, font=\tiny] {si $\ell=f(x_0)$} (cv);
\draw[->, thick, popgreen] (cv) -- node[right, font=\tiny] {hypothèse} (tvi);
\draw[->, thick, poppurple] (limite) -- node[left, font=\tiny] {base de} (prol);
\draw[->, thick, poppink] (prol) |- node[left, font=\tiny, pos=0.25] {régularise} (cv);
\end{tikzpicture}
\end{center}
\legende{Relations entre les notions clés du chapitre : la limite est l'outil de base, la continuité en découle, le TVI utilise la continuité, le prolongement répare un « trou ».}
\end{popfigure}

\begin{aretenir}[Synthèse des liens]
\begin{itemize}
\item La \cle{limite} est l'outil de base : elle décrit le comportement de $f$ \emph{près} de $x_0$ sans exiger que $f$ soit définie en $x_0$.
\item La \cle{continuité} en $x_0$ signifie que la limite existe, est finie, et \emph{coïncide} avec la valeur de la fonction : pas de « trou » ni de « saut ».
\item Le \cle{TVI} est une propriété \emph{globale} des fonctions continues sur un intervalle : l'image d'un intervalle contient toutes les valeurs intermédiaires. Il transforme une information sur les bornes en une information sur \emph{toutes} les valeurs comprises entre elles.
\item Le \cle{prolongement par continuité} répare une discontinuité \emph{évitable} (un « trou ») en définissant la valeur manquante comme la limite. Il ne fonctionne \emph{pas} quand les limites à gauche et à droite sont différentes (saut) ni quand la limite est infinie (asymptote verticale).
\end{itemize}
\end{aretenir}

\courssub{Modélisation 1 : fonctions discontinues — tarification par paliers}

De nombreux services techniques (électricité, transport, télécommunications) utilisent une \cle{tarification par paliers} : le prix change brutalement quand on franchit un seuil. Cela produit une fonction \textbf{discontinue aux seuils}.

\begin{exemplebox}[Tarification de l'électricité (modèle simplifié)]
Chez le distributeur d'électricité, le tarif au kWh dépend de la consommation totale du mois : si l'on dépasse le seuil de la tranche sociale, \emph{toute} la consommation est facturée au tarif supérieur. Modélisons le coût $C(x)$ en francs CFA pour une consommation de $x$ kWh :
\[
C(x) = \begin{cases}
50x & \text{si } 0 \le x \le 110 \quad\text{(tranche sociale)}\\
79x & \text{si } 110 < x \le 400\\
94x & \text{si } x > 400
\end{cases}
\]
\begin{itemize}
\item $C$ est \textbf{continue} sur chacun des intervalles $[0,110]$, $]110,400]$, $]400,+\infty[$ (fonctions affines).
\item En $x=110$ : $\lim\limits_{x\to 110^-}C(x)=50\times110=5500$ et $\lim\limits_{x\to 110^+}C(x)=79\times110=8690$. Les limites à gauche et à droite sont \textbf{différentes} : $C$ est \textbf{discontinue} en $110$ (saut de $3190$ F).
\item En $x=400$ : $\lim\limits_{x\to 400^-}C(x)=79\times400=31600$ et $\lim\limits_{x\to 400^+}C(x)=94\times400=37600$ : $C$ est \textbf{discontinue} en $400$ (saut de $6000$ F).
\item Interprétation : consommer $110$ kWh coûte $5500$ F, mais consommer $111$ kWh coûte $79\times111=8769$ F : un seul kWh supplémentaire augmente la facture de plus de $3000$ F ! C'est l'effet de seuil.
\end{itemize}
\end{exemplebox}

\begin{popfigure}
\begin{center}
\begin{tikzpicture}
\begin{axis}[
    width=\linewidth, height=7.5cm,
    axis lines=middle,
    xlabel={Consommation $x$ (kWh)}, ylabel={Coût $C(x)$ (FCFA)},
    xmin=0, xmax=600, ymin=0, ymax=60000,
    xtick={0,110,400,600},
    ytick={5500,8690,31600,37600},
    yticklabels={$5\,500$,$8\,690$,$31\,600$,$37\,600$},
    grid=major, grid style={popline},
    clip=false
]
% Tranche 1 : 0 à 110
\addplot[popcurve, domain=0:110, samples=2, very thick] {50*x};
\addplot[mark=*, fill=popblue, draw=popblue] coordinates {(0,0) (110,5500)};
% Tranche 2 : 110 à 400
\addplot[popcurve, domain=110:400, samples=2, very thick, poporange] {79*x};
\addplot[mark=o, fill=white, draw=poporange, thick] coordinates {(110,8690)};
\addplot[mark=*, fill=poporange, draw=poporange] coordinates {(400,31600)};
% Tranche 3 : 400 à 600
\addplot[popcurve, domain=400:600, samples=2, very thick, popgreen] {94*x};
\addplot[mark=o, fill=white, draw=popgreen, thick] coordinates {(400,37600)};
% Lignes pointillées de rappel
\addplot[dashed, popmuted, domain=0:110] coordinates {(110,0) (110,5500)};
\addplot[dashed, popmuted] coordinates {(400,0) (400,31600)};
\end{axis}
\end{tikzpicture}
\end{center}
\legende{Coût de l'électricité en fonction de la consommation : fonction affine par morceaux, discontinue aux seuils $x=110$ et $x=400$ (points pleins = valeur prise, points creux = valeur non prise).}
\end{popfigure}

\begin{attention}[Piège : continuité sur chaque morceau ne suffit pas]
Une fonction « continue par morceaux » n'est \emph{pas} nécessairement continue sur tout son domaine. Il faut vérifier \emph{systématiquement} aux points de jonction :
\[
\lim_{x\to x_0^-}f(x) \stackrel{?}{=} \lim_{x\to x_0^+}f(x) \stackrel{?}{=} f(x_0)
\]
Si l'une de ces égalités échoue, la fonction est discontinue en $x_0$. Dans un problème concret, \textbf{précisez toujours l'ensemble de définition} du modèle (par exemple $x\ge 0$ pour une consommation).
\end{attention}

\courssub{Modélisation 2 : fonctions continues — grandeurs physiques sans saut}

À l'inverse, de nombreuses grandeurs physiques évoluent \textbf{sans discontinuité} : température, niveau d'eau, pression, position d'un mobile. La continuité n'est pas qu'une propriété mathématique, c'est un \textbf{principe physique} : une grandeur naturelle ne « saute » pas d'une valeur à une autre sans passer par les valeurs intermédiaires.

\begin{exemplebox}[Niveau d'un réservoir — application du TVI]
Un réservoir d'eau contient $200$ L à $8$ h et $800$ L à $14$ h. Le remplissage se fait par un robinet à débit régulier, sans vidange.
\begin{enumerate}
\item Modéliser le volume $V(t)$ en litres en fonction du temps $t$ en heures ($t=0$ correspond à $8$ h).
\item Montrer qu'il y a eu un instant où le réservoir contenait exactement $500$ L.
\end{enumerate}
\textbf{Solution.}
\begin{itemize}
\item $V$ est définie sur $[0,6]$, avec $V(0)=200$ et $V(6)=800$. Physiquement, $V$ est \textbf{continue} : le volume d'eau ne peut pas passer instantanément de $200$ L à $800$ L.
\item $500\in[200,800]$. Par le \textbf{TVI} appliqué à $V$ continue sur $[0,6]$, il existe $c\in[0,6]$ tel que $V(c)=500$.
\item Concrètement : entre $8$ h et $14$ h, le réservoir a \emph{nécessairement} contenu exactement $500$ L à un certain instant.
\end{itemize}
\end{exemplebox}

\begin{savaistu}[Ingénierie et continuité]
Les ingénieurs civils (ponts, barrages, réseaux d'eau) utilisent la continuité comme \textbf{critère de sûreté}. Une contrainte qui varie brutalement (discontinuité) crée une concentration de contraintes, donc un risque de rupture. Les logiciels de calcul vérifient que les déplacements et les contraintes sont des fonctions \emph{continues} dans la structure. Un pont qui « saute » n'existe pas !
\end{savaistu}

\courssub{Rédiger une démarche complète : la méthode en quatre étapes}

Au Probatoire, une réponse « oui/non » ou un calcul brut ne suffit pas. Il faut une \textbf{argumentation structurée} qui montre que vous maîtrisez les hypothèses des théorèmes.

\begin{methode}[Rédaction d'une justification complète — 4 étapes]
\begin{enumerate}
\item \textbf{Hypothèses / Cadre} : préciser l'ensemble de définition $D_f$, l'intervalle considéré, la continuité (ou non) de $f$ sur cet intervalle. Citer la propriété qui justifie la continuité (fonction polynôme, grandeur physique, etc.).
\item \textbf{Théorème invoqué} : nommer explicitement le théorème utilisé (TVI, prolongement par continuité, etc.) et rappeler ses conditions d'application.
\item \textbf{Calculs / Vérifications} : effectuer les calculs de limites, d'images, de bornes. Montrer que les conditions numériques sont satisfaites (par exemple $k$ compris entre $f(a)$ et $f(b)$).
\item \textbf{Conclusion rédigée} : phrase complète répondant à la question posée, en langage courant \emph{et} mathématique.
\end{enumerate}
\end{methode}

\begin{exoresolu}[Recherche d'un instant précis — TVI]
\textbf{Énoncé.} La température $T(t)$ (en °C) d'un four en chauffe est modélisée par $T(t)=20+15t-t^2$ sur $[0,10]$ ($t$ en minutes). Montrer qu'il existe un instant où la température est exactement $50$ °C, puis déterminer cet instant.
\tcblower
\textbf{Solution rédigée en 4 étapes.}
\begin{enumerate}
\item \textbf{Hypothèses.} $T$ est une fonction polynôme, donc \textbf{continue sur $\mathbb{R}$}, et en particulier sur $[0,10]$. On calcule les bornes : $T(0)=20$ et $T(10)=20+150-100=70$.
\item \textbf{Théorème.} Théorème des valeurs intermédiaires : si $f$ est continue sur $[a,b]$, alors pour tout réel $k$ compris entre $f(a)$ et $f(b)$, il existe au moins un réel $c\in[a,b]$ tel que $f(c)=k$.
\item \textbf{Calculs.} On a bien $50\in[20,70]$, donc le TVI s'applique : il existe $c\in[0,10]$ tel que $T(c)=50$. Pour trouver $c$, on résout :
\begin{align*}
20+15c-c^2=50 &\iff c^2-15c+30=0\\
\Delta &= 15^2-4\times 30 = 225-120=105\\
c &= \frac{15\pm\sqrt{105}}{2}
\end{align*}
Seule la racine $c=\dfrac{15-\sqrt{105}}{2}\approx 2,38$ appartient à $[0,10]$ (l'autre vaut environ $12,6$).
\item \textbf{Conclusion.} Il existe un instant, environ $2,4$ minutes (soit $2$ min $23$ s) après le démarrage, où la température du four est exactement $50$ °C.
\end{enumerate}
\textbf{Analyse de la rédaction :} chaque étape est identifiée, le TVI est nommé, les bornes sont calculées, l'appartenance de $50$ à l'intervalle $[20,70]$ est vérifiée, la conclusion est une phrase complète.
\end{exoresolu}

\courssub{Critères de choix : quel outil pour quelle question ?}

Devant un problème concret, il faut \textbf{identifier la question} pour choisir l'outil adapté. Voici un organigramme de décision.

\begin{popfigure}
\begin{center}
\begin{tikzpicture}[
    decision/.style={diamond, draw, thick, fill=poporangeL, aspect=2.2, inner sep=1pt, font=\small, align=center},
    action/.style={rectangle, rounded corners, draw, thick, fill=popblueL, inner sep=4pt, font=\small, align=center},
    start/.style={ellipse, draw, thick, fill=popgreenL, inner sep=4pt, font=\small, align=center},
    fleche/.style={->, thick, popdark},
    node distance=1.4cm and 2.2cm
]
\node[start, align=center] (start){Question posée\\sur une fonction $f$};
\node[decision, below=of start, align=center] (q1){Comportement\\près d'un point\\ou à l'infini ?};
\node[action, left=of q1, align=center] (lim){Outil : \textbf{LIMITE}\\calculer $\lim\limits_{x\to x_0}f(x)$\\ou $\lim\limits_{x\to\pm\infty}f(x)$};
\node[decision, below=of q1, align=center] (q2){Existence d'une\\valeur prise sur\\un intervalle ?};
\node[action, left=of q2, align=center] (tvi){Outil : \textbf{TVI}\\vérifier $f$ continue sur $[a,b]$,\\vérifier $k$ entre $f(a)$ et $f(b)$,\\conclure : il existe $c$};
\node[decision, right=of q2, align=center] (q3){Fonction non définie\\en $x_0$ mais limite\\finie en $x_0$ ?};
\node[action, right=of q3, align=center] (prol){Outil : \textbf{PROLONGEMENT}\\calculer $\ell=\lim\limits_{x\to x_0}f(x)$,\\poser $\tilde f(x_0)=\ell$ :\\$\tilde f$ est continue en $x_0$};
\node[action, below=of q2, align=center] (cv){Outil : \textbf{CONTINUITÉ}\\vérifier $\lim\limits_{x\to x_0}f(x)=f(x_0)$\\et en déduire les propriétés};

\draw[fleche] (start) -- (q1);
\draw[fleche] (q1) -- node[above, font=\tiny] {OUI} (lim);
\draw[fleche] (q1) -- node[right, font=\tiny] {NON} (q2);
\draw[fleche] (q2) -- node[above, font=\tiny] {OUI} (tvi);
\draw[fleche] (q2) -- node[above, font=\tiny] {NON} (q3);
\draw[fleche] (q3) -- node[above, font=\tiny] {OUI} (prol);
\draw[fleche] (q3) |- node[right, font=\tiny, pos=0.3] {NON} (cv);
\end{tikzpicture}
\end{center}
\legende{Organigramme de choix de l'outil : limite, TVI, prolongement ou continuité selon la question posée.}
\end{popfigure}

\begin{exemplebox}[Application de l'organigramme — quatre questions sur le même modèle]
Soit $h(t)=100-5t^2$ la hauteur (en m) d'un objet en chute libre lâché depuis $100$ m, pour $t\in[0,\sqrt{20}]$ (l'objet touche le sol à $t=\sqrt{20}\approx 4,47$ s).
\begin{enumerate}
\item \textbf{« Quelle est la hauteur au bout de $2$ s ? »} Question sur une \textbf{valeur exacte en un point} : calcul direct $h(2)=100-20=80$ m (la continuité garantit que cette valeur a du sens).
\item \textbf{« L'objet est-il passé par la hauteur $50$ m ? »} Question d'\textbf{existence d'une valeur intermédiaire} : outil \textbf{TVI}. $h$ est continue sur $[0,\sqrt{20}]$, $h(0)=100$, $h(\sqrt{20})=0$, et $50\in[0,100]$, donc oui, il existe un instant $c$ avec $h(c)=50$.
\item \textbf{« Que se passe-t-il quand $t$ se rapproche de $\sqrt{20}$ ? »} Question sur le \textbf{comportement près d'un point} : outil \textbf{LIMITE}. $\lim\limits_{t\to\sqrt{20}}h(t)=0$ : l'objet approche du sol.
\item \textbf{« Peut-on définir la hauteur en $t=\sqrt{20}$ pour que le modèle soit continu ? »} Question de \textbf{régularisation} : outil \textbf{PROLONGEMENT}. La limite vaut $0$, on pose $\tilde h(\sqrt{20})=0$.
\end{enumerate}
\end{exemplebox}

\courssub{Exercices de synthèse}

\begin{exercice}[Tarification transport urbain]
Une société de transport facture un trajet selon la distance $d$ (en km) :
\begin{itemize}
\item si $0<d\le 5$ : forfait de $300$ F ;
\item si $5<d\le 15$ : $300 + 50(d-5)$ F ;
\item si $d>15$ : $800 + 40(d-15)$ F.
\end{itemize}
\begin{enumerate}
\item Écrire la fonction prix $P(d)$ par morceaux et préciser son ensemble de définition.
\item Étudier la continuité de $P$ en $d=5$ et en $d=15$. Interpréter concrètement.
\item Un client paie $1200$ F. Peut-on affirmer qu'il a parcouru \emph{exactement} $17$ km ? Justifier.
\item La fonction $P$ admet-elle un prolongement par continuité en $d=0$ ? Pourquoi ?
\end{enumerate}
\end{exercice}

\begin{corrige}[Exercice Tarification transport urbain]
\begin{enumerate}
\item $D_P = ]0,+\infty[$ et
\[
P(d)=\begin{cases}
300 & \text{si } 0<d\le 5\\
300+50(d-5) & \text{si } 5<d\le 15\\
800+40(d-15) & \text{si } d>15
\end{cases}
\]
\item En $d=5$ : $\lim\limits_{d\to 5^-}P(d)=300$, $\lim\limits_{d\to 5^+}P(d)=300+0=300$ et $P(5)=300$ : $P$ est \textbf{continue} en $5$. Pas de saut à la sortie du forfait.

En $d=15$ : $\lim\limits_{d\to 15^-}P(d)=300+50\times10=800$, $\lim\limits_{d\to 15^+}P(d)=800+0=800$ et $P(15)=800$ : $P$ est \textbf{continue} en $15$.

La tarification est donc continue sur $]0,+\infty[$ : le prix évolue sans « trou » ni saut brutal, ce qui est équitable pour le client.
\item On résout $P(d)=1200$ sur chaque tranche.

Sur la 2\up{e} tranche : $300+50(d-5)=1200 \iff 50(d-5)=900 \iff d=23$, mais $23\notin\,]5,15]$ : rejeté.

Sur la 3\up{e} tranche : $800+40(d-15)=1200 \iff 40(d-15)=400 \iff d=25$, et $25>15$ : accepté.

\textbf{Seule solution : $d=25$ km.} On ne peut pas affirmer $17$ km : d'ailleurs $P(17)=800+40\times2=880$ F $\neq 1200$ F.
\item $P$ n'est pas définie en $0$, mais $\lim\limits_{d\to 0^+}P(d)=300$ (limite finie). On peut donc prolonger $P$ par continuité en posant $\tilde P(0)=300$. Concrètement, cela correspond au prix d'accès au véhicule, même pour un trajet quasi nul.
\end{enumerate}
\end{corrige}

\begin{exercice}[Réservoir et TVI — rédaction complète]
Un réservoir contient $8$ m$^3$ d'eau. On ouvre une vanne au fond ; le volume d'eau restant $V(t)$ (en m$^3$) à l'instant $t$ (en minutes) est donné par $V(t)=8-0,5\sqrt{t}$ pour $t\in[0,256]$.
\begin{enumerate}
\item Justifier que $V$ est continue sur $[0,256]$.
\item Montrer qu'il existe un instant où le réservoir contient exactement $4$ m$^3$ d'eau.
\item Déterminer cet instant.
\item Rédiger la conclusion complète en suivant la méthode en 4 étapes.
\end{enumerate}
\end{exercice}

\begin{corrige}[Exercice Réservoir et TVI]
\begin{enumerate}
\item La fonction $t\mapsto\sqrt{t}$ est continue sur $[0,+\infty[$, donc sur $[0,256]$. $V$ est obtenue par multiplication par $-0,5$ et addition de $8$ : $V$ est \textbf{continue sur $[0,256]$}.
\item $V(0)=8$ et $V(256)=8-0,5\times\sqrt{256}=8-0,5\times16=0$. Comme $4\in[0,8]$ et que $V$ est continue sur $[0,256]$, le \textbf{TVI} garantit qu'il existe $c\in[0,256]$ tel que $V(c)=4$.
\item On résout :
\[
8-0,5\sqrt{c}=4 \iff 0,5\sqrt{c}=4 \iff \sqrt{c}=8 \iff c=64.
\]
\item \textbf{Rédaction en 4 étapes :}
\begin{enumerate}
\item \textbf{Hypothèses :} $V$ est continue sur $[0,256]$, $V(0)=8$, $V(256)=0$.
\item \textbf{Théorème :} TVI — une fonction continue sur un intervalle prend toute valeur comprise entre ses valeurs aux bornes.
\item \textbf{Calculs :} $4\in[0,8]$ ; la résolution de $V(c)=4$ donne $c=64\in[0,256]$.
\item \textbf{Conclusion :} le réservoir a contenu exactement $4$ m$^3$ d'eau au bout de $64$ minutes (soit $1$ h $04$ min) après l'ouverture de la vanne.
\end{enumerate}
\end{enumerate}
\end{corrige}

\begin{aretenir}[Check-list pour le Probatoire — Continuité et applications]
\begin{itemize}
\item[$\square$] Savoir \textbf{définir} la continuité en un point : $f$ définie en $x_0$, limite finie en $x_0$, et $\lim\limits_{x\to x_0}f(x)=f(x_0)$.
\item[$\square$] Savoir \textbf{appliquer} le TVI : vérifier la continuité sur $[a,b]$, vérifier que $k$ est entre $f(a)$ et $f(b)$, conclure à l'existence de $c$.
\item[$\square$] Savoir \textbf{calculer} une limite pour prolonger par continuité (forme indéterminée $\frac{0}{0}$ : factorisation, expression conjuguée).
\item[$\square$] \textbf{Modéliser} : choisir une fonction continue (grandeur physique) ou discontinue (tarifs par paliers, seuils).
\item[$\square$] \textbf{Rédiger} en 4 étapes : Hypothèses — Théorème — Calculs — Conclusion.
\item[$\square$] Toujours préciser \textbf{l'ensemble de définition} du modèle dans un contexte concret.
\end{itemize}
\end{aretenir}

\newpage

\coursec{Activité d'intégration}

\coursec{Limites et continuité d'une fonction}

\courssub{Objectifs de l'activité}

À l'issue de cette activité d'intégration, l'élève sera capable de :
\begin{itemize}
\item calculer la \cle{limite} d'une fonction en un point par \cle{factorisation} (cas d'une forme indéterminée du type $\frac{0}{0}$) ;
\item définir et justifier un \cle{prolongement par continuité} d'une fonction en un point où elle n'est pas définie ;
\item vérifier qu'une fonction est \cle{continue sur un intervalle} ;
\item appliquer le \cle{théorème des valeurs intermédiaires} (TVI) pour justifier l'existence d'une solution à un problème concret ;
\item \cle{rédiger} une démarche complète, justifiée et argumentée, destinée à un destinataire non mathématicien (un technicien de forage).
\end{itemize}

\courssub{Situation}

\textbf{Le forage du quartier : à la recherche du bon débit.}

Dans le quartier de Banengo, à la périphérie de Bafoussam (région de l'Ouest), la communauté vient d'installer un forage équipé d'une pompe électrique pour alimenter les habitants en eau potable. Le technicien hydraulicien de la commune, M.~Tchoupo, effectue des relevés lors de la mise en service.

Après analyse de ses mesures, il constate que le débit instantané de la pompe, exprimé en litres par minute (\si{L/min}), $t$ minutes après la mise en route, est bien modélisé par la fonction :
\[ d(t) = \frac{t^2 - 9}{t - 3}, \qquad \text{valable pour } t \in [0\,;\,10] \text{ et } t \neq 3. \]

Problème : la formule ne donne \textbf{aucune valeur} à l'instant précis $t = 3$ minutes, car le dénominateur s'annule. Or le manomètre du technicien n'a jamais montré de saut brutal : le débit semble varier de façon régulière. M.~Tchoupo affirme donc qu'on peut attribuer une valeur \emph{raisonnable} au débit à l'instant $t = 3$.

Par ailleurs, la quantité totale d'eau pompée depuis la mise en route, en litres, est modélisée par une fonction $Q$ \textbf{continue} sur $[0\,;\,10]$, avec :
\[ Q(0) = \SI{0}{L} \qquad \text{et} \qquad Q(10) = \SI{120}{L}. \]
Le réservoir du quartier a une capacité de \SI{120}{L}. Le chef de quartier veut savoir s'il y a eu un instant où le réservoir était \textbf{exactement à moitié plein}, c'est-à-dire contenant exactement \SI{60}{L}.

Votre groupe d'élèves de Première technique est sollicité pour répondre, par écrit et de façon rigoureuse, aux questions du technicien.

\courssub{Consignes}

Travail en groupes de 3 ou 4 élèves. Durée : une séance, suivie d'une restitution et d'une correction collective de la rédaction.

\begin{enumerate}
\item \textbf{Étude du débit en $t = 3$.}
\begin{enumerate}
\item Vérifier que le calcul direct de $d(3)$ conduit à une forme indéterminée. Préciser laquelle.
\item Factoriser le numérateur $t^2 - 9$, simplifier l'expression de $d(t)$ pour $t \neq 3$, puis calculer $\displaystyle\lim_{t \to 3} d(t)$.
\end{enumerate}
\item \textbf{Prolongement par continuité.} Expliquer pourquoi $d$ admet un prolongement par continuité en $3$. Écrire l'expression complète de la fonction prolongée $\tilde{d}$ sur $[0\,;\,10]$ (valeur en $3$ comprise).
\item \textbf{Continuité sur l'intervalle.} Justifier que la fonction prolongée $\tilde{d}$ est continue sur tout l'intervalle $[0\,;\,10]$.
\item \textbf{Le réservoir à moitié plein.} En appliquant le théorème des valeurs intermédiaires à la fonction $Q$, justifier qu'il existe au moins un instant $t_0 \in [0\,;\,10]$ où le réservoir contenait exactement \SI{60}{L}. Vérifier au préalable que toutes les hypothèses du théorème sont réunies.
\item \textbf{Rédaction finale.} Rédiger une réponse complète à M.~Tchoupo (10 à 15 lignes) : valeur du débit à $t = 3$, justification du prolongement, et réponse au chef de quartier. Chaque affirmation doit être appuyée par un calcul ou un théorème cité.
\end{enumerate}

\courssub{Ressources à mobiliser}

\begin{aretenir}[Boîte à outils de l'activité]
\begin{itemize}
\item \textbf{Limite en un point :} $\displaystyle\lim_{t \to a} f(t) = \ell$ signifie que $f(t)$ est aussi proche que l'on veut de $\ell$ lorsque $t$ est suffisamment proche de $a$.
\item \textbf{Lever une indétermination $\frac{0}{0}$ :} factoriser (identité remarquable $a^2 - b^2 = (a-b)(a+b)$), simplifier, puis calculer la limite.
\item \textbf{Continuité en un point :} $f$ est continue en $a$ si $f$ est définie en $a$ et $\displaystyle\lim_{t \to a} f(t) = f(a)$.
\item \textbf{Prolongement par continuité :} si $f$ n'est pas définie en $a$ mais admet une limite finie $\ell$ en $a$, on prolonge $f$ par continuité en posant $f(a) = \ell$.
\item \textbf{Théorème des valeurs intermédiaires :} si $f$ est continue sur $[a\,;\,b]$, alors pour tout réel $k$ compris entre $f(a)$ et $f(b)$, il existe au moins un réel $c \in [a\,;\,b]$ tel que $f(c) = k$.
\end{itemize}
\end{aretenir}

\begin{propriete}[Complément : limites infinies et limites à l'infini]
\textbf{Rappels utiles pour la suite du programme (Leçon 24).}
\begin{itemize}
\item \textbf{Limite infinie en un point :} $\displaystyle\lim_{x \to a} f(x) = +\infty$ (resp. $-\infty$) signifie que $f(x)$ devient arbitrairement grand (resp. petit) quand $x$ se rapproche de $a$. Exemple : $\lim_{x \to 0} \frac{1}{x^2} = +\infty$.
\item \textbf{Limite à l'infini :} $\displaystyle\lim_{x \to +\infty} f(x) = \ell$ signifie que $f(x)$ se rapproche de $\ell$ quand $x$ devient arbitrairement grand. Exemple : $\lim_{x \to +\infty} \frac{1}{x} = 0$.
\item \textbf{Règles des formes indéterminées à l'infini :} $\frac{\infty}{\infty}$, $\infty - \infty$, $0 \times \infty$ nécessitent une transformation (factorisation par le terme de plus haut degré, changement de variable, etc.).
\end{itemize}
\end{propriete}

\courssub{Démarche et corrigé commenté}

\begin{popfigure}
\begin{tikzpicture}[>=Stealth]
\begin{axis}[
    axis lines=middle,
    xlabel=$t$ (min), ylabel=$d(t)$ (\si{L/min}),
    xmin=-0.5, xmax=10.5, ymin=2, ymax=14,
    xtick={0,3,10}, ytick={3,6,13},
    domain=0:10, samples=100,
    width=\linewidth, height=0.55\linewidth,
    clip=false,
    tick label style={font=\small},
    label style={font=\small}
]
% Courbe de la fonction affine y = x+3
\addplot[popcurve, thick, popblue] {x+3};
% Trou pour la fonction d initiale (non définie en 3)
\addplot[only marks, mark=*, mark size=3, popblue, fill=white, thick] coordinates {(3,6)};
% Point plein pour le prolongement
\addplot[only marks, mark=*, mark size=3, poporange, thick] coordinates {(3,6)};
% Lignes pointillées pour repérer le point
\addplot[dashed, popmuted, domain=0:3] coordinates {(3,0) (3,6)};
\addplot[dashed, popmuted, domain=0:6] coordinates {(0,6) (3,6)};
% Annotations
\node[popblue, anchor=south east, font=\small] at (axis cs:2.8,6.2) {$d$ : trou};
\node[poporange, anchor=north east, font=\small] at (axis cs:2.8,5.8) {$\tilde{d}$ : plein};
\node[popblue, anchor=south west, font=\small] at (axis cs:3.2,6.2) {$(3;6)$};
\end{axis}
\end{tikzpicture}
\legende{Graphique de $d(t)$ (trou en $t=3$) et de son prolongement $\tilde{d}(t)=t+3$ (point plein).}
\end{popfigure}

\begin{exoresolu}[Le forage de Banengo : corrigé complet]
Énoncé : $d(t) = \dfrac{t^2 - 9}{t - 3}$ pour $t \in [0\,;\,10]$, $t \neq 3$ ; $Q$ continue sur $[0\,;\,10]$ avec $Q(0) = \SI{0}{L}$ et $Q(10) = \SI{120}{L}$. Questions : limite de $d$ en 3, prolongement par continuité, continuité sur $[0\,;\,10]$, existence d'un instant où $Q = \SI{60}{L}$, rédaction au technicien.
\tcblower
\textbf{Étape 1 : limite de $d$ en 3.}

Pour $t = 3$ : le numérateur vaut $3^2 - 9 = 0$ et le dénominateur vaut $3 - 3 = 0$. On obtient la forme indéterminée $\frac{0}{0}$ : le calcul direct est impossible, il faut transformer l'expression.

On factorise le numérateur grâce à l'identité remarquable $a^2 - b^2 = (a-b)(a+b)$ :
\[ t^2 - 9 = t^2 - 3^2 = (t - 3)(t + 3). \]
Pour tout $t \neq 3$, on peut alors simplifier par $(t-3)$, qui est non nul :
\[ d(t) = \frac{(t-3)(t+3)}{t-3} = t + 3. \]
Par conséquent :
\[ \lim_{t \to 3} d(t) = \lim_{t \to 3} (t + 3) = 6. \]
La limite existe et est \textbf{finie} : $\displaystyle\lim_{t \to 3} d(t) = \SI{6}{L/min}$.

\textbf{Étape 2 : prolongement par continuité.}

La fonction $d$ n'est pas définie en $3$, mais elle y admet une limite finie égale à $6$. D'après le théorème de prolongement par continuité, on peut définir la fonction $\tilde{d}$ sur $[0\,;\,10]$ tout entier par :
\[ \tilde{d}(t) = \begin{cases} \dfrac{t^2 - 9}{t - 3} = t + 3 & \text{si } t \neq 3, \\ 6 & \text{si } t = 3. \end{cases} \]
En réalité, $\tilde{d}(t) = t + 3$ pour \emph{tout} $t \in [0\,;\,10]$, car $t+3$ vaut exactement $6$ en $t = 3$. Le prolongement est donc simplement la fonction affine $\tilde{d}(t) = t + 3$.

\textbf{Étape 3 : continuité sur $[0\,;\,10]$.}

La fonction $\tilde{d}(t) = t + 3$ est une fonction \cle{affine} (fonction polynôme de degré 1). Or toute fonction polynôme est continue sur $\mathbb{R}$, donc en particulier sur $[0\,;\,10]$. De plus, en $t = 3$ :
\[ \lim_{t \to 3} \tilde{d}(t) = 6 = \tilde{d}(3), \]
ce qui confirme la continuité au point de raccordement. Conclusion : $\tilde{d}$ est \textbf{continue sur $[0\,;\,10]$}.

\begin{popfigure}
\begin{tikzpicture}[>=Stealth]
\begin{axis}[
    axis lines=middle,
    xlabel=$t$ (min), ylabel=$Q(t)$ (L),
    xmin=-0.5, xmax=10.5, ymin=-5, ymax=130,
    xtick={0,10}, ytick={0,60,120},
    domain=0:10, samples=100,
    width=\linewidth, height=0.55\linewidth,
    clip=false,
    tick label style={font=\small},
    label style={font=\small}
]
% Courbe continue générique Q(t) = 1.2 t^2 (vérifie Q(0)=0, Q(10)=120)
\addplot[popcurve, thick, popgreen, domain=0:10] {1.2*x^2};
% Ligne horizontale y = 60
\addplot[dashed, popmuted, domain=0:10] {60};
% Point d'intersection t0 = sqrt(50) approx 7.07
\addplot[only marks, mark=*, mark size=3, popred, thick] coordinates {(7.071,60)};
% Lignes pointillées vers les axes
\addplot[dashed, popmuted] coordinates {(7.071,0) (7.071,60)};
\addplot[dashed, popmuted] coordinates {(0,60) (7.071,60)};
% Annotations
\node[popgreen, anchor=south west, font=\small] at (axis cs:2,50) {$Q$ continue};
\node[popred, anchor=south, font=\small] at (axis cs:7.071,60) {$t_0$};
\node[popmuted, anchor=east, font=\small] at (axis cs:0,60) {\SI{60}{L}};
\node[popmuted, anchor=north, font=\small] at (axis cs:7.071,0) {$t_0$};
\end{axis}
\end{tikzpicture}
\legende{Illustration du TVI : la courbe continue $Q$ relie $(0,0)$ à $(10,120)$ et coupe nécessairement la ligne $y=60$.}
\end{popfigure}

\textbf{Étape 4 : le réservoir à moitié plein (TVI).}

Vérifions les hypothèses du théorème des valeurs intermédiaires pour la fonction $Q$ :
\begin{itemize}
\item $Q$ est \textbf{continue} sur $[0\,;\,10]$ (donnée de l'énoncé) ;
\item $Q(0) = \SI{0}{L}$ et $Q(10) = \SI{120}{L}$ ;
\item le nombre $k = \SI{60}{L}$ est compris entre $Q(0)$ et $Q(10)$ : en effet, $0 < 60 < 120$.
\end{itemize}
Toutes les hypothèses sont réunies. D'après le TVI, il existe \textbf{au moins un} instant $t_0 \in [0\,;\,10]$ tel que :
\[ Q(t_0) = \SI{60}{L}. \]
Il y a donc eu au moins un moment où le réservoir contenait exactement \SI{60}{L}, c'est-à-dire était exactement à moitié plein.

\textbf{Étape 5 : rédaction de la réponse à M.~Tchoupo (modèle attendu).}

\emph{« Monsieur Tchoupo, votre formule $d(t) = \frac{t^2-9}{t-3}$ ne permet pas de calculer directement le débit à $t = \SI{3}{min}$, car numérateur et dénominateur s'annulent simultanément (forme $\frac{0}{0}$). Cependant, en factorisant $t^2 - 9 = (t-3)(t+3)$ et en simplifiant, on obtient $d(t) = t + 3$ pour tout $t \neq 3$. La limite en $3$ existe et vaut $\SI{6}{L/min}$ : on peut donc prolonger le modèle par continuité en posant $d(3) = \SI{6}{L/min}$, ce qui correspond bien à l'absence de saut brutal observée sur votre manomètre. La fonction ainsi prolongée, $\tilde{d}(t) = t + 3$, est affine donc continue sur tout l'intervalle $[0\,;\,10]$. Concernant le réservoir : la quantité pompée $Q$ est continue sur $[0\,;\,10]$, avec $Q(0) = \SI{0}{L}$ et $Q(10) = \SI{120}{L}$. Comme $\SI{60}{L}$ est compris entre $0$ et $120$, le théorème des valeurs intermédiaires garantit qu'il existe un instant $t_0$ où $Q(t_0) = \SI{60}{L}$ : le réservoir est donc passé par un état exactement à moitié plein. »}
\end{exoresolu}

\courssub{Bilan de l'activité}

Cette activité a permis de mobiliser, dans une situation authentique de la vie technique camerounaise, l'ensemble des savoirs de la leçon :
\begin{itemize}
\item le calcul d'une \cle{limite} en un point par factorisation, en levant une forme indéterminée $\frac{0}{0}$ ;
\item le \cle{prolongement par continuité}, qui donne un sens mathématique précis à l'idée intuitive du technicien (« pas de saut brutal ») ;
\item la vérification de la \cle{continuité sur un intervalle}, en distinguant bien la fonction initiale $d$ (non définie en 3) de la fonction prolongée $\tilde{d}$ ;
\item le \cle{théorème des valeurs intermédiaires}, outil d'existence puissant : il assure qu'une grandeur continue passant de $0$ à $120$ franchit nécessairement toutes les valeurs intermédiaires, dont $60$ ;
\item enfin, la \cle{rédaction d'une démarche justifiée} : chaque conclusion pratique (débit de $\SI{6}{L/min}$, réservoir à moitié plein) repose sur un calcul ou un théorème explicitement cité, compétence essentielle au Probatoire et au Baccalauréat technique comme dans la pratique professionnelle.
\end{itemize}

\begin{attention}[Erreurs fréquentes relevées pendant la correction collective]
\begin{itemize}
\item Écrire $d(3) = \frac{0}{0} = 1$ ou $= 0$ : une forme indéterminée n'est \textbf{pas} un résultat, c'est un signal qu'il faut transformer l'écriture.
\item Simplifier par $(t-3)$ sans préciser la condition $t \neq 3$ : la simplification n'est licite que parce que le facteur est non nul.
\item Conclure « $d$ est continue en 3 » sans avoir prolongé : $d$ n'existe pas en 3, c'est $\tilde{d}$ qui est continue.
\item Appliquer le TVI sans vérifier la \textbf{continuité} de $Q$ ni l'encadrement $0 < 60 < 120$ : citer les hypothèses fait partie de la justification exigée.
\item Affirmer que le TVI donne la \emph{valeur} de $t_0$ : le théorème garantit seulement l'\textbf{existence} d'un tel instant, pas son calcul.
\end{itemize}
\end{attention}

\begin{savaistu}[L'eau potable, un enjeu national]
Selon les programmes camerounais d'hydraulique rurale, des milliers de forages sont équipés chaque année dans les régions de l'Ouest, du Nord et de l'Adamaoua. Le dimensionnement des pompes et des réservoirs repose sur des modèles mathématiques de débit et de volume : les notions de continuité et de valeurs intermédiaires, loin d'être abstraites, garantissent concrètement qu'un réservoir se remplit progressivement et atteint tous les niveaux intermédiaires. Les techniciens hydrauliciens formés dans les lycées techniques du Cameroun utilisent quotidiennement ce type de raisonnement.
\end{savaistu}

\newpage

\coursec{Méthodes et exercices résolus}

\coursec{Limites et continuité d'une fonction}

\courssub{Limite d'une fonction en un point}

\begin{definition}[Limite finie en un point]
Soit $f$ une fonction définie sur un intervalle $I$ et $a \in I$ (ou $a$ borne de $I$). On dit que $f$ admet la \cle{limite finie} $\ell \in \mathbb{R}$ en $a$ si, lorsque $x$ se rapproche de $a$ (par valeurs supérieures, inférieures, ou les deux), $f(x)$ se rapproche arbitrairement de $\ell$. On note :
\[ \lim_{x \to a} f(x) = \ell. \]
Si $x$ tend vers $a$ par valeurs $> a$ (resp. $< a$), on note $\lim_{x \to a^+} f(x)$ (resp. $\lim_{x \to a^-} f(x)$).
\end{definition}

\begin{propriete}[Opérations sur les limites (rappel)]
Si $\lim_{x\to a} f(x) = \ell_1$ et $\lim_{x\to a} g(x) = \ell_2$ (finies), alors :
\begin{itemize}
\item $\lim_{x\to a} (f+g)(x) = \ell_1 + \ell_2$,
\item $\lim_{x\to a} (f \times g)(x) = \ell_1 \ell_2$,
\item $\lim_{x\to a} \dfrac{f}{g}(x) = \dfrac{\ell_1}{\ell_2}$ si $\ell_2 \neq 0$.
\end{itemize}
Pour les fonctions usuelles (polynômes, rationnelles, $\sqrt{x}$, $\sin$, $\cos$, $\exp$, $\ln$) \cle{continues sur leur domaine}, $\lim_{x\to a} f(x) = f(a)$ dès que $a$ appartient au domaine.
\end{propriete}

\begin{methode}[Calculer $\lim_{x\to a} f(x)$]
Pour déterminer la limite d'une fonction $f$ en un point $a$ :
\begin{enumerate}
\item \textbf{Vérifier que $a$ appartient au domaine de définition} de $f$ (ou en est une borne d'adhérence).
\item \textbf{Substituer} : si $f$ est une fonction usuelle continue en $a$ (polynôme, rationnelle définie en $a$, $\sqrt{x}$, $\sin$, $\cos$, etc.), alors $\lim_{x\to a} f(x) = f(a)$.
\item Si l'on obtient une \textbf{forme indéterminée} du type $\dfrac{0}{0}$ : factoriser le numérateur et le dénominateur par $(x-a)$ (identités remarquables, factorisation), simplifier, puis substituer à nouveau.
\item Si le dénominateur tend vers $0$ et le numérateur vers un nombre \textbf{non nul} : étudier le \textbf{signe du dénominateur} à gauche ($x \to a^-$) et à droite ($x \to a^+$) de $a$ pour conclure ($+\infty$ ou $-\infty$).
\item \textbf{Conclure par une phrase} : « donc $\lim_{x\to a} f(x) = \ldots$ ».
\end{enumerate}
\textbf{Réflexe} : une forme $\dfrac{0}{0}$ n'est jamais une réponse finale : c'est un signal qu'il faut \emph{transformer l'expression} (factoriser, rationaliser, etc.).
\end{methode}

\begin{exoresolu}[Limite en un point avec forme indéterminée]
Soit $f$ la fonction définie par $f(x)=\dfrac{x^{2}-9}{x-3}$.
\begin{enumerate}
\item Déterminer le domaine de définition de $f$.
\item Calculer $\lim_{x\to 3} f(x)$.
\end{enumerate}
\tcblower
\textbf{1. Domaine de définition.}

$f$ est une fonction rationnelle : elle est définie lorsque son dénominateur est non nul.
\[ x-3=0 \iff x=3. \]
Donc $D_f = \mathbb{R}\setminus\{3\} = ]-\infty\,;\,3[\,\cup\,]3\,;\,+\infty[$.

\textbf{2. Calcul de la limite.}

Substituons $x=3$ : le numérateur vaut $3^{2}-9=0$ et le dénominateur vaut $3-3=0$. On obtient la forme indéterminée $\dfrac{0}{0}$ : il faut factoriser.

Le numérateur est une identité remarquable :
\[ x^{2}-9 = (x-3)(x+3). \]
Pour tout $x\neq 3$, on peut donc simplifier par $(x-3)$ :
\[ f(x)=\dfrac{(x-3)(x+3)}{x-3}=x+3. \]
Ainsi :
\[ \lim_{x\to 3} f(x) = \lim_{x\to 3}(x+3) = 3+3 = 6. \]

\textbf{Conclusion} : $\lim_{x\to 3} f(x) = 6$. La fonction $f$ n'est pas définie en $3$, mais elle admet une limite finie en $3$.
\end{exoresolu}

\courssub{Limites infinies et limites à l'infini}

\begin{methode}[Limite du type $\dfrac{\text{non nul}}{0}$ et asymptote verticale]
Pour étudier $\lim_{x\to a} f(x)$ lorsque le dénominateur tend vers $0$ et le numérateur vers un nombre non nul $L$ :
\begin{enumerate}
\item Calculer la limite du \textbf{numérateur} $L = \lim_{x\to a} \text{num}(x)$ (non nul).
\item Étudier le \textbf{signe du dénominateur} pour $x \to a^-$ (valeurs $< a$) et $x \to a^+$ (valeurs $> a$), à l'aide d'un tableau de signes si nécessaire.
\item Appliquer la règle des signes :
  \begin{itemize}
  \item $\dfrac{L}{0^{+}} = +\infty$ si $L>0$, $-\infty$ si $L<0$,
  \item $\dfrac{L}{0^{-}} = -\infty$ si $L>0$, $+\infty$ si $L<0$.
  \end{itemize}
\item Conclure : si la limite est infinie (au moins d'un côté), la droite d'équation $x=a$ est \textbf{asymptote verticale} à la courbe représentative de $f$.
\end{enumerate}
\textbf{À retenir} : $0^{+}$ signifie « tend vers $0$ par valeurs positives », $0^{-}$ « par valeurs négatives ».
\end{methode}

\begin{exoresolu}[Limite infinie et asymptote verticale]
Soit $g$ la fonction définie sur $]1\,;\,+\infty[$ par $g(x)=\dfrac{2x+1}{x-1}$.
\begin{enumerate}
\item Calculer $\lim_{x\to 1,\ x>1} g(x)$.
\item Interpréter graphiquement le résultat.
\end{enumerate}
\tcblower
\textbf{1. Calcul de la limite.}

\underline{Limite du numérateur} : $\lim_{x\to 1}(2x+1)=2\times 1+1=3 = L \neq 0$.

\underline{Limite du dénominateur} : $\lim_{x\to 1}(x-1)=0$.

Le numérateur tend vers $3>0$ : la limite de $g$ est donc infinie. Il reste à déterminer le signe.

Comme on s'approche de $1$ par valeurs supérieures ($x>1$), on a $x-1>0$ : le dénominateur tend vers $0^{+}$.

Par quotient : $\dfrac{3}{0^{+}} = +\infty$.
\[ \lim_{x\to 1^+} g(x) = +\infty. \]

\textbf{2. Interprétation graphique.}

Puisque la limite de $g$ en $1$ (par la droite) est $+\infty$, la droite d'équation $x=1$ est \textbf{asymptote verticale} à la courbe représentative de $g$ : la courbe « monte » vers $+\infty$ en se rapprochant de cette droite sans jamais la toucher.

\begin{popfigure}
\begin{tikzpicture}
\begin{axis}[
    width=\linewidth, height=6cm,
    axis lines=middle,
    xlabel=$x$, ylabel=$y$,
    xmin=0.5, xmax=4, ymin=-2, ymax=10,
    xtick={1,2,3}, ytick={2,4,6,8},
    domain=1.01:4, samples=200,
    clip=false,
    grid=major
]
\addplot[popcurve, thick, domain=1.01:4] { (2*x+1)/(x-1) };
% Asymptote verticale
\draw[poporange, dashed, thick] (axis cs:1,-2) -- (axis cs:1,10) node[above, pos=0.9, poporange] {$x=1$};
% Point exclus
\addplot[only marks, mark=*, mark size=2, poporange, fill=white] coordinates {(1,3)};
\end{axis}
\end{tikzpicture}
\legende{Courbe de $g(x)=\frac{2x+1}{x-1}$ avec asymptote verticale $x=1$.}
\end{popfigure}

\textbf{Conclusion} : $\lim_{x\to 1^{+}} g(x)=+\infty$ et la droite $x=1$ est asymptote verticale à la courbe.
\end{exoresolu}

\begin{methode}[Limite en $+\infty$ ou $-\infty$ et asymptote horizontale]
\begin{enumerate}
\item \textbf{Fonction polynôme} : la limite en l'infini est celle de son \textbf{terme de plus haut degré} (coefficient dominant $\times$ puissance).
\item \textbf{Fonction rationnelle} $f(x)=\frac{P(x)}{Q(x)}$ : la limite en l'infini est celle du \textbf{quotient des termes de plus haut degré}.
  \begin{itemize}
  \item $\deg P < \deg Q \implies \lim = 0$ (asymptote horizontale $y=0$).
  \item $\deg P = \deg Q \implies \lim = \frac{\text{coeff dominant de }P}{\text{coeff dominant de }Q}$ (asymptote horizontale $y=\ell$).
  \item $\deg P > \deg Q \implies \lim = \pm\infty$ (pas d'asymptote horizontale, éventuellement oblique).
  \end{itemize}
\item \textbf{Technique universelle} : factoriser par $x^{n}$ (plus haut degré du dénominateur ou du numérateur), simplifier, puis faire tendre $x$ vers l'infini en utilisant $\lim_{x\to\pm\infty}\dfrac{1}{x^{n}}=0$.
\item Si la limite est un nombre fini $\ell$, la droite $y=\ell$ est \textbf{asymptote horizontale}.
\end{enumerate}
\end{methode}

\begin{exoresolu}[Limite d'une fonction rationnelle en $+\infty$]
Un atelier produit des pièces. Le coût moyen de production, en milliers de francs CFA, de $x$ centaines de pièces est modélisé par :
\[ C(x)=\dfrac{3x^{2}+2x+1}{x^{2}+5}, \qquad x\geq 1. \]
\begin{enumerate}
\item Calculer $\lim_{x\to +\infty} C(x)$.
\item Interpréter le résultat pour l'atelier.
\end{enumerate}
\tcblower
\textbf{1. Calcul de la limite.}

En $+\infty$, numérateur et dénominateur tendent vers $+\infty$ : forme indéterminée $\dfrac{\infty}{\infty}$. On factorise par le terme de plus haut degré, $x^{2}$ :
\[ C(x)=\dfrac{x^{2}\left(3+\dfrac{2}{x}+\dfrac{1}{x^{2}}\right)}{x^{2}\left(1+\dfrac{5}{x^{2}}\right)}=\dfrac{3+\dfrac{2}{x}+\dfrac{1}{x^{2}}}{1+\dfrac{5}{x^{2}}}. \]
Or $\lim_{x\to+\infty}\dfrac{2}{x}=0$, $\lim_{x\to+\infty}\dfrac{1}{x^{2}}=0$ et $\lim_{x\to+\infty}\dfrac{5}{x^{2}}=0$. Donc :
\[ \lim_{x\to+\infty} C(x)=\dfrac{3+0+0}{1+0}=3. \]

\textbf{2. Interprétation.}

Lorsque la production devient très grande, le coût moyen se rapproche de \SI{3000}{\text{FCFA}} par centaine de pièces. Graphiquement, la droite $y=3$ est asymptote horizontale à la courbe de $C$.

\begin{popfigure}
\begin{tikzpicture}
\begin{axis}[
    width=\linewidth, height=6cm,
    axis lines=middle,
    xlabel=$x$ (centaines de pièces), ylabel=$C(x)$ (kFCFA),
    xmin=0, xmax=10, ymin=0, ymax=5,
    xtick={1,2,5,10}, ytick={1,2,3,4},
    domain=1:10, samples=200,
    grid=major
]
\addplot[popcurve, thick] { (3*x^2+2*x+1)/(x^2+5) };
% Asymptote horizontale
\draw[poporange, dashed, thick] (axis cs:0,3) -- (axis cs:10,3) node[right, poporange] {$y=3$};
\end{axis}
\end{tikzpicture}
\legende{Courbe du coût moyen $C(x)$ et son asymptote horizontale $y=3$.}
\end{popfigure}

\textbf{Conclusion} : $\lim_{x\to+\infty} C(x)=3$ : le coût moyen se stabilise à \SI{3000}{\text{FCFA}}.
\end{exoresolu}

\courssub{Continuité en un point et sur un intervalle}

\begin{definition}[Continuité en un point]
Soit $f$ définie sur un intervalle $I$ et $a \in I$. La fonction $f$ est \cle{continue en $a$} si et seulement si les trois conditions suivantes sont réunies :
\begin{enumerate}
\item $f$ est \textbf{définie en $a$} : $f(a)$ existe ;
\item $f$ admet une \textbf{limite finie en $a$} (les limites à gauche et à droite existent et sont égales) : $\lim_{x\to a} f(x) = \ell \in \mathbb{R}$ ;
\item $\lim_{x\to a} f(x) = f(a)$.
\end{enumerate}
Si l'une des conditions échoue, $f$ est \textbf{discontinue en $a$}.
\end{definition}

\begin{definition}[Continuité sur un intervalle]
Une fonction $f$ est \cle{continue sur un intervalle $I$} si elle est continue en \emph{chaque point} de $I$.
\end{definition}

\begin{propriete}[Opérations sur les fonctions continues]
Si $f$ et $g$ sont continues sur un intervalle $I$, alors :
\begin{itemize}
\item $f+g$, $f-g$, $f \times g$ sont continues sur $I$,
\item $\dfrac{f}{g}$ est continue sur $I \setminus \{x \in I \mid g(x)=0\}$,
\item La fonction composée $f \circ g$ est continue sur $I$ (si $g(I) \subset \text{Dom}(f)$).
\end{itemize}
Les fonctions usuelles (polynômes, rationnelles sur leur domaine, $\sqrt{x}$, $\sin$, $\cos$, $\exp$, $\ln$) sont continues sur chaque intervalle inclus dans leur domaine de définition.
\end{propriete}

\begin{methode}[Étudier la continuité d'une fonction en $a$]
\begin{enumerate}
\item Vérifier que $f(a)$ \textbf{existe} (appartenance au domaine).
\item Calculer $\lim_{x\to a^-} f(x)$ et $\lim_{x\to a^+} f(x)$.
\item Comparer : si les deux limites sont égales à un réel $\ell$ \textbf{et} $\ell = f(a)$, alors $f$ est continue en $a$. Sinon, elle est discontinue.
\end{enumerate}
\textbf{Cas particulier (fonction définie par morceaux)} : calculer la limite à gauche avec l'expression de gauche, la limite à droite avec l'expression de droite, et comparer avec $f(a)$.
\end{methode}

\begin{exoresolu}[Continuité d'une fonction définie par morceaux]
Soit $f$ la fonction définie sur $\mathbb{R}$ par :
\[ f(x)=\begin{cases} 2x+1 & \text{si } x\leq 2 \\ x^{2}-1 & \text{si } x>2 \end{cases} \]
La fonction $f$ est-elle continue en $2$ ?
\tcblower
\textbf{Condition 1 :} $f$ est définie en $2$. Comme $2\leq 2$, on utilise la première expression :
\[ f(2)=2\times 2+1=5. \]

\textbf{Condition 2 : limites à gauche et à droite.}

\underline{Limite à gauche} ($x<2$, donc $f(x)=2x+1$) :
\[ \lim_{x\to 2^-} f(x)=\lim_{x\to 2}(2x+1)=5. \]

\underline{Limite à droite} ($x>2$, donc $f(x)=x^{2}-1$) :
\[ \lim_{x\to 2^+} f(x)=\lim_{x\to 2}(x^{2}-1)=4-1=3. \]

Les limites à gauche ($5$) et à droite ($3$) sont \textbf{différentes} : $f$ n'admet pas de limite en $2$.

\textbf{Condition 3 :} non satisfaite puisque la limite n'existe pas.

\textbf{Conclusion} : $f$ est \textbf{discontinue en $2$}. Graphiquement, la courbe présente un « saut » au point d'abscisse $2$ : on ne peut pas la tracer sans lever le crayon.

\begin{popfigure}
\begin{tikzpicture}
\begin{axis}[
    width=\linewidth, height=6cm,
    axis lines=middle,
    xlabel=$x$, ylabel=$y$,
    xmin=0, xmax=4, ymin=0, ymax=8,
    xtick={1,2,3}, ytick={1,3,5,7},
    domain=0:2, samples=100,
    grid=major
]
% Branche gauche : 2x+1
\addplot[popcurve, thick, domain=0:2] {2*x+1};
% Point plein à (2,5)
\addplot[only marks, mark=*, mark size=2, popblue] coordinates {(2,5)};
% Branche droite : x^2-1
\addplot[popcurve, thick, domain=2:4] {x^2-1};
% Point vide à (2,3)
\addplot[only marks, mark=*, mark size=2, poporange, fill=white] coordinates {(2,3)};
% Flèches indicatives
\draw[popmuted, <->] (axis cs:2,3.2) -- (axis cs:2,4.8) node[midway, right] {saut = 2};
\end{axis}
\end{tikzpicture}
\legende{Discontinuité de première espèce (saut) en $x=2$.}
\end{popfigure}
\end{exoresolu}

\courssub{Propriétés des fonctions continues : Théorème des valeurs intermédiaires}

\begin{propriete}[Théorème des valeurs intermédiaires (TVI)]
Soit $f$ une fonction \textbf{continue sur un intervalle fermé $[a\,;\,b]$}. Pour tout réel $k$ compris entre $f(a)$ et $f(b)$, il existe \textbf{au moins un} $c \in [a\,;\,b]$ tel que $f(c)=k$.
En particulier, si $f(a)$ et $f(b)$ sont de \textbf{signes contraires} ($f(a) \times f(b) < 0$), alors l'équation $f(x)=0$ admet au moins une solution dans $[a\,;\,b]$.
\end{propriete}

\begin{methode}[Montrer qu'une équation admet une solution (TVI)]
Pour montrer que l'équation $f(x)=k$ admet au moins une solution dans $[a\,;\,b]$ :
\begin{enumerate}
\item \textbf{Justifier la continuité} de $f$ sur $[a\,;\,b]$ (fonction polynôme, somme/produit/composée de fonctions continues, etc.) ;
\item \textbf{Calculer $f(a)$ et $f(b)$} ;
\item \textbf{Vérifier que $k$ est compris entre $f(a)$ et $f(b)$} (ou $f(a)f(b)<0$ pour $k=0$) ;
\item \textbf{Conclure} par le TVI : « il existe au moins un réel $c\in[a\,;\,b]$ tel que $f(c)=k$ ».
\end{enumerate}
\textbf{Attention} : Le TVI garantit l'\emph{existence}, pas l'unicité (pour l'unicité, il faut ajouter que $f$ est strictement monotone).
\end{methode}

\begin{exoresolu}[Existence d'une solution par le TVI]
Soit $f$ la fonction définie sur $\mathbb{R}$ par $f(x)=x^{3}+x-3$.
\begin{enumerate}
\item Justifier que $f$ est continue sur $\mathbb{R}$.
\item Montrer que l'équation $f(x)=0$ admet au moins une solution dans l'intervalle $[1\,;\,2]$.
\end{enumerate}
\tcblower
\textbf{1. Continuité.}

$f$ est une fonction \textbf{polynôme} : elle est donc continue sur $\mathbb{R}$, en particulier sur l'intervalle $[1\,;\,2]$.

\textbf{2. Application du théorème des valeurs intermédiaires.}

Calculons les images des bornes :
\begin{align*}
f(1) &= 1^{3}+1-3 = -1, \\
f(2) &= 2^{3}+2-3 = 7.
\end{align*}
On a $f(1)=-1<0$ et $f(2)=7>0$ : $f(1)$ et $f(2)$ sont de signes contraires, autrement dit $0$ est compris entre $f(1)$ et $f(2)$.

Comme $f$ est continue sur $[1\,;\,2]$ et que $0$ est compris entre $f(1)$ et $f(2)$, le \textbf{théorème des valeurs intermédiaires} garantit qu'il existe au moins un réel $c\in[1\,;\,2]$ tel que $f(c)=0$.

\textbf{Conclusion} : l'équation $x^{3}+x-3=0$ admet au moins une solution dans $[1\,;\,2]$.
\end{exoresolu}

\courssub{Prolongement par continuité}

\begin{definition}[Prolongement par continuité]
Soit $f$ une fonction définie sur un intervalle $I$ \textbf{privé d'un point $a$} ($a$ est une borne d'adhérence de $D_f$). Si $\lim_{x\to a} f(x) = \ell \in \mathbb{R}$ (limite \textbf{finie}), on peut définir une nouvelle fonction $\tilde{f}$ sur $I$ tout entier par :
\[ \tilde{f}(x) = \begin{cases} f(x) & \text{si } x \neq a \\ \ell & \text{si } x = a \end{cases} \]
La fonction $\tilde{f}$ est continue en $a$ (et sur $I$ si $f$ l'était par morceaux). C'est le \cle{prolongement par continuité} de $f$ en $a$.
Si la limite est infinie ou n'existe pas, le prolongement par continuité est \textbf{impossible} (on a une asymptote verticale ou un saut).
\end{definition}

\begin{methode}[Prolonger une fonction par continuité en $a$]
\begin{enumerate}
\item Vérifier que $f$ n'est pas définie en $a$ mais l'est autour (point isolé du domaine).
\item \textbf{Calculer $\lim_{x\to a} f(x)$} (factoriser si forme $\dfrac{0}{0}$).
\item Si cette limite est un \textbf{nombre fini $\ell$} : $f$ est prolongeable par continuité.
\item \textbf{Écrire le prolongement} $\tilde{f}$ en donnant sa définition par morceaux (ou l'expression simplifiée valable sur tout l'intervalle).
\item Si la limite est infinie ou n'existe pas : conclure à l'impossibilité.
\end{enumerate}
\end{methode}

\begin{exoresolu}[Prolongement par continuité]
Soit $f$ la fonction définie par $f(x)=\dfrac{x^{2}-4}{x-2}$.
\begin{enumerate}
\item Préciser le domaine de définition de $f$.
\item Montrer que $f$ est prolongeable par continuité en $2$ et définir ce prolongement.
\end{enumerate}
\tcblower
\textbf{1. Domaine de définition.}

$f$ est définie lorsque $x-2\neq 0$, c'est-à-dire $x\neq 2$. Donc $D_f=\mathbb{R}\setminus\{2\}$.

\textbf{2. Prolongement par continuité.}

Calculons la limite de $f$ en $2$. La substitution donne $\dfrac{0}{0}$ : on factorise le numérateur grâce à l'identité remarquable $x^{2}-4=(x-2)(x+2)$ :
\[ f(x)=\dfrac{(x-2)(x+2)}{x-2}=x+2 \quad \text{pour tout } x\neq 2. \]
Donc :
\[ \lim_{x\to 2} f(x)=\lim_{x\to 2}(x+2)=4. \]
Cette limite est \textbf{finie} : $f$ est donc prolongeable par continuité en $2$.

Le prolongement par continuité de $f$ en $2$ est la fonction $\tilde{f}$ définie sur $\mathbb{R}$ tout entier par :
\[ \tilde{f}(x)=\begin{cases} \dfrac{x^{2}-4}{x-2} & \text{si } x\neq 2 \\ 4 & \text{si } x=2 \end{cases} \]
c'est-à-dire simplement $\tilde{f}(x)=x+2$ pour tout réel $x$.

\begin{popfigure}
\begin{tikzpicture}
\begin{axis}[
    width=\linewidth, height=6cm,
    axis lines=middle,
    xlabel=$x$, ylabel=$y$,
    xmin=0, xmax=4, ymin=0, ymax=7,
    xtick={1,2,3}, ytick={2,3,4,5,6},
    domain=0:1.9, samples=100,
    grid=major
]
% Branche gauche du trou
\addplot[popcurve, thick, domain=0:1.9] {x+2};
% Branche droite du trou
\addplot[popcurve, thick, domain=2.1:4] {x+2};
% Trou (point vide)
\addplot[only marks, mark=*, mark size=2.5, poporange, fill=white] coordinates {(2,4)};
% Point plein du prolongement (surimposé en vert)
\addplot[only marks, mark=*, mark size=2, popgreen] coordinates {(2,4)};
% Légendes
\node[poporange, above right] at (axis cs:2,4) {$f$ non définie};
\node[popgreen, below right] at (axis cs:2,4) {$\tilde{f}(2)=4$};
\end{axis}
\end{tikzpicture}
\legende{Prolongement par continuité : le « trou » en $(2,4)$ est bouché.}
\end{popfigure}

\textbf{Conclusion} : $f$ admet un prolongement par continuité en $2$, avec $\tilde{f}(2)=4$. La « cassure » de la courbe en $x=2$ est ainsi réparée.
\end{exoresolu}

\courssub{Synthèse et applications à la vie courante}

\begin{methode}[Modéliser une situation concrète par la continuité]
Face à un problème concret (coût, température, vitesse, tarification, population) :
\begin{enumerate}
\item \textbf{Identifier la fonction} qui modélise la situation et la variable (temps, quantité, distance) ;
\item \textbf{Traduire la question} en langage mathématique :
  \begin{itemize}
  \item Comportement à long terme $\to$ limite à l'infini (asymptote horizontale) ;
  \item Seuil critique, rupture $\to$ limite en un point, continuité ;
  \item Existence d'une valeur atteinte $\to$ TVI (équation $f(x)=k$).
  \end{itemize}
\item \textbf{Appliquer la méthode adaptée} vue dans cette leçon ;
\item \textbf{Conclure par une phrase en français} qui répond à la question posée, avec les unités.
\end{enumerate}
\end{methode}

\begin{exoresolu}[Tarification et continuité]
Une entreprise de transport à Douala facture la location d'un camion selon la durée $x$ (en heures) :
\[ P(x)=\begin{cases} 5000x & \text{si } 0\leq x\leq 4 \\ 4000x+4000 & \text{si } x>4 \end{cases} \]
où $P(x)$ est le prix en francs CFA.
\begin{enumerate}
\item Calculer $P(4)$.
\item La fonction $P$ est-elle continue en $4$ ? Interpréter.
\end{enumerate}
\tcblower
\textbf{1.} Comme $4\leq 4$, on utilise la première formule :
\[ P(4)=5000\times 4=\SI{20000}{\text{FCFA}}. \]

\textbf{2. Continuité en $4$.}

\underline{Limite à gauche} ($x<4$) : $\lim_{x\to 4^-} P(x)=\lim_{x\to 4} 5000x = \SI{20000}{\text{FCFA}}$.

\underline{Limite à droite} ($x>4$) : $\lim_{x\to 4^+} P(x)=\lim_{x\to 4}(4000x+4000)=\SI{16000}{\text{FCFA}}+4000=\SI{20000}{\text{FCFA}}$.

Les deux limites sont égales à $\SI{20000}{\text{FCFA}}=P(4)$ : les trois conditions de continuité sont remplies.

\begin{popfigure}
\begin{tikzpicture}
\begin{axis}[
    width=\linewidth, height=6cm,
    axis lines=middle,
    xlabel=$x$ (heures), ylabel=$P(x)$ (FCFA),
    xmin=0, xmax=8, ymin=0, ymax=40000,
    xtick={1,2,3,4,5,6,7}, ytick={10000,20000,30000},
    domain=0:4, samples=50,
    grid=major,
    yticklabel style={/pgf/number format/fixed, /pgf/number format/precision=0}
]
% Branche 1 : 5000x
\addplot[popcurve, thick, domain=0:4] {5000*x};
% Branche 2 : 4000x+4000
\addplot[popcurve, thick, domain=4:8] {4000*x+4000};
% Point de raccord
\addplot[only marks, mark=*, mark size=2, popgreen] coordinates {(4,20000)};
\end{axis}
\end{tikzpicture}
\legende{Tarification continue en $x=4$ : pas de saut de prix.}
\end{popfigure}

\textbf{Conclusion} : $P$ est \textbf{continue en $4$}. Concrètement, le tarif ne subit aucun saut brutal au passage de la quatrième heure : un client qui loue $4$ h $1$ min ne paie qu'une très légère différence par rapport à $4$ h exactement. Le tarif est donc « juste » au raccordement des deux formules.
\end{exoresolu}

\begin{savaistu}[Le TVI et la recherche de zéros : méthode de la dichotomie]
Le théorème des valeurs intermédiaires est la base mathématique de la \cle{méthode de dichotomie} (ou bisection), algorithme standard pour \emph{approcher} numériquement la solution d'une équation $f(x)=0$ sur un intervalle où $f$ est continue et change de signe. C'est la méthode utilisée par les calculatrices et logiciels (Excel, Python, tableurs) pour « résoudre » une équation. Au Probatoire, on demande souvent de \emph{montrer l'existence} (TVI) puis d'\emph{encadrer} la solution par dichotomie (quelques itérations).
\end{savaistu}

\begin{attention}[Les 5 erreurs qui coûtent des points au Probatoire]
\begin{enumerate}
\item \textbf{Conclure sur une forme indéterminée.} Écrire « $\lim_{x\to 3}\dfrac{x^{2}-9}{x-3}=\dfrac{0}{0}$ donc la limite n'existe pas » est faux : $\dfrac{0}{0}$ est un signal pour \emph{factoriser et simplifier}, jamais une conclusion.
\item \textbf{Oublier le signe du dénominateur pour une limite infinie.} $\dfrac{3}{x-1}$ tend vers $+\infty$ à droite de $1$ mais vers $-\infty$ à gauche : il faut toujours préciser si $x$ tend vers $a$ par valeurs supérieures ($a^+$) ou inférieures ($a^-$).
\item \textbf{Oublier de vérifier la continuité avant d'appliquer le TVI.} Le théorème des valeurs intermédiaires exige que $f$ soit \emph{continue} sur $[a\,;\,b]$ : cette hypothèse doit être écrite et justifiée dans la copie (ex: « $f$ polynôme donc continue sur $\mathbb{R}$ »).
\item \textbf{Confondre limite et valeur.} Une fonction peut avoir une limite en $a$ sans être définie en $a$ (cas du prolongement par continuité) : ne jamais écrire $f(a)$ quand $a\notin D_f$.
\item \textbf{Prolonger par continuité avec une limite infinie.} La condition est une limite \emph{finie} : si $\lim_{x\to a} f(x)=\pm\infty$, le prolongement est impossible (on a une asymptote verticale, pas un prolongement).
\end{enumerate}
\end{attention}

\newpage

\coursec{Exercices}

\courssub{Je vérifie mes connaissances}

\begin{exercice}[QCM : limites et continuité]
Pour chaque question, une seule réponse est exacte. Entoure la bonne réponse.
\begin{enumerate}
\item $\displaystyle\lim_{x \to 2}\ (3x-5)$ est égale à :
\begin{enumerate}
\item $0$ \hspace{1cm} \item $1$ \hspace{1cm} \item $11$
\end{enumerate}
\item $\displaystyle\lim_{x \to 4}\dfrac{x^2-16}{x-4}$ est égale à :
\begin{enumerate}
\item $0$ \hspace{1cm} \item $8$ \hspace{1cm} \item n'existe pas
\end{enumerate}
\item Une fonction $f$ est continue en $a$ si :
\begin{enumerate}
\item $f(a)$ existe ;
\item $\displaystyle\lim_{x \to a} f(x) = f(a)$ ;
\item $f$ est définie sur $\mathbb{R}$.
\end{enumerate}
\item La fonction $f(x)=\dfrac{1}{x-3}$ est continue sur :
\begin{enumerate}
\item $\mathbb{R}$ \hspace{1cm} \item $]-\infty\,;\,3[$ et sur $]3\,;\,+\infty[$ \hspace{1cm} \item $[3\,;\,+\infty[$
\end{enumerate}
\end{enumerate}
\end{exercice}

\begin{exercice}[Vrai ou faux ?]
Réponds par Vrai ou Faux, en justifiant chaque réponse.
\begin{enumerate}
\item Si $f$ est définie en $a$, alors $f$ est continue en $a$.
\item Toute fonction polynôme est continue sur $\mathbb{R}$.
\item Si $f$ est continue sur $[a\,;\,b]$ et si $f(a)\times f(b) < 0$, alors l'équation $f(x)=0$ admet au moins une solution dans $]a\,;\,b[$.
\item Si $\displaystyle\lim_{x \to a} f(x)$ existe, alors $f$ est définie en $a$.
\end{enumerate}
\end{exercice}

\begin{exercice}[Définitions]
\begin{enumerate}
\item Complète : on dit que $f$ est continue en $a$ lorsque $f$ est définie en $a$ et $\displaystyle\lim_{x \to a} f(x) = \ldots$
\item Énonce le théorème des valeurs intermédiaires pour une fonction continue sur un intervalle $[a\,;\,b]$.
\item Qu'appelle-t-on prolongement par continuité d'une fonction en un point $a$ ?
\end{enumerate}
\end{exercice}

\begin{exercice}[Calculs directs de limites]
Calculer les limites suivantes, en levant l'indétermination lorsqu'il y en a une :
\begin{enumerate}
\item $\displaystyle\lim_{x \to 2}\ \frac{x^2-5x+6}{x-2}$ ;
\item $\displaystyle\lim_{x \to 4}\ \frac{x^2-16}{x-4}$ ;
\item $\displaystyle\lim_{x \to -1}\ \frac{x^2-1}{x+1}$ ;
\item $\displaystyle\lim_{x \to 3}\ (2x^2-7x+4)$.
\end{enumerate}
\end{exercice}

\courssub{Je m'entraîne}

\begin{exercice}[Continuité d'une fonction définie par morceaux]
Soit $f$ la fonction définie sur $\mathbb{R}$ par :
\[
f(x) = \begin{cases} 3x-1 & \text{si } x < 2 \\ x+3 & \text{si } x \geq 2 \end{cases}
\]
\begin{enumerate}
\item Calculer $f(2)$.
\item Calculer $\displaystyle\lim_{\substack{x \to 2 \\ x < 2}} f(x)$ et $\displaystyle\lim_{\substack{x \to 2 \\ x > 2}} f(x)$.
\item La fonction $f$ est-elle continue en $2$ ? Justifier.
\item Interpréter graphiquement le résultat obtenu.
\end{enumerate}
\end{exercice}

\begin{exercice}[Détermination d'un paramètre]
Soit $m$ un réel et $f$ la fonction définie sur $\mathbb{R}$ par :
\[
f(x) = \begin{cases} x^2+1 & \text{si } x \leq 1 \\ mx & \text{si } x > 1 \end{cases}
\]
Déterminer la valeur du réel $m$ pour que $f$ soit continue en $1$.
\end{exercice}

\begin{exercice}[Théorème des valeurs intermédiaires]
Soit $f$ la fonction définie sur $\mathbb{R}$ par $f(x) = x^3 - 3x + 1$.
\begin{enumerate}
\item Vérifier que $f$ est continue sur $\mathbb{R}$.
\item Calculer $f(0)$ et $f(1)$.
\item En déduire que l'équation $f(x)=0$ admet au moins une solution $\alpha$ dans l'intervalle $]0\,;\,1[$.
\item Par balayage, donner un encadrement de $\alpha$ à $0{,}1$ près.
\end{enumerate}
\end{exercice}

\begin{exercice}[Prolongement par continuité]
Soit $f$ la fonction définie par $f(x) = \dfrac{x^2 - x - 6}{x - 3}$.
\begin{enumerate}
\item Quel est l'ensemble de définition de $f$ ?
\item Factoriser le numérateur de $f(x)$.
\item Calculer $\displaystyle\lim_{x \to 3} f(x)$.
\item La fonction $f$ admet-elle un prolongement par continuité en $3$ ? Si oui, définir explicitement ce prolongement.
\end{enumerate}
\end{exercice}

\courssub{Je me prépare au Probatoire}

\begin{exercice}[Problème : la course de moto-taxi]
À Douala, le prix d'une course de moto-taxi est fixé ainsi : $500$ FCFA pour toute distance inférieure ou égale à $2$ km ; au-delà de $2$ km, on ajoute $200$ FCFA par kilomètre supplémentaire. On note $P(x)$ le prix en FCFA d'une course de $x$ kilomètres ($x > 0$).
\begin{enumerate}
\item Montrer que $P(x) = \begin{cases} 500 & \text{si } 0 < x \leq 2 \\ 500 + 200(x-2) & \text{si } x > 2 \end{cases}$
\item Calculer $P(2)$, puis $\displaystyle\lim_{\substack{x \to 2 \\ x < 2}} P(x)$ et $\displaystyle\lim_{\substack{x \to 2 \\ x > 2}} P(x)$.
\item La fonction $P$ est-elle continue en $2$ ? Justifier la réponse.
\item Interpréter ce résultat pour le client : le prix subit-il un « saut » brutal au moment où l'on dépasse $2$ km ?
\item Un client dispose de $900$ FCFA. Quelle distance maximale peut-il parcourir ?
\end{enumerate}
\end{exercice}

\begin{exercice}[Problème : le réservoir d'eau]
Le réservoir d'eau d'un atelier de mécanique se vide de façon continue. À $8$ h, il contient $500$ litres ; à $12$ h, il n'en contient plus que $100$ litres. On note $V(t)$ le volume d'eau (en litres) dans le réservoir à l'instant $t$ (en heures), pour $t \in [8\,;\,12]$. On admet que $V$ est une fonction continue sur $[8\,;\,12]$.
\begin{enumerate}
\item Donner $V(8)$ et $V(12)$.
\item Justifier que $250$ est une valeur intermédiaire entre $V(8)$ et $V(12)$.
\item En citant précisément le théorème utilisé, montrer qu'il existe au moins un instant $t_0 \in [8\,;\,12]$ où le réservoir contenait exactement $250$ litres.
\item On suppose maintenant que $V(t) = 500 - 100(t-8)$. Déterminer l'instant $t_0$ exact où le réservoir contient $250$ litres, et donner l'heure correspondante.
\end{enumerate}
\end{exercice}

\begin{exercice}[Problème : étude complète d'une fonction]
Un menuisier fabrique des caisses en bois. Le coût de production, en milliers de FCFA, de $x$ caisses ($x \in [0\,;\,10]$) est modélisé par la fonction :
\[
C(x) = x^3 - 6x^2 + 9x + 2.
\]
\begin{enumerate}
\item Justifier que $C$ est continue sur $[0\,;\,10]$.
\item Calculer $C(0)$ et $C(2)$. En déduire que l'équation $C(x) = 4$ admet au moins une solution dans l'intervalle $]0\,;\,2[$.
\item Par balayage, encadrer cette solution à $0{,}1$ près.
\item On considère la fonction $g$ définie par $g(x) = \dfrac{x^2 - 9}{x - 3}$.
\begin{enumerate}
\item Quel est l'ensemble de définition de $g$ ?
\item Montrer que $g$ admet un prolongement par continuité en $3$ et définir ce prolongement.
\end{enumerate}
\item Le menuisier affirme : « Le coût moyen par caisse, donné par $\dfrac{C(x)}{x}$, n'est pas défini pour $x = 0$, donc on ne peut rien dire en $0$. » Calculer $\displaystyle\lim_{x \to 0^+} \frac{C(x)}{x}$ et commenter cette affirmation.
\end{enumerate}
\end{exercice}

\newpage

\coursec{Corrigés des exercices}

\coursec{Corrigés des exercices — Limites et continuité (Première Technique)}

\begin{corrige}[Exercice 1 — QCM : limites et continuité]
\begin{enumerate}
\item \textbf{Réponse : $1$}. \par
Calcul direct : $\displaystyle\lim_{x \to 2} (3x-5) = 3\times 2 - 5 = 6 - 5 = 1$.
\item \textbf{Réponse : $8$}. \par
On factorise le numérateur : $x^2-16 = (x-4)(x+4)$. Pour $x \neq 4$, $\dfrac{x^2-16}{x-4} = x+4$. Donc $\displaystyle\lim_{x \to 4} \dfrac{x^2-16}{x-4} = \lim_{x \to 4} (x+4) = 8$.
\item \textbf{Réponse : $\displaystyle\lim_{x \to a} f(x) = f(a)$}. \par
C'est la définition même de la continuité en un point : $f$ définie en $a$, la limite existe et est égale à $f(a)$.
\item \textbf{Réponse : $]-\infty\,;\,3[$ et sur $]3\,;\,+\infty[$}. \par
La fonction $f(x)=\dfrac{1}{x-3}$ est une fonction rationnelle, continue sur son ensemble de définition $\mathbb{R}\setminus\{3\}$, c'est-à-dire sur les deux intervalles ouverts $]-\infty\,;\,3[$ et $]3\,;\,+\infty[$.
\end{enumerate}
\end{corrige}

\begin{corrige}[Exercice 2 — Vrai ou faux ?]
\begin{enumerate}
\item \textbf{Faux}. \par
\emph{Justification :} La définition de $f$ en $a$ est nécessaire mais non suffisante. Il faut de plus que la limite en $a$ existe et soit égale à $f(a)$. Contre-exemple : $f(x) = \begin{cases} 0 & \text{si } x \neq 0 \\ 1 & \text{si } x = 0 \end{cases}$. $f$ est définie en $0$ ($f(0)=1$) mais $\lim_{x\to 0}f(x)=0 \neq f(0)$, donc $f$ n'est pas continue en $0$.
\item \textbf{Vrai}. \par
\emph{Justification :} Toute fonction polynôme est continue sur $\mathbb{R}$ (c'est un résultat fondamental du cours).
\item \textbf{Vrai}. \par
\emph{Justification :} C'est l'énoncé exact du \textbf{théorème des valeurs intermédiaires} (TVI) : si $f$ est continue sur $[a\,;\,b]$ et si $f(a)\times f(b) < 0$ (donc $0$ est entre $f(a)$ et $f(b)$), alors il existe au moins un $c \in ]a\,;\,b[$ tel que $f(c)=0$.
\item \textbf{Faux}. \par
\emph{Justification :} L'existence de la limite en $a$ n'implique pas que $f$ soit définie en $a$. Exemple classique : $f(x) = \dfrac{x^2-1}{x-1}$ en $a=1$. $\lim_{x\to 1}f(x) = 2$ existe, mais $f(1)$ n'existe pas (division par zéro).
\end{enumerate}
\end{corrige}

\begin{corrige}[Exercice 3 — Définitions]
\begin{enumerate}
\item On dit que $f$ est continue en $a$ lorsque $f$ est définie en $a$ et $\displaystyle\lim_{x \to a} f(x) = \textbf{$f(a)$}$.
\item \textbf{Théorème des valeurs intermédiaires (TVI) :} Soit $f$ une fonction continue sur un intervalle $[a\,;\,b]$. Pour tout réel $k$ compris entre $f(a)$ et $f(b)$, il existe au moins un $c \in [a\,;\,b]$ tel que $f(c) = k$. \par
(Si $f(a) \neq f(b)$, on peut préciser $c \in ]a\,;\,b[$).
\item On appelle \textbf{prolongement par continuité} de $f$ en $a$ (où $f$ n'est pas définie mais admet une limite finie $\ell$) la fonction $\tilde{f}$ définie sur $D_f \cup \{a\}$ par :
\[
\tilde{f}(x) = \begin{cases} f(x) & \text{si } x \in D_f \\ \ell & \text{si } x = a \end{cases}
\]
Cette fonction $\tilde{f}$ est continue en $a$.
\end{enumerate}
\end{corrige}

\begin{corrige}[Exercice 4 — Calculs directs de limites]
\begin{enumerate}
\item $\displaystyle\lim_{x \to 2} \frac{x^2-5x+6}{x-2}$. \par
Numérateur : $x^2-5x+6 = (x-2)(x-3)$. Pour $x \neq 2$, $\frac{(x-2)(x-3)}{x-2} = x-3$. \par
$\displaystyle\lim_{x \to 2} (x-3) = 2-3 = \boxed{-1}$.
\item $\displaystyle\lim_{x \to 4} \frac{x^2-16}{x-4}$. \par
$x^2-16 = (x-4)(x+4)$. Pour $x \neq 4$, la fraction se simplifie en $x+4$. \par
$\displaystyle\lim_{x \to 4} (x+4) = \boxed{8}$.
\item $\displaystyle\lim_{x \to -1} \frac{x^2-1}{x+1}$. \par
$x^2-1 = (x-1)(x+1)$. Pour $x \neq -1$, la fraction vaut $x-1$. \par
$\displaystyle\lim_{x \to -1} (x-1) = -1-1 = \boxed{-2}$.
\item $\displaystyle\lim_{x \to 3} (2x^2-7x+4)$. \par
Fonction polynôme, continue partout. On calcule directement : $2\times 3^2 - 7\times 3 + 4 = 18 - 21 + 4 = \boxed{1}$.
\end{enumerate}
\end{corrige}

\begin{corrige}[Exercice 5 — Continuité d'une fonction définie par morceaux]
\begin{enumerate}
\item Pour $x \geq 2$, $f(x) = x+3$. Donc $f(2) = 2+3 = \boxed{5}$.
\item \begin{itemize}
\item Limite à gauche ($x < 2$) : $f(x) = 3x-1$. $\displaystyle\lim_{\substack{x \to 2 \\ x < 2}} f(x) = 3\times 2 - 1 = \boxed{5}$.
\item Limite à droite ($x > 2$) : $f(x) = x+3$. $\displaystyle\lim_{\substack{x \to 2 \\ x > 2}} f(x) = 2+3 = \boxed{5}$.
\end{itemize}
\item Les deux limites latérales existent, sont égales entre elles (valeur commune $5$), et cette valeur est égale à $f(2)=5$. \par
Donc $\displaystyle\lim_{x \to 2} f(x) = f(2)$. La fonction $f$ \textbf{est continue en $2$}.
\item \textbf{Interprétation graphique :} Le graphe de $f$ est constitué de deux demi-droites : $y=3x-1$ pour $x<2$ et $y=x+3$ pour $x \geq 2$. Elles se rejoignent exactement au point de coordonnées $(2\,;\,5)$. Il n'y a \textbf{aucun saut} (discontinuité) au point d'abscisse $2$ ; le trait est continu.
\end{enumerate}
\end{corrige}

\begin{corrige}[Exercice 6 — Détermination d'un paramètre]
Pour que $f$ soit continue en $1$, il faut : $f(1) = \displaystyle\lim_{\substack{x \to 1 \\ x > 1}} f(x)$ (la limite à gauche est automatiquement $f(1)$ car $f(x)=x^2+1$ pour $x\leq 1$).
\begin{itemize}
\item $f(1) = 1^2 + 1 = 2$.
\item Pour $x > 1$, $f(x) = mx$. $\displaystyle\lim_{\substack{x \to 1 \\ x > 1}} f(x) = m \times 1 = m$.
\end{itemize}
La condition de continuité donne $m = 2$.
\par
\boxed{m = 2}
\end{corrige}

\begin{corrige}[Exercice 7 — Théorème des valeurs intermédiaires]
\begin{enumerate}
\item $f(x) = x^3 - 3x + 1$ est une fonction polynôme, donc \textbf{continue sur $\mathbb{R}$} (et en particulier sur tout intervalle $[a\,;\,b]$).
\item $f(0) = 0^3 - 0 + 1 = \boxed{1}$. \par
$f(1) = 1^3 - 3\times 1 + 1 = 1 - 3 + 1 = \boxed{-1}$.
\item $f$ est continue sur $[0\,;\,1]$. On a $f(0)=1 > 0$ et $f(1)=-1 < 0$, donc $f(0)\times f(1) < 0$. \par
D'après le \textbf{théorème des valeurs intermédiaires}, il existe au moins un $\alpha \in ]0\,;\,1[$ tel que $f(\alpha) = 0$. L'équation $f(x)=0$ admet donc au moins une solution dans $]0\,;\,1[$.
\item \textbf{Balayage (pas de $0{,}1$) :}
\begin{center}
\begin{tabular}{|c|ccccccccccc|}\hline
$x$ & $0$ & $0,1$ & $0,2$ & $0,3$ & $0,4$ & $0,5$ & $0,6$ & $0,7$ & $0,8$ & $0,9$ & $1$ \\ \hline
$f(x)$ & $1$ & $0,701$ & $0,408$ & $0,127$ & $-0,136$ & $-0,375$ & $-0,584$ & $-0,757$ & $-0,888$ & $-0,971$ & $-1$ \\ \hline
\end{tabular}
\end{center}
$f(0,3) > 0$ et $f(0,4) < 0$. La solution $\alpha$ est encadrée par : \boxed{0,3 < \alpha < 0,4}.
\end{enumerate}
\end{corrige}

\begin{corrige}[Exercice 8 — Prolongement par continuité]
\begin{enumerate}
\item $f(x) = \dfrac{x^2 - x - 6}{x - 3}$. Le dénominateur s'annule pour $x=3$. \par
Ensemble de définition : $\boxed{D_f = \mathbb{R} \setminus \{3\}}$.
\item Numérateur : $x^2 - x - 6 = (x-3)(x+2)$.
\item Pour $x \neq 3$, $f(x) = \dfrac{(x-3)(x+2)}{x-3} = x+2$. \par
$\displaystyle\lim_{x \to 3} f(x) = \lim_{x \to 3} (x+2) = \boxed{5}$.
\item Oui, $f$ admet une limite finie ($5$) en $3$ (point hors de son ensemble de définition). \par
Le \textbf{prolongement par continuité} est la fonction $\tilde{f}$ définie sur $\mathbb{R}$ par :
\[
\tilde{f}(x) = \begin{cases} \dfrac{x^2 - x - 6}{x - 3} = x+2 & \text{si } x \neq 3 \\ 5 & \text{si } x = 3 \end{cases}
\]
Soit simplement $\boxed{\tilde{f}(x) = x+2 \quad \forall x \in \mathbb{R}}$.
\end{enumerate}
\end{corrige}

\begin{corrige}[Exercice 9 — Problème : la course de moto-taxi]
\begin{enumerate}
\item \textbf{Modélisation :} \par
Pour $0 < x \leq 2$ km : tarif forfaitaire $\SI{500}{FCFA}$ $\Rightarrow P(x) = 500$. \par
Pour $x > 2$ km : on paie le forfait $\SI{500}{FCFA}$ pour les 2 premiers km, plus $\SI{200}{FCFA}$ par km supplémentaire. Le nombre de km supplémentaires est $x-2$. \par
Prix $P(x) = 500 + 200(x-2)$. \par
D'où la fonction :
\[
P(x) = \begin{cases} 500 & \text{si } 0 < x \leq 2 \\ 500 + 200(x-2) & \text{si } x > 2 \end{cases}
\]
\item $P(2) = 500$ (première branche). \par
$\displaystyle\lim_{\substack{x \to 2 \\ x < 2}} P(x) = \lim_{x \to 2^-} 500 = \boxed{500}$. \par
$\displaystyle\lim_{\substack{x \to 2 \\ x > 2}} P(x) = \lim_{x \to 2^+} [500 + 200(x-2)] = 500 + 200(0) = \boxed{500}$.
\item Les deux limites latérales existent, sont égales à $500$, et $P(2)=500$. \par
Donc $\displaystyle\lim_{x \to 2} P(x) = P(2)$. La fonction $P$ \textbf{est continue en $2$}.
\item \textbf{Interprétation :} Non, le prix ne subit \textbf{aucun saut brutal} au passage de $\SI{2}{km}$. La tarification est « douce » : le prix évolue de manière continue, ce qui est équitable pour le client (pas de surcoût brutal dès le premier mètre au-delà de $\SI{2}{km}$).
\item Le client a $\SI{900}{FCFA}$. \par
Si $x \leq 2$, $P(x)=500 \leq 900$ : possible. \par
Si $x > 2$, on résout $500 + 200(x-2) \leq 900$. \par
$200(x-2) \leq 400 \iff x-2 \leq 2 \iff x \leq 4$. \par
La distance maximale est \boxed{\SI{4}{km}}.
\end{enumerate}

\textbf{Barème indicatif (Probatoire) :}
\begin{itemize}
\item Modélisation correcte : 2 pts
\item Calculs limites et $P(2)$ : 2 pts
\item Conclusion continuité + justification : 2 pts
\item Interprétation concrète : 1 pt
\item Résolution inéquation distance max : 3 pts
\end{itemize}
\textbf{Points de vigilance :} Bien distinguer $P(2)$ (valeur) et les limites latérales ; citer explicitement la définition de la continuité (limite = valeur) ; pour la question 5, vérifier que la solution $x=4$ appartient bien au domaine de la branche utilisée ($x>2$).
\end{corrige}

\begin{corrige}[Exercice 10 — Problème : le réservoir d'eau]
\begin{enumerate}
\item $V(8) = \boxed{\SI{500}{L}}$ ; $V(12) = \boxed{\SI{100}{L}}$.
\item $V(8)=500$ et $V(12)=100$. Comme $100 < 250 < 500$, la valeur $250$ est bien une \textbf{valeur intermédiaire} entre $V(8)$ et $V(12)$.
\item La fonction $V$ est continue sur l'intervalle $[8\,;\,12]$ (donnée de l'énoncé). \par
D'après le \textbf{théorème des valeurs intermédiaires}, puisque $250$ est compris entre $V(8)=500$ et $V(12)=100$, il existe au moins un instant $t_0 \in [8\,;\,12]$ tel que $V(t_0) = \SI{250}{L}$.
\item On résout $V(t_0) = 250$ avec l'expression $V(t) = 500 - 100(t-8)$. \par
$500 - 100(t_0-8) = 250$ \par
$100(t_0-8) = 250$ \par
$t_0 - 8 = 2,5$ \par
$t_0 = \boxed{10,5}$ heures. \par
Cela correspond à \boxed{10 h 30}.
\end{enumerate}

\textbf{Barème indicatif (Probatoire) :}
\begin{itemize}
\item Lecture valeurs $V(8), V(12)$ : 1 pt
\item Justification valeur intermédiaire : 1 pt
\item Citation TVI + conclusion existence : 3 pts (attention : citer le TVI et vérifier hypothèses : continuité sur $[8,12]$, $250$ entre les bornes)
\item Résolution équation + conversion heure : 2 pts
\end{itemize}
\textbf{Points de vigilance :} Ne pas oublier de rappeler l'hypothèse de continuité de $V$ sur $[8,12]$ avant d'appliquer le TVI. L'heure s'écrit en heures/minutes (10h30), pas en décimal (10,5 h) pour la réponse finale.
\end{corrige}

\begin{corrige}[Exercice 11 — Problème : étude complète d'une fonction]
\begin{enumerate}
\item $C(x) = x^3 - 6x^2 + 9x + 2$ est une fonction polynôme, donc \textbf{continue sur $\mathbb{R}$}, et a fortiori sur l'intervalle $[0\,;\,10]$.
\item $C(0) = 2$. \par
$C(2) = 2^3 - 6\times 2^2 + 9\times 2 + 2 = 8 - 24 + 18 + 2 = \boxed{4}$. \par
On a $C(0)=2$ et $C(2)=4$. La valeur $4$ est atteinte en $x=2$ (borne de l'intervalle). Pour garantir une solution \textbf{dans l'intervalle ouvert $]0\,;\,2[$}, on observe que $C(1) = 1 - 6 + 9 + 2 = 6$. \par
Sur $[0\,;\,1]$, $C$ est continue, $C(0)=2 < 4 < 6 = C(1)$. D'après le \textbf{TVI}, il existe au moins une solution $\alpha \in ]0\,;\,1[ \subset ]0\,;\,2[$ telle que $C(\alpha)=4$. \par
(On peut aussi factoriser $C(x)-4 = (x-2)(x^2-4x+1)$ ; l'autre racine est $2-\sqrt{3} \approx 0,268 \in ]0\,;\,2[$.)
\item \textbf{Balayage pour $\alpha \in ]0\,;\,1[$ (pas $0{,}1$) :}
\begin{center}
\begin{tabular}{|c|ccccccccccc|}\hline
$x$ & $0$ & $0,1$ & $0,2$ & $0,3$ & $0,4$ & $0,5$ & $0,6$ & $0,7$ & $0,8$ & $0,9$ & $1$ \\ \hline
$C(x)$ & $2$ & $2,841$ & $3,568$ & $4,187$ & $4,704$ & $5,125$ & $5,456$ & $5,703$ & $5,872$ & $5,969$ & $6$ \\ \hline
\end{tabular}
\end{center}
$C(0,2) \approx 3,568 < 4$ et $C(0,3) \approx 4,187 > 4$. \par
La solution est dans $]0,2\,;\,0,3[$. \par
\boxed{0,2 < \alpha < 0,3}.
\item \textbf{Étude de $g(x) = \dfrac{x^2 - 9}{x - 3}$ :}
\begin{enumerate}
\item Dénominateur nul pour $x=3$. $\boxed{D_g = \mathbb{R} \setminus \{3\}}$.
\item $x^2-9 = (x-3)(x+3)$. Pour $x \neq 3$, $g(x) = x+3$. \par
$\displaystyle\lim_{x \to 3} g(x) = \lim_{x \to 3} (x+3) = 6$ (limite finie). \par
$g$ admet un prolongement par continuité en $3$ : la fonction $\tilde{g}$ définie sur $\mathbb{R}$ par $\tilde{g}(x) = x+3$ (avec $\tilde{g}(3)=6$).
\end{enumerate}
\item \textbf{Coût moyen :} $\dfrac{C(x)}{x} = \dfrac{x^3 - 6x^2 + 9x + 2}{x} = x^2 - 6x + 9 + \dfrac{2}{x}$ pour $x > 0$. \par
$\displaystyle\lim_{x \to 0^+} \frac{C(x)}{x} = \lim_{x \to 0^+} \left(x^2 - 6x + 9 + \frac{2}{x}\right) = 0 - 0 + 9 + \boxed{+\infty} = \boxed{+\infty}$. \par
\textbf{Commentaire :} L'affirmation du menuisier est \textbf{fausse}. Même si le coût moyen n'est pas \emph{défini} en $x=0$ (division par zéro), la limite à droite existe (elle est $+\infty$). Cela signifie que pour une production très faible, le coût moyen par caisse devient \textbf{arbitrairement grand} (les frais fixes $\SI{2}{kFCFA}$ sont répartis sur très peu de caisses). On peut donc décrire précisément le comportement près de $0$.
\end{enumerate}

\textbf{Barème indicatif (Probatoire) :}
\begin{itemize}
\item Continuité $C$ (polynôme) : 1 pt
\item Calcul $C(0), C(2)$ + TVI correct (utiliser $C(1)$ ou factorisation) : 3 pts
\item Balayage correct + encadrement $0,1$ : 2 pts
\item Ensemble définition $g$ + factorisation + limite + prolongement : 4 pts
\item Limite coût moyen + interprétation : 3 pts
\end{itemize}
\textbf{Points de vigilance :}
\begin{itemize}
\item Pour le TVI question 2 : $C(2)=4$ ne permet pas de conclure à une solution dans $]0,2[$ (c'est la borne). Il faut utiliser $C(1)=6$ (ou factoriser) pour appliquer le TVI sur $[0,1]$.
\item Balayage : bien calculer $C(x)$ pour $x$ variant de $0$ à $1$ (pas $0,1$) et repérer le changement de signe par rapport à $4$ (ici entre $0,2$ et $0,3$).
\item Prolongement par continuité : donner l'expression explicite de la fonction prolongée ($\tilde{g}(x)=x+3$).
\item Limite en $0^+$ : ne pas écrire « n'existe pas », écrire $+\infty$ et l'interpréter (coût moyen qui explose).
\end{itemize}
\end{corrige}

\begin{corrige}[Exercice 12 — Limites infinies et limites à l'infini (complément)]
\textbf{On considère la fonction $h$ définie sur $\mathbb{R}^*_+$ par $h(x) = \dfrac{2x^2 + 5}{x - 1}$.}
\begin{enumerate}
\item \textbf{Limite en $+\infty$ :} \par
On factorise par la plus grande puissance du dénominateur ($x$) : \par
$h(x) = \dfrac{x^2(2 + \frac{5}{x^2})}{x(1 - \frac{1}{x})} = x \cdot \dfrac{2 + \frac{5}{x^2}}{1 - \frac{1}{x}}$. \par
Quand $x \to +\infty$, $\frac{5}{x^2} \to 0$ et $\frac{1}{x} \to 0$, donc la fraction tend vers $2$. \par
$x \to +\infty$ et la fraction tend vers $2 > 0$, donc $\boxed{\lim_{x \to +\infty} h(x) = +\infty}$.
\item \textbf{Limite en $1$ (point d'interdiction) :} \par
Numérateur en $1$ : $2(1)^2 + 5 = 7 \neq 0$. Dénominateur $x-1 \to 0$. \par
\emph{Étude du signe :}
\begin{itemize}
\item $x \to 1^+$ ($x>1$) : $x-1 > 0$, numérateur $>0$ $\Rightarrow h(x) \to +\infty$.
\item $x \to 1^-$ ($x<1$) : $x-1 < 0$, numérateur $>0$ $\Rightarrow h(x) \to -\infty$.
\end{itemize}
Donc $\boxed{\lim_{x \to 1^+} h(x) = +\infty}$ et $\boxed{\lim_{x \to 1^-} h(x) = -\infty}$. \par
La droite $x=1$ est une \textbf{asymptote verticale}.
\item \textbf{Asymptote oblique (complément) :} \par
Division euclidienne : $(2x^2 + 5) \div (x-1) = 2x + 2$ reste $7$. \par
$h(x) = 2x + 2 + \dfrac{7}{x-1}$. \par
$\lim_{x \to \pm\infty} \dfrac{7}{x-1} = 0$. \par
La courbe admet l'asymptote oblique d'équation $\boxed{y = 2x + 2}$.
\end{enumerate}

\textbf{Application concrète (Pharmacocinétique) :} La concentration $C(t)$ (en mg/L) d'un médicament dans le sang $t$ heures après l'injection est modélisée par $C(t) = \dfrac{10t}{t^2+1}$ pour $t \ge 0$.
\begin{enumerate}
\item $\lim_{t \to +\infty} C(t) = \lim_{t \to +\infty} \dfrac{10t}{t^2} = \lim_{t \to +\infty} \dfrac{10}{t} = 0$. \par
\textbf{Interprétation :} À long terme, le médicament est totalement éliminé de l'organisme.
\item Maximum : $C'(t) = \dfrac{10(t^2+1) - 10t(2t)}{(t^2+1)^2} = \dfrac{10 - 10t^2}{(t^2+1)^2}$. \par
$C'(t)=0 \iff t=1$ (car $t\ge 0$). $C(1)=5$. \par
Concentration maximale $\SI{5}{mg/L}$ atteinte à $\SI{1}{h}$.
\end{enumerate}
\end{corrige}

\newpage

\coursec{Fiche bilan}

\begin{popfigure}
\begin{tikzpicture}[
  mindmap,
  concept color=popblue,
  level 1 concept/.append style={level distance=3.5cm, sibling angle=60},
  level 2 concept/.append style={level distance=2.5cm},
  every node/.style={font=\small},
  text=white
]
\node[concept] {Limites \& Continuité}
  child[concept color=popblue] {
    node[concept] {Limite en un point}
    child { node[concept, scale=0.6] {Finie / Infinie} }
    child { node[concept, scale=0.6] {Opérations} }
    child { node[concept, scale=0.6] {Formes indéterminées} }
  }
  child[concept color=poporange] {
    node[concept] {Limite à l'infini}
    child { node[concept, scale=0.6] {$+\infty$ / $-\infty$} }
    child { node[concept, scale=0.6] {Comparaison} }
    child { node[concept, scale=0.6] {Croissances} }
  }
  child[concept color=popgreen] {
    node[concept] {Continuité en un point}
    child { node[concept, scale=0.6] {$\lim=f(a)$} }
    child { node[concept, scale=0.6] {Gauche / Droite} }
    child { node[concept, scale=0.6] {Discontinuités} }
  }
  child[concept color=poppurple] {
    node[concept] {Continuité sur intervalle}
    child { node[concept, scale=0.6] {Th. Valeurs Intermédiaires} }
    child { node[concept, scale=0.6] {Th. Encadrement} }
    child { node[concept, scale=0.6] {Bijection} }
  }
  child[concept color=poppink] {
    node[concept] {Propriétés algébriques}
    child { node[concept, scale=0.6] {Somme, Produit} }
    child { node[concept, scale=0.6] {Composée} }
    child { node[concept, scale=0.6] {Inverse bijection} }
  }
  child[concept color=popgold] {
    node[concept] {Prolongement par continuité}
    child { node[concept, scale=0.6] {Limite finie en $x_0$} }
    child { node[concept, scale=0.6] {Fonction prolongée} }
    child { node[concept, scale=0.6] {Ex : $\frac{\sin x}{x}$} }
  };
\end{tikzpicture}
\legende{Carte mentale : Limites et continuité (Première Technique)}
\end{popfigure}

\begin{bilan}

\textbf{Définitions clés}
\begin{itemize}
  \item \cle{Limite finie en $a$} : $\lim_{x\to a}f(x)=\ell \iff \forall \varepsilon>0,\ \exists \eta>0,\ |x-a|<\eta \Rightarrow |f(x)-\ell|<\varepsilon$.
  \item \cle{Limite infinie} : $\lim_{x\to a}f(x)=+\infty \iff \forall M>0,\ \exists \eta>0,\ |x-a|<\eta \Rightarrow f(x)>M$.
  \item \cle{Limite à $+\infty$} : $\lim_{x\to +\infty}f(x)=\ell \iff \forall \varepsilon>0,\ \exists A>0,\ x>A \Rightarrow |f(x)-\ell|<\varepsilon$.
  \item \cle{Continuité en $a$} : $f$ définie en $a$ ET $\lim_{x\to a}f(x)=f(a)$.
  \item \cle{Prolongement par continuité} : $f$ non définie en $a$ mais $\lim_{x\to a}f(x)=\ell$ fini $\Rightarrow$ $\tilde f(x)=\begin{cases}f(x)&x\neq a\\ \ell&x=a\end{cases}$ continue en $a$.
\end{itemize}

\textbf{Théorèmes \& Formules à connaître par cœur}
\begin{center}
\begin{tabularx}{\linewidth}{|l|X|}\hline
\rowcolor{popblueL} \textbf{Théorème / Règle} & \textbf{Énoncé / Condition} \\ \hline
Opérations sur limites (finies) & $\lim(f\pm g)=\lim f\pm\lim g$, $\lim(fg)=\lim f\lim g$, $\lim(1/g)=1/\lim g$ (si $\lim g\neq0$) \\ \hline
Formes indéterminées & $0\times\infty$, $\infty-\infty$, $\frac{0}{0}$, $\frac{\infty}{\infty}$, $1^\infty$ $\rightarrow$ factoriser, conjuguer, changer de variable \\ \hline
Croissances comparées ($x\to+\infty$) & $\ln x \ll x^a \ll x^b \ll e^x \ll x^x$ ($0<a<b$) \\ \hline
Th. Valeurs Intermédiaires (TVI) & $f$ continue sur $[a,b]$, $\forall k\in[f(a),f(b)],\ \exists c\in[a,b], f(c)=k$ \\ \hline
Th. Encadrement & $f$ continue sur $[a,b] \Rightarrow \exists m,M,\ \forall x\in[a,b],\ m\le f(x)\le M$ \\ \hline
Bijection & $f$ continue et \textbf{strictement monotone} sur intervalle $\Rightarrow$ bijection vers son image \\ \hline
Prolongement par continuité & Si $\lim_{x\to x_0}f(x)=\ell$ (fini) et $x_0\notin D_f$, alors $\tilde f$ continue sur $D_f\cup\{x_0\}$ \\ \hline
\end{tabularx}
\end{center}

\textbf{Méthodes express}
\begin{enumerate}
  \item \textbf{Limite en un point} : Remplacer $x$ par $a$ $\rightarrow$ si forme déterminée : fini. Si $0/0$ ou $\infty/\infty$ : factoriser (racines, identités remarquables), simplifier, puis remplacer.
  \item \textbf{Limite à $\pm\infty$ (fonction rationnelle)} : Factoriser par le terme de plus haut degré numérateur et dénominateur $\rightarrow$ garder les termes dominants.
  \item \textbf{Étudier la continuité en $a$} : 1) $a\in D_f$ ? 2) Calculer $\lim_{x\to a}f(x)$ (gauche/droite si nécessaire). 3) Comparer à $f(a)$.
  \item \textbf{Prolonger par continuité} : Calculer $\lim_{x\to x_0}f(x)=\ell$. Définir $\tilde f(x_0)=\ell$. Vérifier continuité.
  \item \textbf{Résoudre $f(x)=k$ (TVI)} : Vérifier continuité sur $[a,b]$, signes de $f(a)-k$ et $f(b)-k$ opposés $\Rightarrow$ solution existe. Monotonie $\Rightarrow$ unique.
\end{enumerate}

\end{bilan}

\begin{aretenir}[Checklist avant le Probatoire]
$\square$ Savoir écrire la définition $\varepsilon-\eta$ d'une limite finie en un point.\\
$\square$ Reconnaître et lever les 5 formes indéterminées classiques ($0/0$, $\infty/\infty$, $0\times\infty$, $\infty-\infty$, $1^\infty$).\\
$\square$ Appliquer les règles de croissance comparée pour les limites à l'infini.\\
$\square$ Vérifier la continuité en un point : 3 conditions (appartenance, existence limite, égalité).\\
$\square$ Différencier continuité à gauche, à droite et bilatère.\\
$\square$ Citer et appliquer le Théorème des Valeurs Intermédiaires (hypothèses : continuité + intervalle).\\
$\square$ Citer et appliquer le Théorème d'Encadrement (majorant/minorant sur segment).\\
$\square$ Montrer qu'une fonction est bijection sur un intervalle (continuité + monotonie stricte).\\
$\square$ Construire le prolongement par continuité en un point exclu du domaine (limite finie).\\
$\square$ Justifier chaque étape : « car $f$ est continue sur $[a,b]$ comme somme/quotient de fonctions continues... ».\\
$\square$ Présenter un tableau de variations complet (valeurs aux bornes, limites, flèches, image).
\end{aretenir}

\findecours

\end{document}
